% Complete reviewed Unit 16; editorial master: chapters/ch16/student.md
% Markdown SHA-256: 89dbd7124d92ad9cc944f06547a6094e08f2dd25060d86806d04ff0da8679c4d
% One-use initial assembly followed by complete content and page review.
\newcommand{\runningtitle}{Unit 16. Narrative and speech}
% Layout inherited from the reviewed previous course; set a running title per unit.
% This is a small hand-maintained TeX layout, not a conversion pipeline.
\RequirePackage{fix-cm}
\documentclass[11pt,twoside]{article}
\usepackage{fontspec}
\setmainfont{Cambria}[
  Renderer=HarfBuzz,
  BoldFont=cambriab.ttf,
  ItalicFont=cambriai.ttf,
  BoldItalicFont=cambriaz.ttf
]
\setmonofont{lmmono10-regular.otf}[BoldFont=lmmonolt10-bold.otf]
\usepackage[paperwidth=165mm,paperheight=241mm,inner=21mm,outer=19mm,
  top=22mm,bottom=20mm,headheight=14pt,headsep=7mm,footskip=10mm]{geometry}
\usepackage{xcolor,array,booktabs,longtable,tabularx}
\usepackage{enumitem,fancyhdr}
\usepackage[unicode,hidelinks]{hyperref}
\definecolor{ink}{HTML}{000000}
\definecolor{muted}{HTML}{4C5558}
\setlength{\parindent}{0pt}
\setlength{\parskip}{3.5pt}
\AtBeginDocument{\fontsize{11}{14.3}\selectfont}
\setlength{\tabcolsep}{4pt}
\renewcommand{\arraystretch}{1.10}
\setlist[itemize]{leftmargin=1.3em,label=\textbullet,itemsep=2pt,topsep=3pt,parsep=0pt}
\setlist[enumerate]{leftmargin=1.8em,itemsep=3pt,topsep=4pt,parsep=0pt}
\widowpenalty=10000
\clubpenalty=10000
\emergencystretch=1em
\raggedbottom
\pagestyle{fancy}
\fancyhf{}
\providecommand{\runningtitle}{Old English}
\fancyhead[LE,RO]{\small\itshape \runningtitle}
\fancyhead[LO,RE]{\small Old English}
\fancyfoot[LE,RO]{\small\thepage}
\renewcommand{\headrulewidth}{0.3pt}
\newcommand{\booktitle}[2]{
  {\fontsize{9}{11}\selectfont\color{muted}#1\par}
  \vspace{3pt}
  {\fontsize{20}{24}\selectfont\color{ink}#2\par}
  \vspace{6pt}
}
\newcommand{\source}[1]{\par{\footnotesize\color{muted}#1\par}}
\newcommand{\readingtext}[1]{
  \par\addvspace{5pt}
  \begingroup\leftskip=1.5em
    \fontsize{12}{16}\selectfont #1\par
  \endgroup\addvspace{5pt}
}
\makeatletter
\renewcommand{\section}{\@startsection{section}{1}{\z@}
  {-2.4ex \@plus -0.6ex \@minus -0.2ex}{0.8ex \@plus 0.2ex}
  {\normalfont\fontsize{12}{15}\selectfont\bfseries\color{ink}}}
\renewcommand{\subsection}{\@startsection{subsection}{2}{\z@}
  {-1.8ex \@plus -0.5ex \@minus -0.2ex}{0.5ex \@plus 0.2ex}
  {\normalfont\fontsize{11}{14}\selectfont\bfseries\color{ink}}}
\makeatother

\hyphenpenalty=10000
\exhyphenpenalty=10000
\emergencystretch=3em
\newcommand{\itempara}[2]{\par\addvspace{3pt}\noindent\hangindent=1.8em\hangafter=1\makebox[1.8em][l]{#1.}#2\par}
\newcommand{\reserve}[1]{\par\penalty-100\begingroup\dimen0=\pagegoal\advance\dimen0 by-\pagetotal\ifdim\dimen0<#1\newpage\fi\endgroup}
\begin{document}
\addfontfeatures{Ligatures=TeXOff}
\raggedright
\booktitle{UNIT 16}{Extended narrative and speech}

% mdparagraph:P001:begin
A longer narrative combines events, perceptions, reports and quoted speech. The form of each clause establishes its own participants; context connects those participants across the passage. A change of speaker can change the reference of I, you, here and now without changing the people or events described.
% mdparagraph:P001:end

% mdparagraph:P002:begin
The vocabulary and complete systems below prepare the original exercises and Cook’s complete joint Group XII, XIII, Exercises 32–36. The final reading follows Apollonius through a game at the public baths and the invitation that results from it. The four translated examples and the lettered exercises are new material.
% mdparagraph:P002:end

% mdparagraph:P003:begin
\textbf{Scope.} Learn the productive forms for the meanings and constructions stated. A \textbf{reference} form may be consulted; it does not require unaided recall. \textbf{Recognition} forms are identified where they occur. \textbf{Supplied help} remains available during the corresponding task. In verb tables, 1, 2 and 3 denote person; singular you and plural you remain distinct. Preserve vowel length and the distinction between e and ę.
% mdparagraph:P003:end

% teaching: NEW-U16-PRIOR-RECALL
% mdparagraph:P004:begin
The previously learned case functions, strong and weak adjective choice, personal and dual pronouns, direct quotation, finite content clauses, perception with an infinitive, relative clauses, passives and perfect constructions remain in use. A familiar form does not make every meaning or complement of its verb familiar. The new extensions are stated explicitly below.
% mdparagraph:P004:end

\reserve{7\baselineskip}
% mdparagraph:P005:begin
\section*{1. Nouns, adjectives and both}
% mdparagraph:P005:end

% mdparagraph:P006:begin
Use the nominative for a subject and the accusative for a direct object when the verb takes one. An adnominal genitive expresses possession or another noun-dependent relation. A recipient or a preposition takes its own licensed case. These ordinary uses apply to the nouns below; a special government pattern must be learned separately. Identical endings do not remove the distinction between grammatical positions.\textsuperscript{1}
% mdparagraph:P006:end

\reserve{7\baselineskip}
% mdparagraph:P007:begin
\subsection*{Strong masculine nouns}
% mdparagraph:P007:end

% mdparagraph:P008:begin
The nominative and accusative singular have the bare stem; the genitive adds -es and the dative -e. The plural has nominative/\allowbreak{}accusative -as, genitive -a and dative -um.
% mdparagraph:P008:end

% teaching: NEW-U16-VOC-NEAHGEBUR
\par\begin{minipage}{\linewidth}\raggedright
% mdparagraph:P009:begin
\textbf{nēahgebūr} — masculine; neighbour. Stem \textbf{nēahgebūr-}.
% mdparagraph:P009:end

% teaching: NEW-U16-NOM-NEAHGEBUR
% table: NEW-U16-NOM-NEAHGEBUR-DISPLAY
% mdtable:T01:begin
\par\begingroup
\fontsize{10.5}{13.2}\selectfont
\setlength{\LTleft}{0pt}\setlength{\LTright}{\fill}
\noindent
\begin{tabular}{@{}>{\raggedright\arraybackslash}p{36.99mm}>{\raggedright\arraybackslash}p{40.69mm}>{\raggedright\arraybackslash}p{40.69mm}@{}}
\toprule
\textbf{Case} & \textbf{Singular} & \textbf{Plural} \\
\midrule
Nominative & nēahgebūr & nēahgebūras \\
Genitive & nēahgebūres & nēahgebūra \\
Dative & nēahgebūre & nēahgebūrum \\
Accusative & nēahgebūr & nēahgebūras \\
\bottomrule
\end{tabular}
\par\endgroup

% mdtable:T01:end
\end{minipage}\par
% teaching: NEW-U16-VOC-THOTHOR
\par\begin{minipage}{\linewidth}\raggedright
% mdparagraph:P010:begin
\textbf{ðoðor} — masculine; ball. Stem \textbf{ðoðor-}.
% mdparagraph:P010:end

% mdparagraph:P011:begin
The medial o remains before endings. This rare word is learned for the ball exercises and the gymnasium scene.
% mdparagraph:P011:end

% teaching: NEW-U16-NOM-THOTHOR
% table: NEW-U16-NOM-THOTHOR-DISPLAY
% mdtable:T02:begin
\par\begingroup
\fontsize{10.5}{13.2}\selectfont
\setlength{\LTleft}{0pt}\setlength{\LTright}{\fill}
\noindent
\begin{tabular}{@{}>{\raggedright\arraybackslash}p{39.46mm}>{\raggedright\arraybackslash}p{39.46mm}>{\raggedright\arraybackslash}p{39.46mm}@{}}
\toprule
\textbf{Case} & \textbf{Singular} & \textbf{Plural} \\
\midrule
Nominative & ðoðor & ðoðoras \\
Genitive & ðoðores & ðoðora \\
Dative & ðoðore & ðoðorum \\
Accusative & ðoðor & ðoðoras \\
\bottomrule
\end{tabular}
\par\endgroup

% mdtable:T02:end
\end{minipage}\par
% teaching: NEW-U16-VOC-COSS
\par\begin{minipage}{\linewidth}\raggedright
% mdparagraph:P012:begin
\textbf{coss} — masculine; kiss. Stem \textbf{coss-}.
% mdparagraph:P012:end

% teaching: NEW-U16-NOM-COSS
% table: NEW-U16-NOM-COSS-DISPLAY
% mdtable:T03:begin
\par\begingroup
\fontsize{10.5}{13.2}\selectfont
\setlength{\LTleft}{0pt}\setlength{\LTright}{\fill}
\noindent
\begin{tabular}{@{}>{\raggedright\arraybackslash}p{39.46mm}>{\raggedright\arraybackslash}p{39.46mm}>{\raggedright\arraybackslash}p{39.46mm}@{}}
\toprule
\textbf{Case} & \textbf{Singular} & \textbf{Plural} \\
\midrule
Nominative & coss & cossas \\
Genitive & cosses & cossa \\
Dative & cosse & cossum \\
Accusative & coss & cossas \\
\bottomrule
\end{tabular}
\par\endgroup

% mdtable:T03:end
\end{minipage}\par
% teaching: NEW-U16-VOC-DISC
\par\begin{minipage}{\linewidth}\raggedright
% mdparagraph:P013:begin
\textbf{disc} — masculine; dish. Stem \textbf{disc-}.
% mdparagraph:P013:end

% teaching: NEW-U16-NOM-DISC
% table: NEW-U16-NOM-DISC-DISPLAY
% mdtable:T04:begin
\par\begingroup
\fontsize{10.5}{13.2}\selectfont
\setlength{\LTleft}{0pt}\setlength{\LTright}{\fill}
\noindent
\begin{tabular}{@{}>{\raggedright\arraybackslash}p{39.46mm}>{\raggedright\arraybackslash}p{39.46mm}>{\raggedright\arraybackslash}p{39.46mm}@{}}
\toprule
\textbf{Case} & \textbf{Singular} & \textbf{Plural} \\
\midrule
Nominative & disc & discas \\
Genitive & disces & disca \\
Dative & disce & discum \\
Accusative & disc & discas \\
\bottomrule
\end{tabular}
\par\endgroup

% mdtable:T04:end
\end{minipage}\par
% teaching: NEW-U16-VOC-ATH
\par\begin{minipage}{\linewidth}\raggedright
% mdparagraph:P014:begin
\textbf{āð} — masculine; oath. Stem \textbf{āð-}.
% mdparagraph:P014:end

% teaching: NEW-U16-NOM-ATH
% table: NEW-U16-NOM-ATH-DISPLAY
% mdtable:T05:begin
\par\begingroup
\fontsize{10.5}{13.2}\selectfont
\setlength{\LTleft}{0pt}\setlength{\LTright}{\fill}
\noindent
\begin{tabular}{@{}>{\raggedright\arraybackslash}p{39.46mm}>{\raggedright\arraybackslash}p{39.46mm}>{\raggedright\arraybackslash}p{39.46mm}@{}}
\toprule
\textbf{Case} & \textbf{Singular} & \textbf{Plural} \\
\midrule
Nominative & āð & āðas \\
Genitive & āðes & āða \\
Dative & āðe & āðum \\
Accusative & āð & āðas \\
\bottomrule
\end{tabular}
\par\endgroup

% mdtable:T05:end
\end{minipage}\par
% teaching: NEW-U16-VOC-TEAR
\par\begin{minipage}{\linewidth}\raggedright
% mdparagraph:P015:begin
\textbf{tēar} — masculine; tear. Stem \textbf{tēar-}.
% mdparagraph:P015:end

% teaching: NEW-U16-NOM-TEAR
% table: NEW-U16-NOM-TEAR-DISPLAY
% mdtable:T06:begin
\par\begingroup
\fontsize{10.5}{13.2}\selectfont
\setlength{\LTleft}{0pt}\setlength{\LTright}{\fill}
\noindent
\begin{tabular}{@{}>{\raggedright\arraybackslash}p{39.46mm}>{\raggedright\arraybackslash}p{39.46mm}>{\raggedright\arraybackslash}p{39.46mm}@{}}
\toprule
\textbf{Case} & \textbf{Singular} & \textbf{Plural} \\
\midrule
Nominative & tēar & tēaras \\
Genitive & tēares & tēara \\
Dative & tēare & tēarum \\
Accusative & tēar & tēaras \\
\bottomrule
\end{tabular}
\par\endgroup

% mdtable:T06:end
\end{minipage}\par
% teaching: NEW-U16-VOC-MAEGENTHRYMM
\par\begin{minipage}{\linewidth}\raggedright
% mdparagraph:P016:begin
\textbf{mægenðrymm} — masculine; majesty; glory. Stem \textbf{mægenðrymm-}.
% mdparagraph:P016:end

% mdparagraph:P017:begin
Only the singular is learned for this meaning. No count-plural sense is assigned here.
% mdparagraph:P017:end

% teaching: NEW-U16-NOM-MAEGENTHRYMM
% table: NEW-U16-NOM-MAEGENTHRYMM-DISPLAY
% mdtable:T07:begin
\par\begingroup
\fontsize{10.5}{13.2}\selectfont
\setlength{\LTleft}{0pt}\setlength{\LTright}{\fill}
\noindent
\begin{tabular}{@{}>{\raggedright\arraybackslash}p{55.09mm}>{\raggedright\arraybackslash}p{66.1mm}@{}}
\toprule
\textbf{Case} & \textbf{Singular} \\
\midrule
Nominative & mægenðrymm \\
Genitive & mægenðrymmes \\
Dative & mægenðrymme \\
Accusative & mægenðrymm \\
\bottomrule
\end{tabular}
\par\endgroup

% mdtable:T07:end
\end{minipage}\par
\reserve{7\baselineskip}
% mdparagraph:P018:begin
\subsection*{The masculine title in -ere}
% mdparagraph:P018:end

% mdparagraph:P019:begin
The singular retains -e in the nominative, accusative and dative; the genitive ends in -es. Here fuluhtere is learned in the singular title “Baptist”.
% mdparagraph:P019:end

% teaching: NEW-U16-VOC-FULUHTERE
\par\begin{minipage}{\linewidth}\raggedright
% mdparagraph:P020:begin
\textbf{fuluhtere} — masculine; baptizer; Baptist as John’s title. Stem \textbf{fuluhter-}.
% mdparagraph:P020:end

% mdparagraph:P021:begin
Only the singular is learned for this meaning. No count-plural sense is assigned here.
% mdparagraph:P021:end

% teaching: NEW-U16-NOM-FULUHTERE
% table: NEW-U16-NOM-FULUHTERE-DISPLAY
% mdtable:T08:begin
\par\begingroup
\fontsize{10.5}{13.2}\selectfont
\setlength{\LTleft}{0pt}\setlength{\LTright}{\fill}
\noindent
\begin{tabular}{@{}>{\raggedright\arraybackslash}p{60.59mm}>{\raggedright\arraybackslash}p{60.59mm}@{}}
\toprule
\textbf{Case} & \textbf{Singular} \\
\midrule
Nominative & fuluhtere \\
Genitive & fuluhteres \\
Dative & fuluhtere \\
Accusative & fuluhtere \\
\bottomrule
\end{tabular}
\par\endgroup

% mdtable:T08:end
\end{minipage}\par
\reserve{7\baselineskip}
% mdparagraph:P022:begin
\subsection*{Weak masculine nouns}
% mdparagraph:P022:end

% mdparagraph:P023:begin
The nominative singular ends in -a. The other singular cases and the nominative/\allowbreak{}accusative plural end in -an; the genitive plural has -ena and the dative plural -um.
% mdparagraph:P023:end

% teaching: NEW-U16-VOC-CNAPA
\par\begin{minipage}{\linewidth}\raggedright
% mdparagraph:P024:begin
\textbf{cnapa} — masculine; boy. Stem \textbf{cnap-}.
% mdparagraph:P024:end

% teaching: NEW-U16-NOM-CNAPA
% table: NEW-U16-NOM-CNAPA-DISPLAY
% mdtable:T09:begin
\par\begingroup
\fontsize{10.5}{13.2}\selectfont
\setlength{\LTleft}{0pt}\setlength{\LTright}{\fill}
\noindent
\begin{tabular}{@{}>{\raggedright\arraybackslash}p{39.46mm}>{\raggedright\arraybackslash}p{39.46mm}>{\raggedright\arraybackslash}p{39.46mm}@{}}
\toprule
\textbf{Case} & \textbf{Singular} & \textbf{Plural} \\
\midrule
Nominative & cnapa & cnapan \\
Genitive & cnapan & cnapena \\
Dative & cnapan & cnapum \\
Accusative & cnapan & cnapan \\
\bottomrule
\end{tabular}
\par\endgroup

% mdtable:T09:end
\end{minipage}\par
% teaching: NEW-U16-VOC-GEGIERELA
\par\begin{minipage}{\linewidth}\raggedright
% mdparagraph:P025:begin
\textbf{gegierela} — masculine; garment; clothing. Stem \textbf{gegierel-}.
% mdparagraph:P025:end

% teaching: NEW-U16-NOM-GEGIERELA
% table: NEW-U16-NOM-GEGIERELA-DISPLAY
% mdtable:T10:begin
\par\begingroup
\fontsize{10.5}{13.2}\selectfont
\setlength{\LTleft}{0pt}\setlength{\LTright}{\fill}
\noindent
\begin{tabular}{@{}>{\raggedright\arraybackslash}p{38.19mm}>{\raggedright\arraybackslash}p{38.19mm}>{\raggedright\arraybackslash}p{42.0mm}@{}}
\toprule
\textbf{Case} & \textbf{Singular} & \textbf{Plural} \\
\midrule
Nominative & gegierela & gegierelan \\
Genitive & gegierelan & gegierelena \\
Dative & gegierelan & gegierelum \\
Accusative & gegierelan & gegierelan \\
\bottomrule
\end{tabular}
\par\endgroup

% mdtable:T10:end
\end{minipage}\par
\reserve{7\baselineskip}
% mdparagraph:P026:begin
\subsection*{Strong neuter nouns}
% mdparagraph:P026:end

% mdparagraph:P027:begin
Nominative and accusative coincide in each number. The singular genitive and dative use -es and -e; the plural genitive and dative use -a and -um. Learn the nominative/\allowbreak{}accusative plural with the noun. Furlang and nętt have no ending in those plural positions.
% mdparagraph:P027:end

% teaching: NEW-U16-VOC-FURLANG
\par\begin{minipage}{\linewidth}\raggedright
% mdparagraph:P028:begin
\textbf{furlang} — neuter; furlong. Stem \textbf{furlang-}.
% mdparagraph:P028:end

% teaching: NEW-U16-NOM-FURLANG
% table: NEW-U16-NOM-FURLANG-DISPLAY
% mdtable:T11:begin
\par\begingroup
\fontsize{10.5}{13.2}\selectfont
\setlength{\LTleft}{0pt}\setlength{\LTright}{\fill}
\noindent
\begin{tabular}{@{}>{\raggedright\arraybackslash}p{39.46mm}>{\raggedright\arraybackslash}p{39.46mm}>{\raggedright\arraybackslash}p{39.46mm}@{}}
\toprule
\textbf{Case} & \textbf{Singular} & \textbf{Plural} \\
\midrule
Nominative & furlang & furlang \\
Genitive & furlanges & furlanga \\
Dative & furlange & furlangum \\
Accusative & furlang & furlang \\
\bottomrule
\end{tabular}
\par\endgroup

% mdtable:T11:end
\end{minipage}\par
% teaching: NEW-U16-VOC-NETT
\par\begin{minipage}{\linewidth}\raggedright
% mdparagraph:P029:begin
\textbf{nętt} — neuter; net. Stem \textbf{nętt-}.
% mdparagraph:P029:end

% mdparagraph:P030:begin
The zero nominative/\allowbreak{}accusative plural belongs to this noun’s class; it is not inferred merely from the written double consonant.
% mdparagraph:P030:end

% teaching: NEW-U16-NOM-NETT
% table: NEW-U16-NOM-NETT-DISPLAY
% mdtable:T12:begin
\par\begingroup
\fontsize{10.5}{13.2}\selectfont
\setlength{\LTleft}{0pt}\setlength{\LTright}{\fill}
\noindent
\begin{tabular}{@{}>{\raggedright\arraybackslash}p{39.46mm}>{\raggedright\arraybackslash}p{39.46mm}>{\raggedright\arraybackslash}p{39.46mm}@{}}
\toprule
\textbf{Case} & \textbf{Singular} & \textbf{Plural} \\
\midrule
Nominative & nętt & nętt \\
Genitive & nęttes & nętta \\
Dative & nętte & nęttum \\
Accusative & nętt & nętt \\
\bottomrule
\end{tabular}
\par\endgroup

% mdtable:T12:end
\end{minipage}\par
% teaching: NEW-U16-VOC-TEMPEL
\par\begin{minipage}{\linewidth}\raggedright
% mdparagraph:P031:begin
\textbf{tempel} — neuter; temple. Stem \textbf{templ-}.
% mdparagraph:P031:end

% mdparagraph:P032:begin
Use templ- before the displayed vowel endings.
% mdparagraph:P032:end

% teaching: NEW-U16-NOM-TEMPEL
% table: NEW-U16-NOM-TEMPEL-DISPLAY
% mdtable:T13:begin
\par\begingroup
\fontsize{10.5}{13.2}\selectfont
\setlength{\LTleft}{0pt}\setlength{\LTright}{\fill}
\noindent
\begin{tabular}{@{}>{\raggedright\arraybackslash}p{39.46mm}>{\raggedright\arraybackslash}p{39.46mm}>{\raggedright\arraybackslash}p{39.46mm}@{}}
\toprule
\textbf{Case} & \textbf{Singular} & \textbf{Plural} \\
\midrule
Nominative & tempel & templu \\
Genitive & temples & templa \\
Dative & temple & templum \\
Accusative & tempel & templu \\
\bottomrule
\end{tabular}
\par\endgroup

% mdtable:T13:end
\end{minipage}\par
% teaching: NEW-U16-VOC-MAEDEN
\par\begin{minipage}{\linewidth}\raggedright
% mdparagraph:P033:begin
\textbf{mǣden} — neuter; girl; maiden. Stem \textbf{mǣden-}.
% mdparagraph:P033:end

% mdparagraph:P034:begin
Retain mǣden- in the selected forms.
% mdparagraph:P034:end

% teaching: NEW-U16-NOM-MAEDEN
% table: NEW-U16-NOM-MAEDEN-DISPLAY
% mdtable:T14:begin
\par\begingroup
\fontsize{10.5}{13.2}\selectfont
\setlength{\LTleft}{0pt}\setlength{\LTright}{\fill}
\noindent
\begin{tabular}{@{}>{\raggedright\arraybackslash}p{39.46mm}>{\raggedright\arraybackslash}p{39.46mm}>{\raggedright\arraybackslash}p{39.46mm}@{}}
\toprule
\textbf{Case} & \textbf{Singular} & \textbf{Plural} \\
\midrule
Nominative & mǣden & mǣden \\
Genitive & mǣdenes & mǣdena \\
Dative & mǣdene & mǣdenum \\
Accusative & mǣden & mǣden \\
\bottomrule
\end{tabular}
\par\endgroup

% mdtable:T14:end
\end{minipage}\par
% teaching: NEW-U16-VOC-WOLCEN
\par\begin{minipage}{\linewidth}\raggedright
% mdparagraph:P035:begin
\textbf{wolcen} — neuter; cloud. Stem \textbf{wolcn-}.
% mdparagraph:P035:end

% mdparagraph:P036:begin
Use wolcn- before the displayed vowel endings.
% mdparagraph:P036:end

% teaching: NEW-U16-NOM-WOLCEN
% table: NEW-U16-NOM-WOLCEN-DISPLAY
% mdtable:T15:begin
\par\begingroup
\fontsize{10.5}{13.2}\selectfont
\setlength{\LTleft}{0pt}\setlength{\LTright}{\fill}
\noindent
\begin{tabular}{@{}>{\raggedright\arraybackslash}p{39.46mm}>{\raggedright\arraybackslash}p{39.46mm}>{\raggedright\arraybackslash}p{39.46mm}@{}}
\toprule
\textbf{Case} & \textbf{Singular} & \textbf{Plural} \\
\midrule
Nominative & wolcen & wolcnu \\
Genitive & wolcnes & wolcna \\
Dative & wolcne & wolcnum \\
Accusative & wolcen & wolcnu \\
\bottomrule
\end{tabular}
\par\endgroup

% mdtable:T15:end
\end{minipage}\par
\reserve{7\baselineskip}
% mdparagraph:P037:begin
\subsection*{Strong feminine nouns}
% mdparagraph:P037:end

% mdparagraph:P038:begin
The selected nouns have no ending in the nominative singular. Their genitive and dative singular have -e. The accusative singular also has -e, except wēn, whose accusative is wēn. The plural, where assigned, has nominative/\allowbreak{}accusative -a, genitive -a and dative -um. Smyltnes doubles its final s before an ending.
% mdparagraph:P038:end

% teaching: NEW-U16-VOC-SMYLTNES
\par\begin{minipage}{\linewidth}\raggedright
% mdparagraph:P039:begin
\textbf{smyltnes} — feminine; calm; calmness. Stem \textbf{smyltness-}.
% mdparagraph:P039:end

% mdparagraph:P040:begin
Only the singular is learned for this meaning. No count-plural sense is assigned here.
% mdparagraph:P040:end

% teaching: NEW-U16-NOM-SMYLTNES
% table: NEW-U16-NOM-SMYLTNES-DISPLAY
% mdtable:T16:begin
\par\begingroup
\fontsize{10.5}{13.2}\selectfont
\setlength{\LTleft}{0pt}\setlength{\LTright}{\fill}
\noindent
\begin{tabular}{@{}>{\raggedright\arraybackslash}p{60.59mm}>{\raggedright\arraybackslash}p{60.59mm}@{}}
\toprule
\textbf{Case} & \textbf{Singular} \\
\midrule
Nominative & smyltnes \\
Genitive & smyltnesse \\
Dative & smyltnesse \\
Accusative & smyltnesse \\
\bottomrule
\end{tabular}
\par\endgroup

% mdtable:T16:end
\end{minipage}\par
% teaching: NEW-U16-VOC-STEFN
\par\begin{minipage}{\linewidth}\raggedright
% mdparagraph:P041:begin
\textbf{stefn} — feminine; voice. Stem \textbf{stefn-}.
% mdparagraph:P041:end

% teaching: NEW-U16-NOM-STEFN
% table: NEW-U16-NOM-STEFN-DISPLAY
% mdtable:T17:begin
\par\begingroup
\fontsize{10.5}{13.2}\selectfont
\setlength{\LTleft}{0pt}\setlength{\LTright}{\fill}
\noindent
\begin{tabular}{@{}>{\raggedright\arraybackslash}p{39.46mm}>{\raggedright\arraybackslash}p{39.46mm}>{\raggedright\arraybackslash}p{39.46mm}@{}}
\toprule
\textbf{Case} & \textbf{Singular} & \textbf{Plural} \\
\midrule
Nominative & stefn & stefna \\
Genitive & stefne & stefna \\
Dative & stefne & stefnum \\
Accusative & stefne & stefna \\
\bottomrule
\end{tabular}
\par\endgroup

% mdtable:T17:end
\end{minipage}\par
% teaching: NEW-U16-VOC-ANSIEN
\par\begin{minipage}{\linewidth}\raggedright
% mdparagraph:P042:begin
\textbf{ansīen} — feminine; appearance. Stem \textbf{ansīen-}.
% mdparagraph:P042:end

% mdparagraph:P043:begin
Only the singular is learned for this meaning. No count-plural sense is assigned here.
% mdparagraph:P043:end

% teaching: NEW-U16-NOM-ANSIEN
% table: NEW-U16-NOM-ANSIEN-DISPLAY
% mdtable:T18:begin
\par\begingroup
\fontsize{10.5}{13.2}\selectfont
\setlength{\LTleft}{0pt}\setlength{\LTright}{\fill}
\noindent
\begin{tabular}{@{}>{\raggedright\arraybackslash}p{60.59mm}>{\raggedright\arraybackslash}p{60.59mm}@{}}
\toprule
\textbf{Case} & \textbf{Singular} \\
\midrule
Nominative & ansīen \\
Genitive & ansīene \\
Dative & ansīene \\
Accusative & ansīene \\
\bottomrule
\end{tabular}
\par\endgroup

% mdtable:T18:end
\end{minipage}\par
% teaching: NEW-U16-VOC-SEALF
\par\begin{minipage}{\linewidth}\raggedright
% mdparagraph:P044:begin
\textbf{sealf} — feminine; ointment. Stem \textbf{sealf-}.
% mdparagraph:P044:end

% mdparagraph:P045:begin
Only the singular is learned for this meaning. No count-plural sense is assigned here.
% mdparagraph:P045:end

% teaching: NEW-U16-NOM-SEALF
% table: NEW-U16-NOM-SEALF-DISPLAY
% mdtable:T19:begin
\par\begingroup
\fontsize{10.5}{13.2}\selectfont
\setlength{\LTleft}{0pt}\setlength{\LTright}{\fill}
\noindent
\begin{tabular}{@{}>{\raggedright\arraybackslash}p{60.59mm}>{\raggedright\arraybackslash}p{60.59mm}@{}}
\toprule
\textbf{Case} & \textbf{Singular} \\
\midrule
Nominative & sealf \\
Genitive & sealfe \\
Dative & sealfe \\
Accusative & sealfe \\
\bottomrule
\end{tabular}
\par\endgroup

% mdtable:T19:end
\end{minipage}\par
% teaching: NEW-U16-VOC-WEN
\par\begin{minipage}{\linewidth}\raggedright
% mdparagraph:P046:begin
\textbf{wēn} — feminine; expectation; chance, especially wēn is ðæt. Stem \textbf{wēn-}.
% mdparagraph:P046:end

% mdparagraph:P047:begin
Only the singular is learned for this meaning. No count-plural sense is assigned here.
% mdparagraph:P047:end

% teaching: NEW-U16-NOM-WEN
% table: NEW-U16-NOM-WEN-DISPLAY
% mdtable:T20:begin
\par\begingroup
\fontsize{10.5}{13.2}\selectfont
\setlength{\LTleft}{0pt}\setlength{\LTright}{\fill}
\noindent
\begin{tabular}{@{}>{\raggedright\arraybackslash}p{60.59mm}>{\raggedright\arraybackslash}p{60.59mm}@{}}
\toprule
\textbf{Case} & \textbf{Singular} \\
\midrule
Nominative & wēn \\
Genitive & wēne \\
Dative & wēne \\
Accusative & wēn \\
\bottomrule
\end{tabular}
\par\endgroup

% mdtable:T20:end
\end{minipage}\par
% teaching: NEW-U16-VOC-SPRAEC
\par\begin{minipage}{\linewidth}\raggedright
% mdparagraph:P048:begin
\textbf{sprǣc} — feminine; saying; speech. Stem \textbf{sprǣc-}.
% mdparagraph:P048:end

% teaching: NEW-U16-NOM-SPRAEC
% table: NEW-U16-NOM-SPRAEC-DISPLAY
% mdtable:T21:begin
\par\begingroup
\fontsize{10.5}{13.2}\selectfont
\setlength{\LTleft}{0pt}\setlength{\LTright}{\fill}
\noindent
\begin{tabular}{@{}>{\raggedright\arraybackslash}p{39.46mm}>{\raggedright\arraybackslash}p{39.46mm}>{\raggedright\arraybackslash}p{39.46mm}@{}}
\toprule
\textbf{Case} & \textbf{Singular} & \textbf{Plural} \\
\midrule
Nominative & sprǣc & sprǣca \\
Genitive & sprǣce & sprǣca \\
Dative & sprǣce & sprǣcum \\
Accusative & sprǣce & sprǣca \\
\bottomrule
\end{tabular}
\par\endgroup

% mdtable:T21:end
\end{minipage}\par
\reserve{7\baselineskip}
% mdparagraph:P049:begin
\subsection*{Adjectives}
% mdparagraph:P049:end

% mdparagraph:P050:begin
The following tables give the complete strong and weak positive forms. Use the strong or weak series under the conditions already learned, with agreement in case, number and gender. In ordinary predicate use the adjective also agrees with its subject. The strong instrumental singular forms are for recognition only; their appearance here introduces no new instrumental construction. The other displayed positions are productive.\textsuperscript{2}
% mdparagraph:P050:end

% teaching: NEW-U16-VOC-HEALT
\par\begin{minipage}{\linewidth}\raggedright
% mdparagraph:P051:begin
\textbf{healt} — halt; lame.
% mdparagraph:P051:end

% mdparagraph:P052:begin
The heavy stem takes no -u in the strong feminine nominative singular or neuter nominative/\allowbreak{}accusative plural.
% mdparagraph:P052:end

% teaching: NEW-U16-ADJ-HEALT
% table: NEW-U16-ADJ-HEALT-DISPLAY
% mdtable:T22:begin
% course-table-style: 2026-10
% course-panel:T22-P01:begin
\par\begingroup
\fontsize{10.5}{13.2}\selectfont
\setlength{\LTleft}{0pt}\setlength{\LTright}{\fill}
\noindent
\begin{tabular}{@{}>{\raggedright\arraybackslash}p{34.062mm}>{\raggedright\arraybackslash}p{29.341mm}>{\raggedright\arraybackslash}p{25.069mm}>{\raggedright\arraybackslash}p{27.092mm}@{}}
\toprule
\textbf{Series /\allowbreak{} case} & \textbf{Masculine singular} & \textbf{Neuter singular} & \textbf{Feminine singular} \\
\midrule
Strong nominative & healt & healt & healt \\
Strong genitive & healtes & healtes & healtre \\
Strong dative & healtum & healtum & healtre \\
Strong accusative & healtne & healt & healte \\
Strong instrumental & healte & healte & — \\
Weak nominative & healta & healte & healte \\
Weak genitive & healtan & healtan & healtan \\
Weak dative & healtan & healtan & healtan \\
Weak accusative & healtan & healte & healtan \\
\bottomrule
\end{tabular}
\par\endgroup
% course-panel:T22-P01:end
\end{minipage}\par
\par\addvspace{5pt}
% course-panel:T22-P02:begin
\par\begingroup
\fontsize{10.5}{13.2}\selectfont
\setlength{\LTleft}{0pt}\setlength{\LTright}{\fill}
\noindent
\begin{tabular}{@{}>{\raggedright\arraybackslash}p{36.136mm}>{\raggedright\arraybackslash}p{28.98mm}>{\raggedright\arraybackslash}p{23.852mm}>{\raggedright\arraybackslash}p{26.595mm}@{}}
\toprule
\textbf{Series /\allowbreak{} case} & \textbf{Masculine plural} & \textbf{Neuter plural} & \textbf{Feminine plural} \\
\midrule
Strong nominative & healte & healt & healta \\
Strong genitive & healtra & healtra & healtra \\
Strong dative & healtum & healtum & healtum \\
Strong accusative & healte & healt & healta \\
Strong instrumental & — & — & — \\
Weak nominative & healtan & healtan & healtan \\
Weak genitive & healtra & healtra & healtra \\
Weak dative & healtum & healtum & healtum \\
Weak accusative & healtan & healtan & healtan \\
\bottomrule
\end{tabular}
\par\endgroup
% course-panel:T22-P02:end
% mdtable:T22:end

% teaching: NEW-U16-VOC-GELARED
\par\begin{minipage}{\linewidth}\raggedright
% mdparagraph:P053:begin
\textbf{gelǣred} — educated; trained; skilled.
% mdparagraph:P053:end

% mdparagraph:P054:begin
The selected forms keep the full gelǣred- stem before endings. Shorter contracted variants do not replace this table.
% mdparagraph:P054:end

% teaching: NEW-U16-ADJ-GELARED
% table: NEW-U16-ADJ-GELARED-DISPLAY
% mdtable:T23:begin
% course-table-style: 2026-10
% course-panel:T23-P01:begin
\par\begingroup
\fontsize{10.5}{13.2}\selectfont
\setlength{\LTleft}{0pt}\setlength{\LTright}{\fill}
\noindent
\begin{tabular}{@{}>{\raggedright\arraybackslash}p{33.348mm}>{\raggedright\arraybackslash}p{28.726mm}>{\raggedright\arraybackslash}p{25.754mm}>{\raggedright\arraybackslash}p{27.735mm}@{}}
\toprule
\textbf{Series /\allowbreak{} case} & \textbf{Masculine singular} & \textbf{Neuter singular} & \textbf{Feminine singular} \\
\midrule
Strong nominative & gelǣred & gelǣred & gelǣredu \\
Strong genitive & gelǣredes & gelǣredes & gelǣredre \\
Strong dative & gelǣredum & gelǣredum & gelǣredre \\
Strong accusative & gelǣredne & gelǣred & gelǣrede \\
Strong instrumental & gelǣrede & gelǣrede & — \\
Weak nominative & gelǣreda & gelǣrede & gelǣrede \\
Weak genitive & gelǣredan & gelǣredan & gelǣredan \\
Weak dative & gelǣredan & gelǣredan & gelǣredan \\
Weak accusative & gelǣredan & gelǣrede & gelǣredan \\
\bottomrule
\end{tabular}
\par\endgroup
% course-panel:T23-P01:end
\end{minipage}\par
\par\addvspace{5pt}
% course-panel:T23-P02:begin
\par\begingroup
\fontsize{10.5}{13.2}\selectfont
\setlength{\LTleft}{0pt}\setlength{\LTright}{\fill}
\noindent
\begin{tabular}{@{}>{\raggedright\arraybackslash}p{35.157mm}>{\raggedright\arraybackslash}p{28.195mm}>{\raggedright\arraybackslash}p{25.062mm}>{\raggedright\arraybackslash}p{27.151mm}@{}}
\toprule
\textbf{Series /\allowbreak{} case} & \textbf{Masculine plural} & \textbf{Neuter plural} & \textbf{Feminine plural} \\
\midrule
Strong nominative & gelǣrede & gelǣredu & gelǣreda \\
Strong genitive & gelǣredra & gelǣredra & gelǣredra \\
Strong dative & gelǣredum & gelǣredum & gelǣredum \\
Strong accusative & gelǣrede & gelǣredu & gelǣreda \\
Strong instrumental & — & — & — \\
Weak nominative & gelǣredan & gelǣredan & gelǣredan \\
Weak genitive & gelǣredra & gelǣredra & gelǣredra \\
Weak dative & gelǣredum & gelǣredum & gelǣredum \\
Weak accusative & gelǣredan & gelǣredan & gelǣredan \\
\bottomrule
\end{tabular}
\par\endgroup
% course-panel:T23-P02:end
% mdtable:T23:end

% teaching: NEW-U16-VOC-EORTHLIC
\par\begin{minipage}{\linewidth}\raggedright
% mdparagraph:P055:begin
\textbf{eorðlic} — earthly.
% mdparagraph:P055:end

% mdparagraph:P056:begin
The selected forms keep the full eorðlic- stem before endings.
% mdparagraph:P056:end

% teaching: NEW-U16-ADJ-EORTHLIC
% table: NEW-U16-ADJ-EORTHLIC-DISPLAY
% mdtable:T24:begin
% course-table-style: 2026-10
% course-panel:T24-P01:begin
\par\begingroup
\fontsize{10.5}{13.2}\selectfont
\setlength{\LTleft}{0pt}\setlength{\LTright}{\fill}
\noindent
\begin{tabular}{@{}>{\raggedright\arraybackslash}p{33.348mm}>{\raggedright\arraybackslash}p{28.726mm}>{\raggedright\arraybackslash}p{25.754mm}>{\raggedright\arraybackslash}p{27.735mm}@{}}
\toprule
\textbf{Series /\allowbreak{} case} & \textbf{Masculine singular} & \textbf{Neuter singular} & \textbf{Feminine singular} \\
\midrule
Strong nominative & eorðlic & eorðlic & eorðlicu \\
Strong genitive & eorðlices & eorðlices & eorðlicre \\
Strong dative & eorðlicum & eorðlicum & eorðlicre \\
Strong accusative & eorðlicne & eorðlic & eorðlice \\
Strong instrumental & eorðlice & eorðlice & — \\
Weak nominative & eorðlica & eorðlice & eorðlice \\
Weak genitive & eorðlican & eorðlican & eorðlican \\
Weak dative & eorðlican & eorðlican & eorðlican \\
Weak accusative & eorðlican & eorðlice & eorðlican \\
\bottomrule
\end{tabular}
\par\endgroup
% course-panel:T24-P01:end
\end{minipage}\par
\par\addvspace{5pt}
% course-panel:T24-P02:begin
\par\begingroup
\fontsize{10.5}{13.2}\selectfont
\setlength{\LTleft}{0pt}\setlength{\LTright}{\fill}
\noindent
\begin{tabular}{@{}>{\raggedright\arraybackslash}p{35.157mm}>{\raggedright\arraybackslash}p{28.195mm}>{\raggedright\arraybackslash}p{25.062mm}>{\raggedright\arraybackslash}p{27.151mm}@{}}
\toprule
\textbf{Series /\allowbreak{} case} & \textbf{Masculine plural} & \textbf{Neuter plural} & \textbf{Feminine plural} \\
\midrule
Strong nominative & eorðlice & eorðlicu & eorðlica \\
Strong genitive & eorðlicra & eorðlicra & eorðlicra \\
Strong dative & eorðlicum & eorðlicum & eorðlicum \\
Strong accusative & eorðlice & eorðlicu & eorðlica \\
Strong instrumental & — & — & — \\
Weak nominative & eorðlican & eorðlican & eorðlican \\
Weak genitive & eorðlicra & eorðlicra & eorðlicra \\
Weak dative & eorðlicum & eorðlicum & eorðlicum \\
Weak accusative & eorðlican & eorðlican & eorðlican \\
\bottomrule
\end{tabular}
\par\endgroup
% course-panel:T24-P02:end
% mdtable:T24:end

% teaching: NEW-U16-VOC-NEARU
\par\begin{minipage}{\linewidth}\raggedright
% mdparagraph:P057:begin
\textbf{nearu} — narrow.
% mdparagraph:P057:end

% mdparagraph:P058:begin
Learn the alternating stems together: nearu in the strong nominative singular, neuter accusative singular and neuter nominative/\allowbreak{}accusative plural; nearw- before vowel endings; nearo- before consonantal -ne, -re and -ra. The weak forms use nearw- except for genitive plural nearora.
% mdparagraph:P058:end

% teaching: NEW-U16-ADJ-NEARU
% table: NEW-U16-ADJ-NEARU-DISPLAY
% mdtable:T25:begin
% course-table-style: 2026-10
% course-panel:T25-P01:begin
\par\begingroup
\fontsize{10.5}{13.2}\selectfont
\setlength{\LTleft}{0pt}\setlength{\LTright}{\fill}
\noindent
\begin{tabular}{@{}>{\raggedright\arraybackslash}p{34.062mm}>{\raggedright\arraybackslash}p{29.341mm}>{\raggedright\arraybackslash}p{25.069mm}>{\raggedright\arraybackslash}p{27.092mm}@{}}
\toprule
\textbf{Series /\allowbreak{} case} & \textbf{Masculine singular} & \textbf{Neuter singular} & \textbf{Feminine singular} \\
\midrule
Strong nominative & nearu & nearu & nearu \\
Strong genitive & nearwes & nearwes & nearore \\
Strong dative & nearwum & nearwum & nearore \\
Strong accusative & nearone & nearu & nearwe \\
Strong instrumental & nearwe & nearwe & — \\
Weak nominative & nearwa & nearwe & nearwe \\
Weak genitive & nearwan & nearwan & nearwan \\
Weak dative & nearwan & nearwan & nearwan \\
Weak accusative & nearwan & nearwe & nearwan \\
\bottomrule
\end{tabular}
\par\endgroup
% course-panel:T25-P01:end
\end{minipage}\par
\par\addvspace{5pt}
% course-panel:T25-P02:begin
\par\begingroup
\fontsize{10.5}{13.2}\selectfont
\setlength{\LTleft}{0pt}\setlength{\LTright}{\fill}
\noindent
\begin{tabular}{@{}>{\raggedright\arraybackslash}p{36.136mm}>{\raggedright\arraybackslash}p{28.98mm}>{\raggedright\arraybackslash}p{23.852mm}>{\raggedright\arraybackslash}p{26.595mm}@{}}
\toprule
\textbf{Series /\allowbreak{} case} & \textbf{Masculine plural} & \textbf{Neuter plural} & \textbf{Feminine plural} \\
\midrule
Strong nominative & nearwe & nearu & nearwa \\
Strong genitive & nearora & nearora & nearora \\
Strong dative & nearwum & nearwum & nearwum \\
Strong accusative & nearwe & nearu & nearwa \\
Strong instrumental & — & — & — \\
Weak nominative & nearwan & nearwan & nearwan \\
Weak genitive & nearora & nearora & nearora \\
Weak dative & nearwum & nearwum & nearwum \\
Weak accusative & nearwan & nearwan & nearwan \\
\bottomrule
\end{tabular}
\par\endgroup
% course-panel:T25-P02:end
% mdtable:T25:end

\reserve{18\baselineskip}
% mdparagraph:P059:begin
\subsection*{Both}
% mdparagraph:P059:end

% teaching: NEW-U16-VOC-BEGEN
\par\begin{minipage}{\linewidth}\raggedright
% mdparagraph:P060:begin
\textbf{bēgen} — both. Its case and gender follow the noun or pronoun to which it refers. It has a plural declension, although it denotes two: it is not a dual inflection. It can accompany a noun or stand in apposition to a pronoun; \textbf{unc bǣm} means “to us both” and \textbf{inc bǣm}, “to you both”. The dative \textbf{bām} is also valid; neuter \textbf{bū} is an alternative to \textbf{bā}. The table selects one spelling for each position.
% mdparagraph:P060:end

% teaching: NEW-U16-NOM-BEGEN
% table: NEW-U16-NOM-BEGEN-DISPLAY
% mdtable:T26:begin
% course-table-style: 2026-10
% course-panel:T26-P01:begin
\par\begingroup
\fontsize{10.5}{13.2}\selectfont
\setlength{\LTleft}{0pt}\setlength{\LTright}{\fill}
\noindent
\begin{tabular}{@{}>{\raggedright\arraybackslash}p{22.0mm}>{\raggedright\arraybackslash}p{31.188mm}>{\raggedright\arraybackslash}p{31.188mm}>{\raggedright\arraybackslash}p{31.188mm}@{}}
\toprule
\textbf{Case} & \textbf{Masculine plural} & \textbf{Neuter plural} & \textbf{Feminine plural} \\
\midrule
Nominative & bēgen & bā & bā \\
Genitive & bēgra & bēgra & bēgra \\
Dative & bǣm & bǣm & bǣm \\
Accusative & bēgen & bā & bā \\
\bottomrule
\end{tabular}
\par\endgroup
% course-panel:T26-P01:end
% mdtable:T26:end
\end{minipage}\par
% teaching: NEW-U16-VOC-DOM
% mdparagraph:P061:begin
\textbf{dōm} — judgment; a judgment made in judging. Use the previously learned singular forms, including accusative \textbf{dōm}. This ordinary judgment sense is additional to the earlier religious and fate-related uses; no new abstract plural is required.
% mdparagraph:P061:end

\reserve{7\baselineskip}
% mdparagraph:P062:begin
\section*{2. Verbal systems and their basic uses}
% mdparagraph:P062:end

% teaching: NEW-U16-GRAMMAR
% mdparagraph:P063:begin
Each indicative table gives all six person-and-number positions in the present and preterite. The basic uses beside it state the objects and complements licensed here. A present can refer to a future event when the context supplies that viewpoint; English will does not by itself require willan. Commands, subjunctives and nonfinite forms are treated separately.\textsuperscript{3}
% mdparagraph:P063:end

% teaching: NEW-U16-PRINCIPAL-PARTS
\reserve{18\baselineskip}
% mdparagraph:P064:begin
\subsection*{Principal parts for reference}
% mdparagraph:P064:end

% mdparagraph:P065:begin
These dictionary forms show the stems from which the selected systems are formed. They remain available for reference. Productive finite forms are given in the complete tables that follow; productive nonfinite forms are identified in their own tables.
% mdparagraph:P065:end

% mdtable:T27:begin
\par\begingroup
\fontsize{10.5}{13.2}\selectfont
\setlength{\LTleft}{0pt}\setlength{\LTright}{\fill}
\noindent
\begin{tabular}{@{}>{\raggedright\arraybackslash}p{24.84mm}>{\raggedright\arraybackslash}p{36.02mm}>{\raggedright\arraybackslash}p{57.51mm}@{}}
\toprule
\textbf{Verb} & \textbf{Class} & \textbf{Principal parts: reference} \\
\midrule
ǣswician & weak II & ǣswician; ǣswicode; ǣswicod \\
blissian & weak II & blissian; blissode; blissod \\
wuldrian & weak II & wuldrian; wuldrode; wuldrod \\
plegian & weak II & plegian; plegode; plegod \\
behēafdian & weak II & behēafdian; behēafdode; behēafdod \\
cleopian & weak II & cleopian; cleopode; cleopod \\
ārīsan & strong I & ārīsan; ārās; ārison; ārisen \\
āceorfan & strong III & āceorfan; ācearf; ācurfon; ācorfen \\
forlēosan & strong II & forlēosan; forlēas; forluron; forloren \\
iernan & strong III & iernan; arn; urnon; urnen \\
ðwēan & strong VI contracted & ðwēan; ðwōh; ðwōgon; ðwegen \\
bewęndan & weak I & bewęndan; bewęnde; bewęnd \\
smierwan & weak I rw & smierwan; smierede; smiered \\
ālīesan & weak I & ālīesan; ālīesde; ālīesed \\
\bottomrule
\end{tabular}
\par\endgroup

% mdtable:T27:end
\reserve{7\baselineskip}
% mdparagraph:P066:begin
\subsection*{Weak verbs in -ian}
% mdparagraph:P066:end

% teaching: u16-aeswician-frame
\par\begin{minipage}{\linewidth}\raggedright
% mdparagraph:P067:begin
\textbf{ǣswician} — offend; cause someone to stumble. The person affected is accusative. In the exercise, hand or foot is a disjunctive singular subject and takes a singular verb; the meaning is moral offence, not simply physical injury.
% mdparagraph:P067:end

% mdparagraph:P068:begin
Weak II. The selected present endings are -ie, -ast, -að, -iað; the preterite uses -ode, -odest, -odon. Keep ǣswic- before these endings.
% mdparagraph:P068:end

% teaching: NEW-U16-PAR-AESWICIAN-IND
% table: NEW-U16-PAR-AESWICIAN-IND-DISPLAY
% mdtable:T28:begin
% course-table-style: 2026-10
% course-panel:T28-P01:begin
\par\begingroup
\fontsize{10.5}{13.2}\selectfont
\setlength{\LTleft}{0pt}\setlength{\LTright}{\fill}
\noindent
\begin{tabular}{@{}>{\raggedright\arraybackslash}p{39.459mm}>{\raggedright\arraybackslash}p{39.459mm}>{\raggedright\arraybackslash}p{39.459mm}@{}}
\toprule
\textbf{Person and number} & \textbf{Present} & \textbf{Preterite} \\
\midrule
1 sg. & ǣswicie & ǣswicode \\
2 sg. & ǣswicast & ǣswicodest \\
3 sg. & ǣswicað & ǣswicode \\
1–3 pl. & ǣswiciað & ǣswicodon \\
\bottomrule
\end{tabular}
\par\endgroup
% course-panel:T28-P01:end
% mdtable:T28:end
\end{minipage}\par
% teaching: u16-blissian-frame
\par\begin{minipage}{\linewidth}\raggedright
% mdparagraph:P069:begin
\textbf{blissian} — rejoice, intransitive. A companion may be introduced by mid with the dative. Ofer with the accusative introduces what one rejoices over; this use is explained again with comparison below.
% mdparagraph:P069:end

% mdparagraph:P070:begin
Weak II. Add the same selected endings to bliss-.
% mdparagraph:P070:end

% teaching: NEW-U16-PAR-BLISSIAN-IND
% table: NEW-U16-PAR-BLISSIAN-IND-DISPLAY
% mdtable:T29:begin
% course-table-style: 2026-10
% course-panel:T29-P01:begin
\par\begingroup
\fontsize{10.5}{13.2}\selectfont
\setlength{\LTleft}{0pt}\setlength{\LTright}{\fill}
\noindent
\begin{tabular}{@{}>{\raggedright\arraybackslash}p{39.459mm}>{\raggedright\arraybackslash}p{39.459mm}>{\raggedright\arraybackslash}p{39.459mm}@{}}
\toprule
\textbf{Person and number} & \textbf{Present} & \textbf{Preterite} \\
\midrule
1 sg. & blissie & blissode \\
2 sg. & blissast & blissodest \\
3 sg. & blissað & blissode \\
1–3 pl. & blissiað & blissodon \\
\bottomrule
\end{tabular}
\par\endgroup
% course-panel:T29-P01:end
% mdtable:T29:end
\end{minipage}\par
% teaching: u16-wuldrian-frame
\par\begin{minipage}{\linewidth}\raggedright
% mdparagraph:P071:begin
\textbf{wuldrian} — glorify, with an accusative person. A masculine speaker glorifying himself uses the accusative reflexive phrase mē selfne.
% mdparagraph:P071:end

% mdparagraph:P072:begin
Weak II. Add the selected endings to wuldr-.
% mdparagraph:P072:end

% teaching: NEW-U16-PAR-WULDRIAN-IND
% table: NEW-U16-PAR-WULDRIAN-IND-DISPLAY
% mdtable:T30:begin
% course-table-style: 2026-10
% course-panel:T30-P01:begin
\par\begingroup
\fontsize{10.5}{13.2}\selectfont
\setlength{\LTleft}{0pt}\setlength{\LTright}{\fill}
\noindent
\begin{tabular}{@{}>{\raggedright\arraybackslash}p{39.459mm}>{\raggedright\arraybackslash}p{39.459mm}>{\raggedright\arraybackslash}p{39.459mm}@{}}
\toprule
\textbf{Person and number} & \textbf{Present} & \textbf{Preterite} \\
\midrule
1 sg. & wuldrie & wuldrode \\
2 sg. & wuldrast & wuldrodest \\
3 sg. & wuldrað & wuldrode \\
1–3 pl. & wuldriað & wuldrodon \\
\bottomrule
\end{tabular}
\par\endgroup
% course-panel:T30-P01:end
% mdtable:T30:end
\end{minipage}\par
% teaching: u16-plegian-play
\par\begin{minipage}{\linewidth}\raggedright
% mdparagraph:P073:begin
\textbf{plegian} — play, intransitive. It can be modified by time and manner adverbs. The reading’s expressions for fellow players and the ball receive their own local help.
% mdparagraph:P073:end

% mdparagraph:P074:begin
Weak II. The selected stems are plegi- in the plural present and pleg- before -ode/\allowbreak{}-odon. Keep the selected spelling throughout the table.
% mdparagraph:P074:end

% teaching: NEW-U16-PAR-PLEGIAN-IND
% table: NEW-U16-PAR-PLEGIAN-IND-DISPLAY
% mdtable:T31:begin
% course-table-style: 2026-10
% course-panel:T31-P01:begin
\par\begingroup
\fontsize{10.5}{13.2}\selectfont
\setlength{\LTleft}{0pt}\setlength{\LTright}{\fill}
\noindent
\begin{tabular}{@{}>{\raggedright\arraybackslash}p{39.459mm}>{\raggedright\arraybackslash}p{39.459mm}>{\raggedright\arraybackslash}p{39.459mm}@{}}
\toprule
\textbf{Person and number} & \textbf{Present} & \textbf{Preterite} \\
\midrule
1 sg. & plegie & plegode \\
2 sg. & plegast & plegodest \\
3 sg. & plegað & plegode \\
1–3 pl. & plegiað & plegodon \\
\bottomrule
\end{tabular}
\par\endgroup
% course-panel:T31-P01:end
% mdtable:T31:end
\end{minipage}\par
% teaching: u16-beheafdian-frame
\par\begin{minipage}{\linewidth}\raggedright
% mdparagraph:P075:begin
\textbf{behēafdian} — behead, with an accusative person.
% mdparagraph:P075:end

% mdparagraph:P076:begin
Weak II. Retain the prefix and behēafd- stem in the selected forms.
% mdparagraph:P076:end

% teaching: NEW-U16-PAR-BEHEAFDIAN-IND
% table: NEW-U16-PAR-BEHEAFDIAN-IND-DISPLAY
% mdtable:T32:begin
% course-table-style: 2026-10
% course-panel:T32-P01:begin
\par\begingroup
\fontsize{10.5}{13.2}\selectfont
\setlength{\LTleft}{0pt}\setlength{\LTright}{\fill}
\noindent
\begin{tabular}{@{}>{\raggedright\arraybackslash}p{39.459mm}>{\raggedright\arraybackslash}p{39.459mm}>{\raggedright\arraybackslash}p{39.459mm}@{}}
\toprule
\textbf{Person and number} & \textbf{Present} & \textbf{Preterite} \\
\midrule
1 sg. & behēafdie & behēafdode \\
2 sg. & behēafdast & behēafdodest \\
3 sg. & behēafdað & behēafdode \\
1–3 pl. & behēafdiað & behēafdodon \\
\bottomrule
\end{tabular}
\par\endgroup
% course-panel:T32-P01:end
% mdtable:T32:end
\end{minipage}\par
% teaching: u16-cleopian-cry
\par\begin{minipage}{\linewidth}\raggedright
% mdparagraph:P077:begin
\textbf{cleopian} — cry aloud. The use here can introduce quoted words. The command forms are already familiar; the complete indicative is now learned.
% mdparagraph:P077:end

% mdparagraph:P078:begin
Weak II, with the same selected present and preterite endings as the preceding verbs.
% mdparagraph:P078:end

% teaching: NEW-U16-PAR-CLEOPIAN-IND
% table: NEW-U16-PAR-CLEOPIAN-IND-DISPLAY
% mdtable:T33:begin
% course-table-style: 2026-10
% course-panel:T33-P01:begin
\par\begingroup
\fontsize{10.5}{13.2}\selectfont
\setlength{\LTleft}{0pt}\setlength{\LTright}{\fill}
\noindent
\begin{tabular}{@{}>{\raggedright\arraybackslash}p{39.459mm}>{\raggedright\arraybackslash}p{39.459mm}>{\raggedright\arraybackslash}p{39.459mm}@{}}
\toprule
\textbf{Person and number} & \textbf{Present} & \textbf{Preterite} \\
\midrule
1 sg. & cleopie & cleopode \\
2 sg. & cleopast & cleopodest \\
3 sg. & cleopað & cleopode \\
1–3 pl. & cleopiað & cleopodon \\
\bottomrule
\end{tabular}
\par\endgroup
% course-panel:T33-P01:end
% mdtable:T33:end
\end{minipage}\par
\reserve{7\baselineskip}
% mdparagraph:P079:begin
\subsection*{Strong and contracted verbs}
% mdparagraph:P079:end

% mdparagraph:P080:begin
The ordinary strong present has singular endings -e, -st, -ð and plural -að. The preterite first and third singular have no added ending; second singular has -e on the plural stem, and all plural persons have -on. The stem changes and contractions stated below modify these ordinary combinations.
% mdparagraph:P080:end

\par\begin{minipage}{\linewidth}\raggedright
% mdparagraph:P081:begin
\textbf{ārīsan} — arise; get up, intransitive. This indicative system supplements the previously learned subjunctive.
% mdparagraph:P081:end

% mdparagraph:P082:begin
Strong I: ī /\allowbreak{} ā /\allowbreak{} i /\allowbreak{} i. Present -s plus -st or -ð gives the selected -st in both second and third singular. Preterite second singular uses the plural stem.
% mdparagraph:P082:end

% teaching: NEW-U16-PAR-ARISAN-IND
% table: NEW-U16-PAR-ARISAN-IND-DISPLAY
% mdtable:T34:begin
% course-table-style: 2026-10
% course-panel:T34-P01:begin
\par\begingroup
\fontsize{10.5}{13.2}\selectfont
\setlength{\LTleft}{0pt}\setlength{\LTright}{\fill}
\noindent
\begin{tabular}{@{}>{\raggedright\arraybackslash}p{39.459mm}>{\raggedright\arraybackslash}p{39.459mm}>{\raggedright\arraybackslash}p{39.459mm}@{}}
\toprule
\textbf{Person and number} & \textbf{Present} & \textbf{Preterite} \\
\midrule
1 sg. & ārīse & ārās \\
2 sg. & ārīst & ārise \\
3 sg. & ārīst & ārās \\
1–3 pl. & ārīsað & ārison \\
\bottomrule
\end{tabular}
\par\endgroup
% course-panel:T34-P01:end
% mdtable:T34:end
\end{minipage}\par
% teaching: u16-aceorfan-frame
\par\begin{minipage}{\linewidth}\raggedright
% mdparagraph:P083:begin
\textbf{āceorfan} — cut off, with an accusative object. The separate particle of can reinforce “off”; it has no governed noun after it in that use.
% mdparagraph:P083:end

% mdparagraph:P084:begin
Strong III: eo /\allowbreak{} ea /\allowbreak{} u /\allowbreak{} o. The second and third singular present have ie. The second singular preterite has the u-stem of the plural. Retain ā- throughout.
% mdparagraph:P084:end

% teaching: NEW-U16-PAR-ACEORFAN-IND
% table: NEW-U16-PAR-ACEORFAN-IND-DISPLAY
% mdtable:T35:begin
% course-table-style: 2026-10
% course-panel:T35-P01:begin
\par\begingroup
\fontsize{10.5}{13.2}\selectfont
\setlength{\LTleft}{0pt}\setlength{\LTright}{\fill}
\noindent
\begin{tabular}{@{}>{\raggedright\arraybackslash}p{39.459mm}>{\raggedright\arraybackslash}p{39.459mm}>{\raggedright\arraybackslash}p{39.459mm}@{}}
\toprule
\textbf{Person and number} & \textbf{Present} & \textbf{Preterite} \\
\midrule
1 sg. & āceorfe & ācearf \\
2 sg. & ācierfst & ācurfe \\
3 sg. & ācierfð & ācearf \\
1–3 pl. & āceorfað & ācurfon \\
\bottomrule
\end{tabular}
\par\endgroup
% course-panel:T35-P01:end
% mdtable:T35:end
\end{minipage}\par
% teaching: u16-forleosan-frame
\par\begin{minipage}{\linewidth}\raggedright
% mdparagraph:P085:begin
\textbf{forlēosan} — lose, with an accusative thing lost. Its participle can describe that thing as lost.
% mdparagraph:P085:end

% mdparagraph:P086:begin
Strong II: ēo /\allowbreak{} ēa /\allowbreak{} u /\allowbreak{} o. The preterite plural and past participle also change s to r. Both second and third singular present are forlīest; preterite second singular uses forlur-.
% mdparagraph:P086:end

% teaching: NEW-U16-PAR-FORLEOSAN-IND
% table: NEW-U16-PAR-FORLEOSAN-IND-DISPLAY
% mdtable:T36:begin
% course-table-style: 2026-10
% course-panel:T36-P01:begin
\par\begingroup
\fontsize{10.5}{13.2}\selectfont
\setlength{\LTleft}{0pt}\setlength{\LTright}{\fill}
\noindent
\begin{tabular}{@{}>{\raggedright\arraybackslash}p{39.459mm}>{\raggedright\arraybackslash}p{39.459mm}>{\raggedright\arraybackslash}p{39.459mm}@{}}
\toprule
\textbf{Person and number} & \textbf{Present} & \textbf{Preterite} \\
\midrule
1 sg. & forlēose & forlēas \\
2 sg. & forlīest & forlure \\
3 sg. & forlīest & forlēas \\
1–3 pl. & forlēosað & forluron \\
\bottomrule
\end{tabular}
\par\endgroup
% course-panel:T36-P01:end
% mdtable:T36:end
\end{minipage}\par
% teaching: u16-iernan-frame
\par\begin{minipage}{\linewidth}\raggedright
% mdparagraph:P087:begin
\textbf{iernan} — run, intransitive. An accusative actor with its plain infinitive can follow a verb of perception, as in the construction recalled below.
% mdparagraph:P087:end

% mdparagraph:P088:begin
Strong III. The selected stems are iern-, arn, urn- and urnen. Retain the position of r and the present ie spelling; do not replace it with a later irn- or yrn- series.
% mdparagraph:P088:end

% teaching: NEW-U16-PAR-IERNAN-IND
% table: NEW-U16-PAR-IERNAN-IND-DISPLAY
% mdtable:T37:begin
% course-table-style: 2026-10
% course-panel:T37-P01:begin
\par\begingroup
\fontsize{10.5}{13.2}\selectfont
\setlength{\LTleft}{0pt}\setlength{\LTright}{\fill}
\noindent
\begin{tabular}{@{}>{\raggedright\arraybackslash}p{39.459mm}>{\raggedright\arraybackslash}p{39.459mm}>{\raggedright\arraybackslash}p{39.459mm}@{}}
\toprule
\textbf{Person and number} & \textbf{Present} & \textbf{Preterite} \\
\midrule
1 sg. & ierne & arn \\
2 sg. & iernst & urne \\
3 sg. & iernð & arn \\
1–3 pl. & iernað & urnon \\
\bottomrule
\end{tabular}
\par\endgroup
% course-panel:T37-P01:end
% mdtable:T37:end
\end{minipage}\par
% teaching: u16-thwean-frame
\par\begin{minipage}{\linewidth}\raggedright
% mdparagraph:P089:begin
\textbf{ðwēan} — wash, with an accusative thing or body part. A possessor or recipient may retain the case independently required by the clause.
% mdparagraph:P089:end

% mdparagraph:P090:begin
Contracted strong VI, like slēan. The contracted present alternates ēa and ie. The selected preterite singular has h, while the second singular and plural have g; the past participle is ðwegen.
% mdparagraph:P090:end

% teaching: NEW-U16-PAR-THWEAN-IND
% table: NEW-U16-PAR-THWEAN-IND-DISPLAY
% mdtable:T38:begin
% course-table-style: 2026-10
% course-panel:T38-P01:begin
\par\begingroup
\fontsize{10.5}{13.2}\selectfont
\setlength{\LTleft}{0pt}\setlength{\LTright}{\fill}
\noindent
\begin{tabular}{@{}>{\raggedright\arraybackslash}p{39.459mm}>{\raggedright\arraybackslash}p{39.459mm}>{\raggedright\arraybackslash}p{39.459mm}@{}}
\toprule
\textbf{Person and number} & \textbf{Present} & \textbf{Preterite} \\
\midrule
1 sg. & ðwēa & ðwōh \\
2 sg. & ðwiehst & ðwōge \\
3 sg. & ðwiehð & ðwōh \\
1–3 pl. & ðwēað & ðwōgon \\
\bottomrule
\end{tabular}
\par\endgroup
% course-panel:T38-P01:end
% mdtable:T38:end
\end{minipage}\par
\reserve{7\baselineskip}
% mdparagraph:P091:begin
\subsection*{Further weak verbs}
% mdparagraph:P091:end

% mdparagraph:P092:begin
The ordinary weak-I present uses -e, -st, -ð and plural -að, with the contractions shown below. Its dental preterite stem takes singular -e, -est, -e and plural -on. Learn the selected dental stem from the reference principal part and the complete preterite column.
% mdparagraph:P092:end

% teaching: u16-bewendan-frame
\par\begin{minipage}{\linewidth}\raggedright
% mdparagraph:P093:begin
\textbf{bewęndan} — turn. In the subject-directed movement used here, dative him can refer back to the subject: the subject turns toward someone. This particular dative reflexive does not make all reflexives dative.
% mdparagraph:P093:end

% mdparagraph:P094:begin
Weak I with a d-stem, like sęndan. The present contracts to -ntst and -nt. The dental preterite contracts to bewęnde: present first singular and preterite first/\allowbreak{}third singular therefore coincide in spelling.
% mdparagraph:P094:end

% teaching: NEW-U16-PAR-BEWENDAN-IND
% table: NEW-U16-PAR-BEWENDAN-IND-DISPLAY
% mdtable:T39:begin
% course-table-style: 2026-10
% course-panel:T39-P01:begin
\par\begingroup
\fontsize{10.5}{13.2}\selectfont
\setlength{\LTleft}{0pt}\setlength{\LTright}{\fill}
\noindent
\begin{tabular}{@{}>{\raggedright\arraybackslash}p{39.459mm}>{\raggedright\arraybackslash}p{39.459mm}>{\raggedright\arraybackslash}p{39.459mm}@{}}
\toprule
\textbf{Person and number} & \textbf{Present} & \textbf{Preterite} \\
\midrule
1 sg. & bewęnde & bewęnde \\
2 sg. & bewęntst & bewęndest \\
3 sg. & bewęnt & bewęnde \\
1–3 pl. & bewęndað & bewęndon \\
\bottomrule
\end{tabular}
\par\endgroup
% course-panel:T39-P01:end
% mdtable:T39:end
\end{minipage}\par
% teaching: u16-smierwan-frame
\par\begin{minipage}{\linewidth}\raggedright
% mdparagraph:P095:begin
\textbf{smierwan} — anoint. The body part or person anointed is accusative; mid with the dative introduces the substance used.
% mdparagraph:P095:end

% mdparagraph:P096:begin
Weak I, like gierwan. Retain w in present first singular and plural; omit it in second/\allowbreak{}third singular and in the preterite. The selected series has ie, including smierede; the singular command is smiere.
% mdparagraph:P096:end

% teaching: NEW-U16-PAR-SMIERWAN-IND
% table: NEW-U16-PAR-SMIERWAN-IND-DISPLAY
% mdtable:T40:begin
% course-table-style: 2026-10
% course-panel:T40-P01:begin
\par\begingroup
\fontsize{10.5}{13.2}\selectfont
\setlength{\LTleft}{0pt}\setlength{\LTright}{\fill}
\noindent
\begin{tabular}{@{}>{\raggedright\arraybackslash}p{39.459mm}>{\raggedright\arraybackslash}p{39.459mm}>{\raggedright\arraybackslash}p{39.459mm}@{}}
\toprule
\textbf{Person and number} & \textbf{Present} & \textbf{Preterite} \\
\midrule
1 sg. & smierwe & smierede \\
2 sg. & smierest & smieredest \\
3 sg. & smiereð & smierede \\
1–3 pl. & smierwað & smieredon \\
\bottomrule
\end{tabular}
\par\endgroup
% course-panel:T40-P01:end
% mdtable:T40:end
\end{minipage}\par
\par\begin{minipage}{\linewidth}\raggedright
% mdparagraph:P097:begin
\textbf{ālīesan} — release; deliver, with an accusative person or thing. The meaning and commands continue the preceding unit; the indicative is now complete.
% mdparagraph:P097:end

% mdparagraph:P098:begin
Weak I. Continue the selected ālīesan, ālīesde, ālīesed system. In the second and third singular present, contraction after s gives -st. Do not substitute a strong-verb system.
% mdparagraph:P098:end

% teaching: NEW-U16-PAR-ALIESAN-IND
% table: NEW-U16-PAR-ALIESAN-IND-DISPLAY
% mdtable:T41:begin
% course-table-style: 2026-10
% course-panel:T41-P01:begin
\par\begingroup
\fontsize{10.5}{13.2}\selectfont
\setlength{\LTleft}{0pt}\setlength{\LTright}{\fill}
\noindent
\begin{tabular}{@{}>{\raggedright\arraybackslash}p{39.459mm}>{\raggedright\arraybackslash}p{39.459mm}>{\raggedright\arraybackslash}p{39.459mm}@{}}
\toprule
\textbf{Person and number} & \textbf{Present} & \textbf{Preterite} \\
\midrule
1 sg. & ālīese & ālīesde \\
2 sg. & ālīest & ālīesdest \\
3 sg. & ālīest & ālīesde \\
1–3 pl. & ālīesað & ālīesdon \\
\bottomrule
\end{tabular}
\par\endgroup
% course-panel:T41-P01:end
% mdtable:T41:end
\end{minipage}\par
\reserve{18\baselineskip}
% mdparagraph:P099:begin
\subsection*{Commands}
% mdparagraph:P099:end

% mdparagraph:P100:begin
The first form addresses one person; the second addresses more than one. Use ne before a negative command. The bare-stem commands preserve their own stem changes; weak-II blissian instead has singular blissa, and smierwan has singular smiere. These are commands, not present indicative forms inferred from a matching spelling.
% mdparagraph:P100:end

% teaching: NEW-U16-PAR-ACEORFAN-IMP
% table: NEW-U16-PAR-ACEORFAN-IMP-DISPLAY
% mdtable:T42:begin
% course-table-style: 2026-10
% course-panel:T42-P01:begin
\par\begingroup
\fontsize{10.5}{13.2}\selectfont
\setlength{\LTleft}{0pt}\setlength{\LTright}{\fill}
\noindent
\begin{tabular}{@{}>{\raggedright\arraybackslash}p{63.342mm}>{\raggedright\arraybackslash}p{57.846mm}@{}}
\toprule
\textbf{Person and number} & \textbf{Imperative} \\
\midrule
2 sg. & āceorf \\
2 pl. & āceorfað \\
\bottomrule
\end{tabular}
\par\endgroup
% course-panel:T42-P01:end
% mdtable:T42:end
% teaching: NEW-U16-PAR-BLISSIAN-IMP
% table: NEW-U16-PAR-BLISSIAN-IMP-DISPLAY
% mdtable:T43:begin
% course-table-style: 2026-10
% course-panel:T43-P01:begin
\par\begingroup
\fontsize{10.5}{13.2}\selectfont
\setlength{\LTleft}{0pt}\setlength{\LTright}{\fill}
\noindent
\begin{tabular}{@{}>{\raggedright\arraybackslash}p{63.342mm}>{\raggedright\arraybackslash}p{57.846mm}@{}}
\toprule
\textbf{Person and number} & \textbf{Imperative} \\
\midrule
2 sg. & blissa \\
2 pl. & blissiað \\
\bottomrule
\end{tabular}
\par\endgroup
% course-panel:T43-P01:end
% mdtable:T43:end
% teaching: NEW-U16-PAR-SMIERWAN-IMP
% table: NEW-U16-PAR-SMIERWAN-IMP-DISPLAY
% mdtable:T44:begin
% course-table-style: 2026-10
% course-panel:T44-P01:begin
\par\begingroup
\fontsize{10.5}{13.2}\selectfont
\setlength{\LTleft}{0pt}\setlength{\LTright}{\fill}
\noindent
\begin{tabular}{@{}>{\raggedright\arraybackslash}p{63.342mm}>{\raggedright\arraybackslash}p{57.846mm}@{}}
\toprule
\textbf{Person and number} & \textbf{Imperative} \\
\midrule
2 sg. & smiere \\
2 pl. & smierwað \\
\bottomrule
\end{tabular}
\par\endgroup
% course-panel:T44-P01:end
% mdtable:T44:end
% teaching: NEW-U16-PAR-ARISAN-IMP
% table: NEW-U16-PAR-ARISAN-IMP-DISPLAY
% mdtable:T45:begin
% course-table-style: 2026-10
% course-panel:T45-P01:begin
\par\begingroup
\fontsize{10.5}{13.2}\selectfont
\setlength{\LTleft}{0pt}\setlength{\LTright}{\fill}
\noindent
\begin{tabular}{@{}>{\raggedright\arraybackslash}p{63.342mm}>{\raggedright\arraybackslash}p{57.846mm}@{}}
\toprule
\textbf{Person and number} & \textbf{Imperative} \\
\midrule
2 sg. & ārīs \\
2 pl. & ārīsað \\
\bottomrule
\end{tabular}
\par\endgroup
% course-panel:T45-P01:end
% mdtable:T45:end
\par\begin{minipage}{\linewidth}\raggedright
% mdparagraph:P101:begin
\textbf{forlǣtan} — leave; let go, with the previously learned objects and sense.
% mdparagraph:P101:end

% teaching: NEW-U16-PAR-FORLAETAN-IMP
% table: NEW-U16-PAR-FORLAETAN-IMP-DISPLAY
% mdtable:T46:begin
% course-table-style: 2026-10
% course-panel:T46-P01:begin
\par\begingroup
\fontsize{10.5}{13.2}\selectfont
\setlength{\LTleft}{0pt}\setlength{\LTright}{\fill}
\noindent
\begin{tabular}{@{}>{\raggedright\arraybackslash}p{63.342mm}>{\raggedright\arraybackslash}p{57.846mm}@{}}
\toprule
\textbf{Person and number} & \textbf{Imperative} \\
\midrule
2 sg. & forlǣt \\
2 pl. & forlǣtað \\
\bottomrule
\end{tabular}
\par\endgroup
% course-panel:T46-P01:end
% mdtable:T46:end
\end{minipage}\par
\par\begin{minipage}{\linewidth}\raggedright
% mdparagraph:P102:begin
\textbf{geswīcan} — cease. It may be used without an expressed complement, including with a wind as subject or in an addressed command. Its genitive complement remains a separate use.
% mdparagraph:P102:end

% teaching: NEW-U16-PAR-GESWICAN-IMP
% table: NEW-U16-PAR-GESWICAN-IMP-DISPLAY
% mdtable:T47:begin
% course-table-style: 2026-10
% course-panel:T47-P01:begin
\par\begingroup
\fontsize{10.5}{13.2}\selectfont
\setlength{\LTleft}{0pt}\setlength{\LTright}{\fill}
\noindent
\begin{tabular}{@{}>{\raggedright\arraybackslash}p{63.342mm}>{\raggedright\arraybackslash}p{57.846mm}@{}}
\toprule
\textbf{Person and number} & \textbf{Imperative} \\
\midrule
2 sg. & geswīc \\
2 pl. & geswīcað \\
\bottomrule
\end{tabular}
\par\endgroup
% course-panel:T47-P01:end
% mdtable:T47:end
\end{minipage}\par
\par\begin{minipage}{\linewidth}\raggedright
% mdparagraph:P103:begin
\textbf{ęfthweorfan} — return; turn back. Retain the prefix in both commands.
% mdparagraph:P103:end

% teaching: NEW-U16-PAR-EFTHWEORFAN-IMP
% table: NEW-U16-PAR-EFTHWEORFAN-IMP-DISPLAY
% mdtable:T48:begin
% course-table-style: 2026-10
% course-panel:T48-P01:begin
\par\begingroup
\fontsize{10.5}{13.2}\selectfont
\setlength{\LTleft}{0pt}\setlength{\LTright}{\fill}
\noindent
\begin{tabular}{@{}>{\raggedright\arraybackslash}p{60.456mm}>{\raggedright\arraybackslash}p{60.732mm}@{}}
\toprule
\textbf{Person and number} & \textbf{Imperative} \\
\midrule
2 sg. & ęfthweorf \\
2 pl. & ęfthweorfað \\
\bottomrule
\end{tabular}
\par\endgroup
% course-panel:T48-P01:end
% mdtable:T48:end
\end{minipage}\par
\par\begin{minipage}{\linewidth}\raggedright
% mdparagraph:P104:begin
\textbf{dēman} — judge. A judgment itself may be expressed as an accusative object, as in the instruction to judge a right judgment.
% mdparagraph:P104:end

% teaching: NEW-U16-PAR-DEMAN-IMP
% table: NEW-U16-PAR-DEMAN-IMP-DISPLAY
% mdtable:T49:begin
% course-table-style: 2026-10
% course-panel:T49-P01:begin
\par\begingroup
\fontsize{10.5}{13.2}\selectfont
\setlength{\LTleft}{0pt}\setlength{\LTright}{\fill}
\noindent
\begin{tabular}{@{}>{\raggedright\arraybackslash}p{63.342mm}>{\raggedright\arraybackslash}p{57.846mm}@{}}
\toprule
\textbf{Person and number} & \textbf{Imperative} \\
\midrule
2 sg. & dēm \\
2 pl. & dēmað \\
\bottomrule
\end{tabular}
\par\endgroup
% course-panel:T49-P01:end
% mdtable:T49:end
\end{minipage}\par
% mdparagraph:P105:begin
\textbf{Immediately before gē}, use the shortened plural command \textbf{dēme gē}. Without that following pronoun the ordinary plural remains \textbf{dēmað}. This positional exception is required in the first clause of Cook 36.1; it does not replace the ordinary pair or introduce a new general command table.\textsuperscript{4}
% mdparagraph:P105:end

\reserve{18\baselineskip}
% mdparagraph:P106:begin
\subsection*{Present subjunctives: expectation and an unspecified provider}
% mdparagraph:P106:end

% mdparagraph:P107:begin
Findan “find” and fēdan “feed” retain their previously learned indicative systems. The complete present subjunctive has -e in every singular person and -en in every plural person. The six positions remain distinct even when the forms coincide.
% mdparagraph:P107:end

% teaching: NEW-U16-PAR-FINDAN-PRES-SUBJ
% table: NEW-U16-PAR-FINDAN-PRES-SUBJ-DISPLAY
% mdtable:T50:begin
% course-table-style: 2026-10
% course-panel:T50-P01:begin
\par\begingroup
\fontsize{10.5}{13.2}\selectfont
\setlength{\LTleft}{0pt}\setlength{\LTright}{\fill}
\noindent
\begin{tabular}{@{}>{\raggedright\arraybackslash}p{63.342mm}>{\raggedright\arraybackslash}p{57.846mm}@{}}
\toprule
\textbf{Person and number} & \textbf{Present} \\
\midrule
1 sg. & finde \\
2 sg. & finde \\
3 sg. & finde \\
1–3 pl. & finden \\
\bottomrule
\end{tabular}
\par\endgroup
% course-panel:T50-P01:end
% mdtable:T50:end
% teaching: NEW-U16-PAR-FEDAN-PRES-SUBJ
% table: NEW-U16-PAR-FEDAN-PRES-SUBJ-DISPLAY
% mdtable:T51:begin
% course-table-style: 2026-10
% course-panel:T51-P01:begin
\par\begingroup
\fontsize{10.5}{13.2}\selectfont
\setlength{\LTleft}{0pt}\setlength{\LTright}{\fill}
\noindent
\begin{tabular}{@{}>{\raggedright\arraybackslash}p{63.342mm}>{\raggedright\arraybackslash}p{57.846mm}@{}}
\toprule
\textbf{Person and number} & \textbf{Present} \\
\midrule
1 sg. & fēde \\
2 sg. & fēde \\
3 sg. & fēde \\
1–3 pl. & fēden \\
\bottomrule
\end{tabular}
\par\endgroup
% course-panel:T51-P01:end
% mdtable:T51:end
% mdparagraph:P108:begin
\textbf{Wēn is ðæt} with a present subjunctive expresses the bounded expectation “perhaps; there is a prospect that …”. In the hypothetical conditional “if you find no one”, use \textbf{finde}; a relative clause identifying a possible, as-yet-unfound provider uses \textbf{fēde}. These choices concern the status of the situation, not a mechanical conversion of an English future.
% mdparagraph:P108:end

% mdparagraph:P109:begin
Keep the previously selected preterites of findan distinct: first singular \textbf{fand}, third singular \textbf{funde}. A spelling shared with another person or mood must be interpreted from its clause.
% mdparagraph:P109:end

\reserve{7\baselineskip}
% mdparagraph:P110:begin
\subsection*{Participles and infinitives}
% mdparagraph:P110:end

% mdparagraph:P111:begin
The following pairs give the productive participial bases. These present participles add -ende to the present stem and describe ongoing action. Learn each past-participle stem and prefix as shown. The past participle can enter the previously learned passive or perfect constructions; in a passive, the affected participant is the subject. These bases do not themselves supply every possible adjective inflection.
% mdparagraph:P111:end

\reserve{10\baselineskip}
% mdparagraph:P112:begin
Forlēosan retains the r-stem in its past participle. In the strong neuter accusative singular, \textbf{forloren} can modify a lost thing; this exact attributive position is also learned here.
% mdparagraph:P112:end

% teaching: NEW-U16-PAR-FORLEOSAN-PART
% table: NEW-U16-PAR-FORLEOSAN-PART-DISPLAY
% mdtable:T52:begin
\par\begingroup
\fontsize{10.5}{13.2}\selectfont
\setlength{\LTleft}{0pt}\setlength{\LTright}{\fill}
\noindent
\begin{tabular}{@{}>{\raggedright\arraybackslash}p{69.12mm}>{\raggedright\arraybackslash}p{52.07mm}@{}}
\toprule
\textbf{Position} & \textbf{Form} \\
\midrule
participle, present & forlēosende \\
participle, past & forloren \\
\bottomrule
\end{tabular}
\par\endgroup

% mdtable:T52:end
% teaching: u16-laedan-passive-carry
\reserve{10\baselineskip}
% mdparagraph:P113:begin
Lǣdan here includes “carry; take a person”. A passive makes the person carried its subject. The endpoint in the assigned carrying event uses ōð with the accusative.
% mdparagraph:P113:end

% teaching: NEW-U16-PAR-LAEDAN-PART
% table: NEW-U16-PAR-LAEDAN-PART-DISPLAY
% mdtable:T53:begin
\par\begingroup
\fontsize{10.5}{13.2}\selectfont
\setlength{\LTleft}{0pt}\setlength{\LTright}{\fill}
\noindent
\begin{tabular}{@{}>{\raggedright\arraybackslash}p{71.92mm}>{\raggedright\arraybackslash}p{49.26mm}@{}}
\toprule
\textbf{Position} & \textbf{Form} \\
\midrule
participle, present & lǣdende \\
participle, past & lǣded \\
\bottomrule
\end{tabular}
\par\endgroup

% mdtable:T53:end
% teaching: u16-gefyllan-utterance
\reserve{10\baselineskip}
% mdparagraph:P114:begin
Gefyllan here means “fulfil a spoken prediction”. That prediction is accusative in an active clause and becomes the subject in a passive. This is distinct from physically filling a container.
% mdparagraph:P114:end

% teaching: NEW-U16-PAR-GEFYLLAN-PART
% table: NEW-U16-PAR-GEFYLLAN-PART-DISPLAY
% mdtable:T54:begin
\par\begingroup
\fontsize{10.5}{13.2}\selectfont
\setlength{\LTleft}{0pt}\setlength{\LTright}{\fill}
\noindent
\begin{tabular}{@{}>{\raggedright\arraybackslash}p{71.92mm}>{\raggedright\arraybackslash}p{49.26mm}@{}}
\toprule
\textbf{Position} & \textbf{Form} \\
\midrule
participle, present & gefyllende \\
participle, past & gefylled \\
\bottomrule
\end{tabular}
\par\endgroup

% mdtable:T54:end
% teaching: u16-forgiefan-forgive
\reserve{11\baselineskip}
% mdparagraph:P115:begin
Forgiefan here means “forgive”: the sin or offence is accusative, the person forgiven dative. In a passive the sin becomes nominative and the person remains dative. The uninflected past participle is \textbf{forgiefen} in the assigned predicative use.
% mdparagraph:P115:end

% teaching: NEW-U16-PAR-FORGIEFAN-PART
% table: NEW-U16-PAR-FORGIEFAN-PART-DISPLAY
% mdtable:T55:begin
\par\begingroup
\fontsize{10.5}{13.2}\selectfont
\setlength{\LTleft}{0pt}\setlength{\LTright}{\fill}
\noindent
\begin{tabular}{@{}>{\raggedright\arraybackslash}p{69.12mm}>{\raggedright\arraybackslash}p{52.07mm}@{}}
\toprule
\textbf{Position} & \textbf{Form} \\
\midrule
participle, present & forgiefende \\
participle, past & forgiefen \\
\bottomrule
\end{tabular}
\par\endgroup

% mdtable:T55:end
\reserve{13\baselineskip}
% mdparagraph:P116:begin
The plain infinitive and the inflected infinitive after tō are separate forms. A perception construction takes the plain infinitive and an accusative actor. The previously learned purpose construction takes tō and the inflected form. A new meaning such as “is to come” still requires separate local help; it does not follow merely from knowing the ending.
% mdparagraph:P116:end

% teaching: NEW-U16-PAR-IERNAN-INF
% table: NEW-U16-PAR-IERNAN-INF-DISPLAY
% mdtable:T56:begin
\par\begingroup
\fontsize{10.5}{13.2}\selectfont
\setlength{\LTleft}{0pt}\setlength{\LTright}{\fill}
\noindent
\begin{tabular}{@{}>{\raggedright\arraybackslash}p{68.04mm}>{\raggedright\arraybackslash}p{53.15mm}@{}}
\toprule
\textbf{Position} & \textbf{Form} \\
\midrule
infinitive & iernan \\
infinitive, dat & iernanne \\
\bottomrule
\end{tabular}
\par\endgroup

% mdtable:T56:end
% teaching: NEW-U16-PAR-SENDAN-INF
% table: NEW-U16-PAR-SENDAN-INF-DISPLAY
% mdtable:T57:begin
\par\begingroup
\fontsize{10.5}{13.2}\selectfont
\setlength{\LTleft}{0pt}\setlength{\LTright}{\fill}
\noindent
\begin{tabular}{@{}>{\raggedright\arraybackslash}p{68.04mm}>{\raggedright\arraybackslash}p{53.15mm}@{}}
\toprule
\textbf{Position} & \textbf{Form} \\
\midrule
infinitive & sęndan \\
infinitive, dat & sęndanne \\
\bottomrule
\end{tabular}
\par\endgroup

% mdtable:T57:end
\reserve{7\baselineskip}
% mdparagraph:P117:begin
\section*{3. Clause structure, government and viewpoint}
% mdparagraph:P117:end

\reserve{7\baselineskip}
% mdparagraph:P118:begin
\subsection*{Reference and finite content}
% mdparagraph:P118:end

% mdparagraph:P119:begin
A pronoun does not automatically refer to the nearest noun. Establish the speaker, addressee and continuing participant from the whole passage. A pronoun’s ending can establish its role without uniquely establishing its referent. When a new composition requires an unambiguous reference, an explicit noun phrase can resolve it.
% mdparagraph:P119:end

% teaching: NEW-U16-EXAMPLE-01
\reserve{8\baselineskip}
% example: NEW-U16-EXAMPLE-01
% mdparagraph:P120:begin
Se fæder sægde: ‘Ic eom hēr.’ Se fæder sægde ðæt hē ðǣr wǣre.
% mdparagraph:P120:end

% mdparagraph:P121:begin
The father said, 'I am here.' The father said that he was there.
% mdparagraph:P121:end

% mdparagraph:P122:begin
The quotation uses the father’s I and here. The report uses third-person hē and, under the stated change of reporting place, ðǣr. The subjunctive wǣre belongs to reported content; it does not mean that every English past verb must be replaced by an Old English subjunctive.
% mdparagraph:P122:end

% teaching: NEW-U16-EXAMPLE-02
\reserve{8\baselineskip}
% example: NEW-U16-EXAMPLE-02
% mdparagraph:P123:begin
Hē geseah ðone mann cuman. Hē wiste ðæt se mann cōme.
% mdparagraph:P123:end

% mdparagraph:P124:begin
He saw the man come. He knew that the man came.
% mdparagraph:P124:end

% mdparagraph:P125:begin
With perception and a plain infinitive, the actor of the perceived action is accusative. A finite ðæt-clause has its own nominative subject and finite verb. Some finite spellings can represent more than one mood; identify only what the ending and context together warrant.
% mdparagraph:P125:end

% mdparagraph:P126:begin
An earlier speech event does not make every situation in its content earlier still. A reported state can be simultaneous with the saying; habban with a past participle can mark an already completed event; willan with an infinitive can express an intention whose action is still ahead. Interpret the whole verb group and the stated context. Directly quoted now or today retains the quoted speaker’s viewpoint.
% mdparagraph:P126:end

\reserve{7\baselineskip}
% mdparagraph:P127:begin
\subsection*{Knowing, action comparison and causation}
% mdparagraph:P127:end

% teaching: u16-cunnan-person
% teaching: u16-witan-person
% mdparagraph:P128:begin
\textbf{Cunnan} can mean “know; be acquainted with” a person and take that person in the accusative. Its earlier use with a plain infinitive expresses knowing how to do something. \textbf{Witan} can introduce factual content or a question, but can also take a personal accusative object. Thus the contrast is useful without making persons and facts an exclusive division between the two verbs.\textsuperscript{5}
% mdparagraph:P128:end

% teaching: NEW-U16-EXAMPLE-03
\reserve{8\baselineskip}
% example: NEW-U16-EXAMPLE-03
% mdparagraph:P129:begin
Gē cunnon ðone cyning; gē cunnon singan.
% mdparagraph:P129:end

% mdparagraph:P130:begin
You know the king; you know how to sing.
% mdparagraph:P130:end

% teaching: u16-don-substitute
% teaching: u16-don-finite-result
% mdparagraph:P131:begin
\textbf{Dōn} can repeat an action already expressed. In an action comparison, \textbf{bęt} “better” modifies the action and \textbf{ðonne} introduces the comparison clause. Dōn then carries the tense, person and number of its own subject. It is not an auxiliary inserted to imitate modern English questions or negatives.
% mdparagraph:P131:end

% teaching: NEW-U16-EXAMPLE-04
\reserve{8\baselineskip}
% example: NEW-U16-EXAMPLE-04
% mdparagraph:P132:begin
Hē plegode wel, and hīe dydon swā.
% mdparagraph:P132:end

% mdparagraph:P133:begin
He played well, and they did so.
% mdparagraph:P133:end

% mdparagraph:P134:begin
A separate causative pattern is \textbf{dōn + ðæt + finite clause}, “cause that; make it be that”. The subject inside that clause remains nominative. With dual \textbf{git}, use a plural finite verb, such as \textbf{bēoð}; a noun describing both participants is plural. Do not turn the dual subject into an accusative object of dōn in this pattern.
% mdparagraph:P134:end

\reserve{7\baselineskip}
% mdparagraph:P135:begin
\subsection*{Government and lexical meaning}
% mdparagraph:P135:end

% mdparagraph:P136:begin
A case after a verb belongs to that verb’s construction. It is not determined by the nearest English preposition.
% mdparagraph:P136:end

% mdparagraph:P137:begin
\textbf{Biddan} takes the requested thing in the genitive; the person asked can simultaneously be accusative. Interrogative \textbf{hwæs} asks for that requested thing. \textbf{Gīeman} “take care of” and \textbf{geswīcan} “cease from” retain their learned genitive government. \textbf{Wēnan} “expect” now also takes a genitive thing expected, including a genitive negative or a coordinated exception. Its earlier finite content-clause use does not replace this pattern.
% mdparagraph:P137:end

% teaching: u16-halsian-oath
% mdparagraph:P138:begin
\textbf{Hālsian} “conjure; solemnly charge” takes the person charged in the accusative. \textbf{Ðurh} with an accusative names the oath standard, and a requested \textbf{ðæt}-clause follows. An uncertain embedded identification uses a subjunctive. This person is not an ordinary dative recipient.
% mdparagraph:P138:end

% teaching: u16-healdan-keep
% teaching: u16-weorthan-arise
% mdparagraph:P139:begin
\textbf{Healdan} now includes “keep; observe” a saying, with accusative sprǣc. \textbf{Weorðan} can mean “arise; come about” with an event as its subject. If a calm comes about and no participle is present, this is an intransitive lexical use, not the passive auxiliary.
% mdparagraph:P139:end

% mdparagraph:P140:begin
With \textbf{bewęndan}, the dative reflexive can identify the person who turns toward someone. With \textbf{gān}, the genitive route pattern \textbf{ðæs weges ðe …} means “by the way by which …”. Neither construction should be generalized to every motion verb. In Cook 36.10, \textbf{āsęndan} has no expressed patient: understand “sent someone”, followed by the coordinated action.
% mdparagraph:P140:end

% teaching: u16-ofer-rejoicing
% teaching: u16-degree-comparison
% teaching: u16-of-particle
% mdparagraph:P141:begin
With \textbf{blissian}, \textbf{ofer + accusative} names the object of rejoicing. \textbf{Swīðor} here compares its degree, “more”, and \textbf{ðonne} introduces the second object of comparison, with its own ofer phrase. It is not the earlier “rather” meaning. The separate particle \textbf{of}, “off”, has no following governed complement; prepositional \textbf{of} still governs the dative in its own uses.
% mdparagraph:P141:end

\reserve{7\baselineskip}
% mdparagraph:P142:begin
\subsection*{Direction, relation and time words}
% mdparagraph:P142:end

% teaching: NEW-U16-VOC-HWANON
% teaching: NEW-U16-VOC-HWIDER
% teaching: NEW-U16-VOC-TOGAEDERE
% teaching: NEW-U16-VOC-GEOND
% teaching: NEW-U16-VOC-INNAN
% teaching: NEW-U16-VOC-SITHTHAN
% teaching: NEW-U16-VOC-THRIWA
% teaching: NEW-U16-VOC-SONA
% teaching: NEW-U16-VOC-AEFTER
% teaching: u16-aefter-follow
% teaching: NEW-U16-VOC-INTO
% teaching: NEW-U16-VOC-BE
% teaching: NEW-U16-VOC-THURH
% teaching: NEW-U16-VOC-FOR-THON-THE
% mdtable:T58:begin
\par\begingroup
\fontsize{10.5}{13.2}\selectfont
\setlength{\LTleft}{0pt}\setlength{\LTright}{\fill}
\noindent
\begin{tabular}{@{}>{\raggedright\arraybackslash}p{18.5mm}>{\raggedright\arraybackslash}p{26.31mm}>{\raggedright\arraybackslash}p{73.57mm}@{}}
\toprule
\textbf{Word} & \textbf{Meaning} & \textbf{Use} \\
\midrule
hwanon & whence; from where & Interrogative adverb of origin. \\
hwider & whither; to where & Interrogative adverb of destination. \\
tōgædere & together & Adverb. \\
geond & along; through & Preposition with accusative. \\
innan & within; inside & Adverb in the assigned use. \\
siððan & since; after; afterwards & As a conjunction, introduces a temporal clause; as an adverb, means afterwards. \\
ðrīwa & thrice; three times & Adverb. \\
sōna & immediately; at once & Adverb; a complete as-soon-as construction is separately supplied where required. \\
æfter & after; following; according to & Dative in the three stated uses: time after an event, following a person, or conformity to something. \\
intō & into & Destination with dative in this selected use. \\
be & according to; concerning & Dative in both senses. The reading’s “by the hand” receives separate local help. \\
ðurh & through; by way of a person & Accusative for the mediating person. Oath use is separately identified; source-local means or evidence receives help. \\
for-ðon & therefore; for that reason & Adverbial consequence here. Do not expand it into a because-clause merely from the related longer expression. \\
\bottomrule
\end{tabular}
\par\endgroup

% mdtable:T58:end
% mdparagraph:P143:begin
The following exercises reuse the previously learned masculine \textbf{sǣ}. Preserve the appropriate masculine article or demonstrative with it. \textbf{Bōc} retains its learned feminine plural \textbf{bēc}. Source variants in the reading are supplied locally and do not replace the forms learned for new composition.
% mdparagraph:P143:end

\reserve{7\baselineskip}
% mdparagraph:P144:begin
\section*{4. Form and clause practice}
% mdparagraph:P144:end

% teaching: NEW-U16-DIRECTION-HEADWORDS
% mdparagraph:P145:begin
The dictionary forms printed with the form exercises are given references. Produce the requested inflections from the learned systems. If a requested position has the same form as its given headword, identify that position; it is not a test of recalling an unseen word. Explain the requested contrasts in English.
% mdparagraph:P145:end

\reserve{5\baselineskip}
\par\begin{minipage}{\linewidth}\raggedright
% mdparagraph:P146:begin
\textbf{F01}
% mdparagraph:P146:end

% item: NEW-U16-F01
% mdparagraph:P147:begin
Give the selected forms requested below. For each verb, explain the stem changes linking your three answers. Retain the prefixes and the selected spelling.
% mdparagraph:P147:end

% mdparagraph:P148:begin
(1) āceorfan
% mdparagraph:P148:end

% mdparagraph:P149:begin
(2) forlēosan
% mdparagraph:P149:end

% mdparagraph:P150:begin
(3) ðwēan
% mdparagraph:P150:end

% mdparagraph:P151:begin
(4) iernan
% mdparagraph:P151:end

% mdparagraph:P152:begin
\textbf{Supplied words.} These forms and meanings remain available for this task.
% mdparagraph:P152:end

% mdtable:T59:begin
\par\begingroup
\fontsize{10.5}{13.2}\selectfont
\setlength{\LTleft}{0pt}\setlength{\LTright}{\fill}
\noindent
\begin{tabular}{@{}>{\raggedright\arraybackslash}p{33.87mm}>{\raggedright\arraybackslash}p{33.87mm}>{\raggedright\arraybackslash}p{50.64mm}@{}}
\toprule
\textbf{Form} & \textbf{Meaning} & \textbf{Role} \\
\midrule
āceorfan & cut off & given dictionary form \\
forlēosan & lose & given dictionary form \\
ðwēan & wash & given dictionary form \\
\bottomrule
\end{tabular}
\par\endgroup

% mdtable:T59:end
% output: NEW-U16-F01-O001
% mdparagraph:P153:begin
(1) āceorfan: present indicative, you singular: \_\_\_\_\_\_\_\_\_\_\_\_\_\_\_\_\_\_\_\_
% mdparagraph:P153:end

% output: NEW-U16-F01-O002
% mdparagraph:P154:begin
(2) āceorfan: preterite indicative, you singular: \_\_\_\_\_\_\_\_\_\_\_\_\_\_\_\_\_\_\_\_
% mdparagraph:P154:end

% output: NEW-U16-F01-O003
% mdparagraph:P155:begin
(3) āceorfan: preterite indicative, they: \_\_\_\_\_\_\_\_\_\_\_\_\_\_\_\_\_\_\_\_
% mdparagraph:P155:end

% output: NEW-U16-F01-O004
% mdparagraph:P156:begin
(4) āceorfan: explain the stem relations: \_\_\_\_\_\_\_\_\_\_\_\_\_\_\_\_\_\_\_\_
% mdparagraph:P156:end

% output: NEW-U16-F01-O005
% mdparagraph:P157:begin
(5) forlēosan: present indicative, he: \_\_\_\_\_\_\_\_\_\_\_\_\_\_\_\_\_\_\_\_
% mdparagraph:P157:end

% output: NEW-U16-F01-O006
% mdparagraph:P158:begin
(6) forlēosan: preterite indicative, you singular: \_\_\_\_\_\_\_\_\_\_\_\_\_\_\_\_\_\_\_\_
% mdparagraph:P158:end

% output: NEW-U16-F01-O007
% mdparagraph:P159:begin
(7) forlēosan: preterite indicative, they: \_\_\_\_\_\_\_\_\_\_\_\_\_\_\_\_\_\_\_\_
% mdparagraph:P159:end

% output: NEW-U16-F01-O008
% mdparagraph:P160:begin
(8) forlēosan: explain the stem relations: \_\_\_\_\_\_\_\_\_\_\_\_\_\_\_\_\_\_\_\_
% mdparagraph:P160:end

% output: NEW-U16-F01-O009
% mdparagraph:P161:begin
(9) ðwēan: present indicative, he: \_\_\_\_\_\_\_\_\_\_\_\_\_\_\_\_\_\_\_\_
% mdparagraph:P161:end

% output: NEW-U16-F01-O010
% mdparagraph:P162:begin
(10) ðwēan: preterite indicative, you singular: \_\_\_\_\_\_\_\_\_\_\_\_\_\_\_\_\_\_\_\_
% mdparagraph:P162:end

% output: NEW-U16-F01-O011
% mdparagraph:P163:begin
(11) ðwēan: preterite indicative, they: \_\_\_\_\_\_\_\_\_\_\_\_\_\_\_\_\_\_\_\_
% mdparagraph:P163:end

% output: NEW-U16-F01-O012
% mdparagraph:P164:begin
(12) ðwēan: explain the stem relations: \_\_\_\_\_\_\_\_\_\_\_\_\_\_\_\_\_\_\_\_
% mdparagraph:P164:end

% output: NEW-U16-F01-O013
% mdparagraph:P165:begin
(13) iernan: present indicative, he: \_\_\_\_\_\_\_\_\_\_\_\_\_\_\_\_\_\_\_\_
% mdparagraph:P165:end

% output: NEW-U16-F01-O014
% mdparagraph:P166:begin
(14) iernan: preterite indicative, he: \_\_\_\_\_\_\_\_\_\_\_\_\_\_\_\_\_\_\_\_
% mdparagraph:P166:end

% output: NEW-U16-F01-O015
% mdparagraph:P167:begin
(15) iernan: preterite indicative, they: \_\_\_\_\_\_\_\_\_\_\_\_\_\_\_\_\_\_\_\_
% mdparagraph:P167:end

% output: NEW-U16-F01-O016
% mdparagraph:P168:begin
(16) iernan: explain the stem relations: \_\_\_\_\_\_\_\_\_\_\_\_\_\_\_\_\_\_\_\_
% mdparagraph:P168:end

\end{minipage}\par\addvspace{4pt}

\par\begin{minipage}{\linewidth}\raggedright
% mdparagraph:P169:begin
\textbf{F02}
% mdparagraph:P169:end

% item: NEW-U16-F02
% mdparagraph:P170:begin
Give the selected forms requested below. Then explain the two stated contrasts. These are separate positions even when their spelling coincides.
% mdparagraph:P170:end

% mdparagraph:P171:begin
(1) smierwan
% mdparagraph:P171:end

% mdparagraph:P172:begin
(2) bewęndan
% mdparagraph:P172:end

% mdparagraph:P173:begin
(3) plegian
% mdparagraph:P173:end

% mdparagraph:P174:begin
(4) behēafdian
% mdparagraph:P174:end

% mdparagraph:P175:begin
(5) gefyllan
% mdparagraph:P175:end

% mdparagraph:P176:begin
(6) forgiefan
% mdparagraph:P176:end

% mdparagraph:P177:begin
\textbf{Supplied words.} These forms and meanings remain available for this task.
% mdparagraph:P177:end

% mdtable:T60:begin
\par\begingroup
\fontsize{10.5}{13.2}\selectfont
\setlength{\LTleft}{0pt}\setlength{\LTright}{\fill}
\noindent
\begin{tabular}{@{}>{\raggedright\arraybackslash}p{33.87mm}>{\raggedright\arraybackslash}p{33.87mm}>{\raggedright\arraybackslash}p{50.64mm}@{}}
\toprule
\textbf{Form} & \textbf{Meaning} & \textbf{Role} \\
\midrule
smierwan & anoint & given dictionary form \\
bewęndan & turn & given dictionary form \\
plegian & play & given dictionary form \\
behēafdian & behead & given dictionary form \\
forgiefan & forgive & given dictionary form \\
\bottomrule
\end{tabular}
\par\endgroup

% mdtable:T60:end
% output: NEW-U16-F02-O001
% mdparagraph:P178:begin
(1) smierwan: present indicative, I: \_\_\_\_\_\_\_\_\_\_\_\_\_\_\_\_\_\_\_\_
% mdparagraph:P178:end

% output: NEW-U16-F02-O002
% mdparagraph:P179:begin
(2) smierwan: present indicative, you singular: \_\_\_\_\_\_\_\_\_\_\_\_\_\_\_\_\_\_\_\_
% mdparagraph:P179:end

% output: NEW-U16-F02-O003
% mdparagraph:P180:begin
(3) smierwan: preterite indicative, she: \_\_\_\_\_\_\_\_\_\_\_\_\_\_\_\_\_\_\_\_
% mdparagraph:P180:end

% output: NEW-U16-F02-O004
% mdparagraph:P181:begin
(4) bewęndan: present indicative, I: \_\_\_\_\_\_\_\_\_\_\_\_\_\_\_\_\_\_\_\_
% mdparagraph:P181:end

% output: NEW-U16-F02-O005
% mdparagraph:P182:begin
(5) bewęndan: preterite indicative, he: \_\_\_\_\_\_\_\_\_\_\_\_\_\_\_\_\_\_\_\_
% mdparagraph:P182:end

% output: NEW-U16-F02-O006
% mdparagraph:P183:begin
(6) plegian: preterite indicative, they: \_\_\_\_\_\_\_\_\_\_\_\_\_\_\_\_\_\_\_\_
% mdparagraph:P183:end

% output: NEW-U16-F02-O007
% mdparagraph:P184:begin
(7) behēafdian: preterite indicative, they: \_\_\_\_\_\_\_\_\_\_\_\_\_\_\_\_\_\_\_\_
% mdparagraph:P184:end

% output: NEW-U16-F02-O008
% mdparagraph:P185:begin
(8) gefyllan: uninflected past participle: \_\_\_\_\_\_\_\_\_\_\_\_\_\_\_\_\_\_\_\_
% mdparagraph:P185:end

% output: NEW-U16-F02-O009
% mdparagraph:P186:begin
(9) forgiefan: uninflected past participle: \_\_\_\_\_\_\_\_\_\_\_\_\_\_\_\_\_\_\_\_
% mdparagraph:P186:end

% output: NEW-U16-F02-O010
% mdparagraph:P187:begin
(10) Explain the smierwan stem distribution shown by your answers: \_\_\_\_\_\_\_\_\_\_\_\_\_\_\_\_\_\_\_\_
% mdparagraph:P187:end

% output: NEW-U16-F02-O011
% mdparagraph:P188:begin
(11) Explain what the two bewęndan answers establish about identifying tense from spelling alone: \_\_\_\_\_\_\_\_\_\_\_\_\_\_\_\_\_\_\_\_
% mdparagraph:P188:end

\end{minipage}\par\addvspace{4pt}

\par\begin{minipage}{\linewidth}\raggedright
% mdparagraph:P189:begin
\textbf{F03}
% mdparagraph:P189:end

% item: NEW-U16-F03
% mdparagraph:P190:begin
For each listed verb, give a command to one person and then to more than one. Use the selected command pair taught in this unit.
% mdparagraph:P190:end

% mdparagraph:P191:begin
(1) āceorfan
% mdparagraph:P191:end

% mdparagraph:P192:begin
(2) blissian
% mdparagraph:P192:end

% mdparagraph:P193:begin
(3) smierwan
% mdparagraph:P193:end

% mdparagraph:P194:begin
(4) ārīsan
% mdparagraph:P194:end

% mdparagraph:P195:begin
(5) forlǣtan
% mdparagraph:P195:end

% mdparagraph:P196:begin
(6) geswīcan
% mdparagraph:P196:end

% mdparagraph:P197:begin
(7) ęfthweorfan
% mdparagraph:P197:end

% mdparagraph:P198:begin
\textbf{Supplied words.} These forms and meanings remain available for this task.
% mdparagraph:P198:end

% mdtable:T61:begin
\par\begingroup
\fontsize{10.5}{13.2}\selectfont
\setlength{\LTleft}{0pt}\setlength{\LTright}{\fill}
\noindent
\begin{tabular}{@{}>{\raggedright\arraybackslash}p{34.27mm}>{\raggedright\arraybackslash}p{37.54mm}>{\raggedright\arraybackslash}p{46.57mm}@{}}
\toprule
\textbf{Form} & \textbf{Meaning} & \textbf{Role} \\
\midrule
āceorfan & cut off & given dictionary form \\
blissian & rejoice & given dictionary form \\
smierwan & anoint & given dictionary form \\
ārīsan & arise & given dictionary form \\
forlǣtan & abandon; forsake & given dictionary form \\
geswīcan & cease & given dictionary form \\
ęfthweorfan & return; turn back & given dictionary form \\
\bottomrule
\end{tabular}
\par\endgroup

% mdtable:T61:end
% output: NEW-U16-F03-O001
% mdparagraph:P199:begin
(1) āceorfan: one addressee: \_\_\_\_\_\_\_\_\_\_\_\_\_\_\_\_\_\_\_\_
% mdparagraph:P199:end

% output: NEW-U16-F03-O002
% mdparagraph:P200:begin
(2) āceorfan: more than one addressee: \_\_\_\_\_\_\_\_\_\_\_\_\_\_\_\_\_\_\_\_
% mdparagraph:P200:end

% output: NEW-U16-F03-O003
% mdparagraph:P201:begin
(3) blissian: one addressee: \_\_\_\_\_\_\_\_\_\_\_\_\_\_\_\_\_\_\_\_
% mdparagraph:P201:end

% output: NEW-U16-F03-O004
% mdparagraph:P202:begin
(4) blissian: more than one addressee: \_\_\_\_\_\_\_\_\_\_\_\_\_\_\_\_\_\_\_\_
% mdparagraph:P202:end

% output: NEW-U16-F03-O005
% mdparagraph:P203:begin
(5) smierwan: one addressee: \_\_\_\_\_\_\_\_\_\_\_\_\_\_\_\_\_\_\_\_
% mdparagraph:P203:end

% output: NEW-U16-F03-O006
% mdparagraph:P204:begin
(6) smierwan: more than one addressee: \_\_\_\_\_\_\_\_\_\_\_\_\_\_\_\_\_\_\_\_
% mdparagraph:P204:end

% output: NEW-U16-F03-O007
% mdparagraph:P205:begin
(7) ārīsan: one addressee: \_\_\_\_\_\_\_\_\_\_\_\_\_\_\_\_\_\_\_\_
% mdparagraph:P205:end

% output: NEW-U16-F03-O008
% mdparagraph:P206:begin
(8) ārīsan: more than one addressee: \_\_\_\_\_\_\_\_\_\_\_\_\_\_\_\_\_\_\_\_
% mdparagraph:P206:end

% output: NEW-U16-F03-O009
% mdparagraph:P207:begin
(9) forlǣtan: one addressee: \_\_\_\_\_\_\_\_\_\_\_\_\_\_\_\_\_\_\_\_
% mdparagraph:P207:end

% output: NEW-U16-F03-O010
% mdparagraph:P208:begin
(10) forlǣtan: more than one addressee: \_\_\_\_\_\_\_\_\_\_\_\_\_\_\_\_\_\_\_\_
% mdparagraph:P208:end

% output: NEW-U16-F03-O011
% mdparagraph:P209:begin
(11) geswīcan: one addressee: \_\_\_\_\_\_\_\_\_\_\_\_\_\_\_\_\_\_\_\_
% mdparagraph:P209:end

% output: NEW-U16-F03-O012
% mdparagraph:P210:begin
(12) geswīcan: more than one addressee: \_\_\_\_\_\_\_\_\_\_\_\_\_\_\_\_\_\_\_\_
% mdparagraph:P210:end

% output: NEW-U16-F03-O013
% mdparagraph:P211:begin
(13) ęfthweorfan: one addressee: \_\_\_\_\_\_\_\_\_\_\_\_\_\_\_\_\_\_\_\_
% mdparagraph:P211:end

% output: NEW-U16-F03-O014
% mdparagraph:P212:begin
(14) ęfthweorfan: more than one addressee: \_\_\_\_\_\_\_\_\_\_\_\_\_\_\_\_\_\_\_\_
% mdparagraph:P212:end

\end{minipage}\par\addvspace{4pt}

\par\begin{minipage}{\linewidth}\raggedright
% mdparagraph:P213:begin
\textbf{F04}
% mdparagraph:P213:end

% item: NEW-U16-F04
% mdparagraph:P214:begin
Give the requested present-subjunctive forms and the two infinitive forms. The labels specify the form, not an English tense-conversion rule.
% mdparagraph:P214:end

% mdparagraph:P215:begin
(1) findan
% mdparagraph:P215:end

% mdparagraph:P216:begin
(2) fēdan
% mdparagraph:P216:end

% mdparagraph:P217:begin
(3) iernan
% mdparagraph:P217:end

% mdparagraph:P218:begin
\textbf{Supplied words.} These forms and meanings remain available for this task.
% mdparagraph:P218:end

% mdtable:T62:begin
\par\begingroup
\fontsize{10.5}{13.2}\selectfont
\setlength{\LTleft}{0pt}\setlength{\LTright}{\fill}
\noindent
\begin{tabular}{@{}>{\raggedright\arraybackslash}p{33.87mm}>{\raggedright\arraybackslash}p{33.87mm}>{\raggedright\arraybackslash}p{50.64mm}@{}}
\toprule
\textbf{Form} & \textbf{Meaning} & \textbf{Role} \\
\midrule
findan & find & given dictionary form \\
fēdan & feed; nourish & given dictionary form \\
\bottomrule
\end{tabular}
\par\endgroup

% mdtable:T62:end
% output: NEW-U16-F04-O001
% mdparagraph:P219:begin
(1) findan: present subjunctive, third singular: \_\_\_\_\_\_\_\_\_\_\_\_\_\_\_\_\_\_\_\_
% mdparagraph:P219:end

% output: NEW-U16-F04-O002
% mdparagraph:P220:begin
(2) findan: present subjunctive, third plural: \_\_\_\_\_\_\_\_\_\_\_\_\_\_\_\_\_\_\_\_
% mdparagraph:P220:end

% output: NEW-U16-F04-O003
% mdparagraph:P221:begin
(3) fēdan: present subjunctive, third singular: \_\_\_\_\_\_\_\_\_\_\_\_\_\_\_\_\_\_\_\_
% mdparagraph:P221:end

% output: NEW-U16-F04-O004
% mdparagraph:P222:begin
(4) fēdan: present subjunctive, third plural: \_\_\_\_\_\_\_\_\_\_\_\_\_\_\_\_\_\_\_\_
% mdparagraph:P222:end

% output: NEW-U16-F04-O005
% mdparagraph:P223:begin
(5) iernan: plain infinitive: \_\_\_\_\_\_\_\_\_\_\_\_\_\_\_\_\_\_\_\_
% mdparagraph:P223:end

% output: NEW-U16-F04-O006
% mdparagraph:P224:begin
(6) iernan: inflected infinitive after tō (give the infinitive only): \_\_\_\_\_\_\_\_\_\_\_\_\_\_\_\_\_\_\_\_
% mdparagraph:P224:end

\end{minipage}\par\addvspace{4pt}

\par\begin{minipage}{\linewidth}\raggedright
% mdparagraph:P225:begin
\textbf{F05}
% mdparagraph:P225:end

% item: NEW-U16-F05
% mdparagraph:P226:begin
Give the four requested strong forms of nearu. Then translate the three phrases using bēgen and sunu, without adding an article. Finally explain the stem alternation and the number of both.
% mdparagraph:P226:end

% mdparagraph:P227:begin
(1) nearu
% mdparagraph:P227:end

% mdparagraph:P228:begin
(2) bēgen
% mdparagraph:P228:end

% mdparagraph:P229:begin
(3) sunu
% mdparagraph:P229:end

% output: NEW-U16-F05-O001
% mdparagraph:P230:begin
(1) nearu: masculine nominative singular: \_\_\_\_\_\_\_\_\_\_\_\_\_\_\_\_\_\_\_\_
% mdparagraph:P230:end

% output: NEW-U16-F05-O002
% mdparagraph:P231:begin
(2) nearu: masculine accusative singular: \_\_\_\_\_\_\_\_\_\_\_\_\_\_\_\_\_\_\_\_
% mdparagraph:P231:end

% output: NEW-U16-F05-O003
% mdparagraph:P232:begin
(3) nearu: masculine genitive singular: \_\_\_\_\_\_\_\_\_\_\_\_\_\_\_\_\_\_\_\_
% mdparagraph:P232:end

% output: NEW-U16-F05-O004
% mdparagraph:P233:begin
(4) nearu: masculine dative plural: \_\_\_\_\_\_\_\_\_\_\_\_\_\_\_\_\_\_\_\_
% mdparagraph:P233:end

% output: NEW-U16-F05-O005
% mdparagraph:P234:begin
(5) both sons, accusative: \_\_\_\_\_\_\_\_\_\_\_\_\_\_\_\_\_\_\_\_
% mdparagraph:P234:end

% output: NEW-U16-F05-O006
% mdparagraph:P235:begin
(6) of both sons, genitive: \_\_\_\_\_\_\_\_\_\_\_\_\_\_\_\_\_\_\_\_
% mdparagraph:P235:end

% output: NEW-U16-F05-O007
% mdparagraph:P236:begin
(7) to both sons, dative: \_\_\_\_\_\_\_\_\_\_\_\_\_\_\_\_\_\_\_\_
% mdparagraph:P236:end

% output: NEW-U16-F05-O008
% mdparagraph:P237:begin
(8) Explain the nearu stem alternation shown by your answers: \_\_\_\_\_\_\_\_\_\_\_\_\_\_\_\_\_\_\_\_
% mdparagraph:P237:end

% output: NEW-U16-F05-O009
% mdparagraph:P238:begin
(9) Explain the grammatical number of both in the phrases: \_\_\_\_\_\_\_\_\_\_\_\_\_\_\_\_\_\_\_\_
% mdparagraph:P238:end

\end{minipage}\par\addvspace{4pt}

\par\begin{minipage}{\linewidth}\raggedright
% mdparagraph:P239:begin
\textbf{A01}
% mdparagraph:P239:end

% item: NEW-U16-A01
% mdparagraph:P240:begin
Give the two defensible readings of the final hē. Then rewrite the second sentence twice, replacing only hē with an explicit noun phrase so that each reading is unambiguous. The first sentence identifies the king's son.
% mdparagraph:P240:end

% mdparagraph:P241:begin
(1) Se cyning geseah his sunu. Hē ēode.
% mdparagraph:P241:end

% output: NEW-U16-A01-O001
% mdparagraph:P242:begin
(1) First reading: \_\_\_\_\_\_\_\_\_\_\_\_\_\_\_\_\_\_\_\_
% mdparagraph:P242:end

% output: NEW-U16-A01-O002
% mdparagraph:P243:begin
(2) Second reading: \_\_\_\_\_\_\_\_\_\_\_\_\_\_\_\_\_\_\_\_
% mdparagraph:P243:end

% output: NEW-U16-A01-O003
% mdparagraph:P244:begin
(3) First explicit version: \_\_\_\_\_\_\_\_\_\_\_\_\_\_\_\_\_\_\_\_
% mdparagraph:P244:end

% output: NEW-U16-A01-O004
% mdparagraph:P245:begin
(4) Second explicit version: \_\_\_\_\_\_\_\_\_\_\_\_\_\_\_\_\_\_\_\_
% mdparagraph:P245:end

\end{minipage}\par\addvspace{4pt}

\par\begin{minipage}{\linewidth}\raggedright
% mdparagraph:P246:begin
\textbf{A02}
% mdparagraph:P246:end

% item: NEW-U16-A02
% mdparagraph:P247:begin
Translate both sentences. For each, identify the expression denoting the person who went, its case and its relation to the verb of going. Then compare the two forms of go: which is nonfinite, which is finite, and does the spelling of the finite form distinguish its mood here?
% mdparagraph:P247:end

% mdparagraph:P248:begin
(1) Se cyning geseah ðone sunu gān.
% mdparagraph:P248:end

% mdparagraph:P249:begin
(2) Se cyning wiste ðæt se sunu ēode.
% mdparagraph:P249:end

% output: NEW-U16-A02-O001
% mdparagraph:P250:begin
(1) Sentence 1 translation: \_\_\_\_\_\_\_\_\_\_\_\_\_\_\_\_\_\_\_\_
% mdparagraph:P250:end

% output: NEW-U16-A02-O002
% mdparagraph:P251:begin
(2) Sentence 2 translation: \_\_\_\_\_\_\_\_\_\_\_\_\_\_\_\_\_\_\_\_
% mdparagraph:P251:end

% output: NEW-U16-A02-O003
% mdparagraph:P252:begin
(3) Sentence 1 actor: \_\_\_\_\_\_\_\_\_\_\_\_\_\_\_\_\_\_\_\_
% mdparagraph:P252:end

% output: NEW-U16-A02-O004
% mdparagraph:P253:begin
(4) Sentence 2 actor: \_\_\_\_\_\_\_\_\_\_\_\_\_\_\_\_\_\_\_\_
% mdparagraph:P253:end

% output: NEW-U16-A02-O005
% mdparagraph:P254:begin
(5) Verb-form comparison: \_\_\_\_\_\_\_\_\_\_\_\_\_\_\_\_\_\_\_\_
% mdparagraph:P254:end

\end{minipage}\par\addvspace{4pt}

\par\begin{minipage}{\linewidth}\raggedright
% mdparagraph:P255:begin
\textbf{A03}
% mdparagraph:P255:end

% item: NEW-U16-A03
% mdparagraph:P256:begin
Translate each report. For each, identify and parse every finite verb, identify any participle or infinitive (or state that there is none), and explain what the construction and context show about the time of the reported situation. In sentence 1 the king is there at the time of saying; in sentence 2 the ship-making is complete; in sentence 3 the singing is intended but has not yet occurred. Your explanation must connect these contexts with the actual verb groups and explain why automatic English backshift is insufficient; merely repeating the time hint is not an answer.
% mdparagraph:P256:end

% mdparagraph:P257:begin
(1) Hē sægde ðæt se cyning ðǣr wǣre.
% mdparagraph:P257:end

% mdparagraph:P258:begin
(2) Hē sægde ðæt se cyning ðæt scip geworht hæfde.
% mdparagraph:P258:end

% mdparagraph:P259:begin
(3) Hē sægde ðæt se cyning singan wolde.
% mdparagraph:P259:end

% output: NEW-U16-A03-O001
% mdparagraph:P260:begin
(1) Sentence 1 translation: \_\_\_\_\_\_\_\_\_\_\_\_\_\_\_\_\_\_\_\_
% mdparagraph:P260:end

% output: NEW-U16-A03-O002
% mdparagraph:P261:begin
(2) Sentence 1 finite forms: \_\_\_\_\_\_\_\_\_\_\_\_\_\_\_\_\_\_\_\_
% mdparagraph:P261:end

% output: NEW-U16-A03-O003
% mdparagraph:P262:begin
(3) Sentence 1 nonfinite forms: \_\_\_\_\_\_\_\_\_\_\_\_\_\_\_\_\_\_\_\_
% mdparagraph:P262:end

% output: NEW-U16-A03-O004
% mdparagraph:P263:begin
(4) Sentence 1 construction and time: \_\_\_\_\_\_\_\_\_\_\_\_\_\_\_\_\_\_\_\_
% mdparagraph:P263:end

% output: NEW-U16-A03-O005
% mdparagraph:P264:begin
(5) Sentence 2 translation: \_\_\_\_\_\_\_\_\_\_\_\_\_\_\_\_\_\_\_\_
% mdparagraph:P264:end

% output: NEW-U16-A03-O006
% mdparagraph:P265:begin
(6) Sentence 2 finite forms: \_\_\_\_\_\_\_\_\_\_\_\_\_\_\_\_\_\_\_\_
% mdparagraph:P265:end

% output: NEW-U16-A03-O007
% mdparagraph:P266:begin
(7) Sentence 2 nonfinite form: \_\_\_\_\_\_\_\_\_\_\_\_\_\_\_\_\_\_\_\_
% mdparagraph:P266:end

% output: NEW-U16-A03-O008
% mdparagraph:P267:begin
(8) Sentence 2 construction and time: \_\_\_\_\_\_\_\_\_\_\_\_\_\_\_\_\_\_\_\_
% mdparagraph:P267:end

% output: NEW-U16-A03-O009
% mdparagraph:P268:begin
(9) Sentence 3 translation: \_\_\_\_\_\_\_\_\_\_\_\_\_\_\_\_\_\_\_\_
% mdparagraph:P268:end

% output: NEW-U16-A03-O010
% mdparagraph:P269:begin
(10) Sentence 3 finite forms: \_\_\_\_\_\_\_\_\_\_\_\_\_\_\_\_\_\_\_\_
% mdparagraph:P269:end

% output: NEW-U16-A03-O011
% mdparagraph:P270:begin
(11) Sentence 3 nonfinite form: \_\_\_\_\_\_\_\_\_\_\_\_\_\_\_\_\_\_\_\_
% mdparagraph:P270:end

% output: NEW-U16-A03-O012
% mdparagraph:P271:begin
(12) Sentence 3 construction and time: \_\_\_\_\_\_\_\_\_\_\_\_\_\_\_\_\_\_\_\_
% mdparagraph:P271:end

\end{minipage}\par\addvspace{4pt}

\reserve{7\baselineskip}
% mdparagraph:P272:begin
\section*{5. New compositions}
% mdparagraph:P272:end

% mdparagraph:P273:begin
Translate every clause and identify the requested relations. Keep a grammatical alternative distinct from compliance with an explicitly requested verb or construction.
% mdparagraph:P273:end

\reserve{5\baselineskip}
\par\begin{minipage}{\linewidth}\raggedright
% mdparagraph:P274:begin
\textbf{C01}
% mdparagraph:P274:end

% item: NEW-U16-C01
% mdparagraph:P275:begin
Translate the complete exchange. Use sęcgan for the reporting verb and sęllan for give: The king said to his son, 'Give me the book.' The son gave him the book. Then identify the referents of the first-person and third-person recipient pronouns in your translation.
% mdparagraph:P275:end

% output: NEW-U16-C01-O001
% mdparagraph:P276:begin
(1) Speech and command: \_\_\_\_\_\_\_\_\_\_\_\_\_\_\_\_\_\_\_\_
% mdparagraph:P276:end

% output: NEW-U16-C01-O002
% mdparagraph:P277:begin
(2) Executed action: \_\_\_\_\_\_\_\_\_\_\_\_\_\_\_\_\_\_\_\_
% mdparagraph:P277:end

% output: NEW-U16-C01-O003
% mdparagraph:P278:begin
(3) Recipient reference: \_\_\_\_\_\_\_\_\_\_\_\_\_\_\_\_\_\_\_\_
% mdparagraph:P278:end

\end{minipage}\par\addvspace{4pt}

\par\begin{minipage}{\linewidth}\raggedright
% mdparagraph:P279:begin
\textbf{C02}
% mdparagraph:P279:end

% item: NEW-U16-C02
% mdparagraph:P280:begin
Translate these three statements separately. Use cunnan in the first two and witan in the third; then explain the complement selected in each: I know that man. I know how to sing. I know that he is here.
% mdparagraph:P280:end

% output: NEW-U16-C02-O001
% mdparagraph:P281:begin
(1) Statement 1: \_\_\_\_\_\_\_\_\_\_\_\_\_\_\_\_\_\_\_\_
% mdparagraph:P281:end

% output: NEW-U16-C02-O002
% mdparagraph:P282:begin
(2) Statement 2: \_\_\_\_\_\_\_\_\_\_\_\_\_\_\_\_\_\_\_\_
% mdparagraph:P282:end

% output: NEW-U16-C02-O003
% mdparagraph:P283:begin
(3) Statement 3: \_\_\_\_\_\_\_\_\_\_\_\_\_\_\_\_\_\_\_\_
% mdparagraph:P283:end

% output: NEW-U16-C02-O004
% mdparagraph:P284:begin
(4) Complement comparison: \_\_\_\_\_\_\_\_\_\_\_\_\_\_\_\_\_\_\_\_
% mdparagraph:P284:end

\end{minipage}\par\addvspace{4pt}

\par\begin{minipage}{\linewidth}\raggedright
% mdparagraph:P285:begin
\textbf{C03}
% mdparagraph:P285:end

% item: NEW-U16-C03
% mdparagraph:P286:begin
Translate both clauses, using plegian for the first two expressed actions and dōn for the final repeated action: The men played well; the king played better than they did. Explain what the last verb replaces and why it is plural.
% mdparagraph:P286:end

% output: NEW-U16-C03-O001
% mdparagraph:P287:begin
(1) Whole comparison: \_\_\_\_\_\_\_\_\_\_\_\_\_\_\_\_\_\_\_\_
% mdparagraph:P287:end

% output: NEW-U16-C03-O002
% mdparagraph:P288:begin
(2) Recovered action and agreement: \_\_\_\_\_\_\_\_\_\_\_\_\_\_\_\_\_\_\_\_
% mdparagraph:P288:end

\end{minipage}\par\addvspace{4pt}

\reserve{7\baselineskip}
% mdparagraph:P289:begin
\section*{6. Cook: Group XII, XIII, Exercises 32–36}
% mdparagraph:P289:end

% teaching: NEW-U16-COOK-SUPPORT
% mdparagraph:P290:begin
Translate the complete original prompts. Parenthesized words, omissions, number instructions and order directions belong to Cook’s exercises. The modern helps below are additional and remain available; a supplied expression does not require independent recall. Cook’s paragraph numbers refer to his grammar, but the material needed here is stated in this unit.\textsuperscript{6}
% mdparagraph:P290:end

% mdparagraph:P291:begin
\textbf{Source correction for 34.9.} Cook prints “reward (āgyltan, 113)”. The verb meaning “repay; reward” is strong \textbf{āgyldan}, belonging with §104; \textbf{āgylt} is a valid contracted third singular. The complete reward clause is supplied with the item. No additional full conjugation is required. The original prompt is retained.\textsuperscript{7}
% mdparagraph:P291:end

% mdparagraph:P292:begin
\textbf{Names in 36.9–10.} The supplied \textbf{Johannes} is genitive and \textbf{Johannem} accusative. Their endings are directly attested with initial I; initial J is used here to follow Cook’s headword. These name forms remain supplied.\textsuperscript{8}
% mdparagraph:P292:end

\reserve{7\baselineskip}
% mdparagraph:P293:begin
\subsection*{Exercise 32}
% mdparagraph:P293:end

\reserve{5\baselineskip}
\par\begin{minipage}{\linewidth}\raggedright
% mdparagraph:P294:begin
\textbf{Cook XII, XIII.32.1}
% mdparagraph:P294:end

% item: EX-XII-XIII-32-01
% mdparagraph:P295:begin
I bid you all farewell (grētan).
% mdparagraph:P295:end

% output: EX-XII-XIII-32-01-O001
\end{minipage}\par\addvspace{4pt}

\par\begin{minipage}{\linewidth}\raggedright
% mdparagraph:P296:begin
\textbf{Cook XII, XIII.32.2}
% mdparagraph:P296:end

% item: EX-XII-XIII-32-02
% mdparagraph:P297:begin
When they had rowed about (swilce) twenty furlongs (furlang, 47), or thirty, then saw they Jesus walking (gān) on (uppan, 166) the sea (verb last).
% mdparagraph:P297:end

% mdparagraph:P298:begin
\textbf{Supplied help:} swilce twēntig furlanga oððe ðrītig; uppan ðǣm sǣ
% mdparagraph:P298:end

% mdparagraph:P299:begin
About twenty furlongs or thirty; upon the sea. The supplied numerical phrase takes genitive furlanga; the prepositional phrase has the dative.
% mdparagraph:P299:end

\end{minipage}\par\addvspace{4pt}

% teaching: NEW-U16-SUP-SWILCE
% teaching: NEW-U16-VOC-SWILCE
% teaching: NEW-U16-SUP-TWENTIG
% teaching: NEW-U16-VOC-TWENTIG
% teaching: NEW-U16-SUP-THRITIG
% teaching: NEW-U16-VOC-THRITIG
% teaching: NEW-U16-SUP-UPPAN
% teaching: NEW-U16-VOC-UPPAN
\reserve{11\baselineskip}
% mdparagraph:P300:begin
\textbf{Supplied words.} These forms and meanings remain available for this task.
% mdparagraph:P300:end

% mdtable:T63:begin
\par\begingroup
\fontsize{10.5}{13.2}\selectfont
\setlength{\LTleft}{0pt}\setlength{\LTright}{\fill}
\noindent
\begin{tabular}{@{}>{\raggedright\arraybackslash}p{31.95mm}>{\raggedright\arraybackslash}p{51.6mm}>{\raggedright\arraybackslash}p{34.83mm}@{}}
\toprule
\textbf{Form} & \textbf{Meaning} & \textbf{Role} \\
\midrule
swilce & about; approximately & invariable \\
twēntig & twenty & as used here \\
ðrītig & thirty & as used here \\
uppan & upon; on top of & invariable \\
\bottomrule
\end{tabular}
\par\endgroup

% mdtable:T63:end
% output: EX-XII-XIII-32-02-O001
\reserve{5\baselineskip}
\par\begin{minipage}{\linewidth}\raggedright
% mdparagraph:P301:begin
\textbf{Cook XII, XIII.32.3}
% mdparagraph:P301:end

% item: EX-XII-XIII-32-03
% mdparagraph:P302:begin
He arose and commanded (bebēodan) the wind (164. b), and said (cweðan) to the sea, Subside.
% mdparagraph:P302:end

% output: EX-XII-XIII-32-03-O001
\end{minipage}\par\addvspace{4pt}

\par\begin{minipage}{\linewidth}\raggedright
% mdparagraph:P303:begin
\textbf{Cook XII, XIII.32.4}
% mdparagraph:P303:end

% item: EX-XII-XIII-32-04
% mdparagraph:P304:begin
The wind ceased (geswīcan), and there was (weorðan) (om. a) great calm.
% mdparagraph:P304:end

% output: EX-XII-XIII-32-04-O001
\end{minipage}\par\addvspace{4pt}

\par\begin{minipage}{\linewidth}\raggedright
% mdparagraph:P305:begin
\textbf{Cook XII, XIII.32.5}
% mdparagraph:P305:end

% item: EX-XII-XIII-32-05
% mdparagraph:P306:begin
Thou hearest his voice, and thou knowest not whence he cometh, nor whither he goeth.
% mdparagraph:P306:end

% output: EX-XII-XIII-32-05-O001
\end{minipage}\par\addvspace{4pt}

\par\begin{minipage}{\linewidth}\raggedright
% mdparagraph:P307:begin
\textbf{Cook XII, XIII.32.6}
% mdparagraph:P307:end

% item: EX-XII-XIII-32-06
% mdparagraph:P308:begin
He saw two brethren casting (sęndan) their net (nętt, 47) into (on, 172) the sea; they were fishers.
% mdparagraph:P308:end

% output: EX-XII-XIII-32-06-O001
\end{minipage}\par\addvspace{4pt}

\par\begin{minipage}{\linewidth}\raggedright
% mdparagraph:P309:begin
\textbf{Cook XII, XIII.32.7}
% mdparagraph:P309:end

% item: EX-XII-XIII-32-07
% mdparagraph:P310:begin
And he said (sęcgean) to them (164. b), Come after (æfter, 166) me, and I (om. will) make (dōn) that ye (dual) are fishers of men (153. d); and they immediately left their nets.
% mdparagraph:P310:end

% output: EX-XII-XIII-32-07-O001
\end{minipage}\par\addvspace{4pt}

\par\begin{minipage}{\linewidth}\raggedright
% mdparagraph:P311:begin
\textbf{Cook XII, XIII.32.8}
% mdparagraph:P311:end

% item: EX-XII-XIII-32-08
% mdparagraph:P312:begin
If thy hand or thy foot offend (ǣswician, 118) thee, cut it (āceorfan, 104) off (of), and cast (om. it) from thee; (om. it) is better (om. for) thee (165) that thou go halt (healt, 58) into life, than (om. that) thou have two hands and two feet, and be cast into everlasting fire.
% mdparagraph:P312:end

% mdparagraph:P313:begin
\textbf{Supplied help:} bętere ðē is ðæt ðū ... gā, ðonne ðū ... hæbbe and ... sīe worpen
% mdparagraph:P313:end

% mdparagraph:P314:begin
It is better for you to go ... than to have ... and be thrown ...
% mdparagraph:P314:end

% mdparagraph:P315:begin
The experiencer is dative ðē; bętere is the familiar neuter comparative predicate. Keep the requested ðæt in the first complement and omit it after ðonne. Earlier hit resumes whichever of the differently gendered body parts is at issue.
% mdparagraph:P315:end

% output: EX-XII-XIII-32-08-O001
\end{minipage}\par\addvspace{4pt}

\par\begin{minipage}{\linewidth}\raggedright
% mdparagraph:P316:begin
\textbf{Cook XII, XIII.32.9}
% mdparagraph:P316:end

% item: EX-XII-XIII-32-09
% mdparagraph:P317:begin
If any (hwilc) man hath (om. a) hundred sheep (154. c), and loseth (forlēosan) one of (of, 166) them [those], (om. doth) he not leave the nine and ninety in the mountains, and goeth and seeketh (sēcean) the one that (87. d) was perishing (forwearð)?
% mdparagraph:P317:end

% mdparagraph:P318:begin
\textbf{Supplied help:} hund scēapa; nigon and hundnigontig; ān of ðǣm
% mdparagraph:P318:end

% mdparagraph:P319:begin
A hundred sheep; ninety-nine; one of those. These numerical and partitive expressions remain supplied.
% mdparagraph:P319:end

\end{minipage}\par\addvspace{4pt}

% teaching: NEW-U16-SUP-HUND
% teaching: NEW-U16-VOC-HUND
% teaching: NEW-U16-SUP-HUNDNIGONTIG
% teaching: NEW-U16-VOC-HUNDNIGONTIG
% teaching: NEW-U16-VOC-OF
\reserve{10\baselineskip}
% mdparagraph:P320:begin
\textbf{Supplied words.} These forms and meanings remain available for this task.
% mdparagraph:P320:end

% mdtable:T64:begin
\par\begingroup
\fontsize{10.5}{13.2}\selectfont
\setlength{\LTleft}{0pt}\setlength{\LTright}{\fill}
\noindent
\begin{tabular}{@{}>{\raggedright\arraybackslash}p{22.44mm}>{\raggedright\arraybackslash}p{77.23mm}>{\raggedright\arraybackslash}p{18.7mm}@{}}
\toprule
\textbf{Form} & \textbf{Meaning} & \textbf{Role} \\
\midrule
hund & hundred, in the supplied hund scēapa phrase & as used here \\
hundnigontig & ninety & as used here \\
of & of; from a portion of a whole, in the supplied phrases ān of ðǣm and of ēowrum ęle & as used here \\
\bottomrule
\end{tabular}
\par\endgroup

% mdtable:T64:end
% output: EX-XII-XIII-32-09-O001
\reserve{5\baselineskip}
\par\begin{minipage}{\linewidth}\raggedright
% mdparagraph:P321:begin
\textbf{Cook XII, XIII.32.10}
% mdparagraph:P321:end

% item: EX-XII-XIII-32-10
% mdparagraph:P322:begin
And if it happeneth that he finds it, truly I say (sęcgean) to you that he rejoiceth more (swīðor) over (ofer, 172) that one than over the nine and ninety.
% mdparagraph:P322:end

% mdparagraph:P323:begin
\textbf{Supplied help:} gif hit gelimpð ðæt hē hit finde
% mdparagraph:P323:end

% mdparagraph:P324:begin
if it happens that he finds it
% mdparagraph:P324:end

% mdparagraph:P325:begin
The first hit is impersonal; the second hit is the sheep. Gelimpð is supplied; finde is the newly taught present subjunctive. No full gelimpan person paradigm is assigned.
% mdparagraph:P325:end

% mdparagraph:P326:begin
\textbf{Supplied help:} hund scēapa; nigon and hundnigontig; ān of ðǣm
% mdparagraph:P326:end

% mdparagraph:P327:begin
A hundred sheep; ninety-nine; one of those. These numerical and partitive expressions remain supplied.
% mdparagraph:P327:end

% mdparagraph:P328:begin
\textbf{Supplied words.} These forms and meanings remain available for this task.
% mdparagraph:P328:end

% mdtable:T65:begin
\par\begingroup
\fontsize{10.5}{13.2}\selectfont
\setlength{\LTleft}{0pt}\setlength{\LTright}{\fill}
\noindent
\begin{tabular}{@{}>{\raggedright\arraybackslash}p{24mm}>{\raggedright\arraybackslash}p{58mm}>{\raggedright\arraybackslash}p{36.38mm}@{}}
\toprule
\textbf{Form} & \textbf{Meaning} & \textbf{Role} \\
\midrule
gelimpð & happen; come about & third person, present, indicative, singular \\
hund & hundred, in the supplied hund scēapa phrase & as used here \\
hundnigontig & ninety & as used here \\
of & of; from a portion of a whole, in the supplied phrases ān of ðǣm and of ēowrum ęle & as used here \\
\bottomrule
\end{tabular}
\par\endgroup

% mdtable:T65:end
% output: EX-XII-XIII-32-10-O001
\end{minipage}\par\addvspace{4pt}

\par\begin{minipage}{\linewidth}\raggedright
% mdparagraph:P329:begin
\textbf{Cook XII, XIII.32.11}
% mdparagraph:P329:end

% item: EX-XII-XIII-32-11
% mdparagraph:P330:begin
When he cometh home (hām), he calleth (cleopian) together his friends and his neighbors (nēahgebūr, 43), and saith (cweðan), Rejoice with me, because I (om. have) found my lost sheep.
% mdparagraph:P330:end

% output: EX-XII-XIII-32-11-O001
\end{minipage}\par\addvspace{4pt}

\reserve{7\baselineskip}
% mdparagraph:P331:begin
\subsection*{Exercise 33}
% mdparagraph:P331:end

\reserve{5\baselineskip}
\par\begin{minipage}{\linewidth}\raggedright
% mdparagraph:P332:begin
\textbf{Cook XII, XIII.33.1}
% mdparagraph:P332:end

% item: EX-XII-XIII-33-01
% mdparagraph:P333:begin
And, going (gangan) into (intō, 166) the house, they found the child with Mary (Maria, 53) his mother.
% mdparagraph:P333:end

% mdparagraph:P334:begin
\textbf{Supplied help:} intō ðǣm hūse gangende
% mdparagraph:P334:end

% mdparagraph:P335:begin
going into the house
% mdparagraph:P335:end

% mdparagraph:P336:begin
Gangende accompanies the plural actors hīe. The participial phrase is supplied as a whole; it does not assign a new complete gangan finite system.
% mdparagraph:P336:end

% mdparagraph:P337:begin
\textbf{Supplied help:} mid Marian his mēder
% mdparagraph:P337:end

% mdparagraph:P338:begin
With Mary his mother. Marian is the supplied weak-feminine form after mid here; the earlier genitive use of the name had a different function.
% mdparagraph:P338:end

\end{minipage}\par\addvspace{4pt}

% teaching: NEW-U16-SUP-MARIA-NAME
% teaching: NEW-U16-VOC-MARIA-NAME
\reserve{9\baselineskip}
% mdparagraph:P339:begin
\textbf{Supplied words.} These forms and meanings remain available for this task.
% mdparagraph:P339:end

% mdtable:T66:begin
\par\begingroup
\fontsize{10.5}{13.2}\selectfont
\setlength{\LTleft}{0pt}\setlength{\LTright}{\fill}
\noindent
\begin{tabular}{@{}>{\raggedright\arraybackslash}p{17.92mm}>{\raggedright\arraybackslash}p{51.71mm}>{\raggedright\arraybackslash}p{48.75mm}@{}}
\toprule
\textbf{Form} & \textbf{Meaning} & \textbf{Role} \\
\midrule
gangende & go; enter in the supplied circumstantial phrase & present, participle, nominative, plural, strong \\
Marian & Mary & dative, singular, feminine \\
\bottomrule
\end{tabular}
\par\endgroup

% mdtable:T66:end
% output: EX-XII-XIII-33-01-O001
\reserve{5\baselineskip}
\par\begin{minipage}{\linewidth}\raggedright
% mdparagraph:P340:begin
\textbf{Cook XII, XIII.33.2}
% mdparagraph:P340:end

% item: EX-XII-XIII-33-02
% mdparagraph:P341:begin
Arise, and take (niman) the child and his mother by night, and flee to (tō, 166) (om. the) land of the Egyptians, and be there until (that) I tell thee (verb last).
% mdparagraph:P341:end

% mdparagraph:P342:begin
\textbf{Supplied help:} Egypta land /\allowbreak{} tō Egypta lande /\allowbreak{} of Egyptum
% mdparagraph:P342:end

% mdparagraph:P343:begin
Land of the Egyptians; to their land; from the Egyptians. These foreign-name forms and their place-reference uses remain supplied.
% mdparagraph:P343:end

\end{minipage}\par\addvspace{4pt}

% teaching: NEW-U16-SUP-EGYPTA
% teaching: NEW-U16-VOC-EGYPTA
% teaching: NEW-U16-VOC-ON
\reserve{10\baselineskip}
% mdparagraph:P344:begin
\textbf{Supplied words.} These forms and meanings remain available for this task.
% mdparagraph:P344:end

% mdtable:T67:begin
\par\begingroup
\fontsize{10.5}{13.2}\selectfont
\setlength{\LTleft}{0pt}\setlength{\LTright}{\fill}
\noindent
\begin{tabular}{@{}>{\raggedright\arraybackslash}p{19.8mm}>{\raggedright\arraybackslash}p{64.73mm}>{\raggedright\arraybackslash}p{33.85mm}@{}}
\toprule
\textbf{Form} & \textbf{Meaning} & \textbf{Role} \\
\midrule
Egypta & Egyptians; country-name phrases Egypta land and of Egyptum & genitive, plural, masculine \\
Egyptum & Egyptians; country-name phrases Egypta land and of Egyptum & dative, plural, masculine \\
on & during, with dative time & as used here \\
\bottomrule
\end{tabular}
\par\endgroup

% mdtable:T67:end
% output: EX-XII-XIII-33-02-O001
\reserve{5\baselineskip}
\par\begin{minipage}{\linewidth}\raggedright
% mdparagraph:P345:begin
\textbf{Cook XII, XIII.33.3}
% mdparagraph:P345:end

% item: EX-XII-XIII-33-03
% mdparagraph:P346:begin
He arose then, and took the child and his mother by night, and journeyed into (on, 172) the land of the Egyptians; that it (om. might) be fulfilled (gefyllan) which (87. a) was spoken (cweðan) of (fram, 166) (om. the) Lord by (ðurh, 172) the prophet, From (of, 166) (om. the) Egyptians (om. have) I called (cleopian) my son.
% mdparagraph:P346:end

% mdparagraph:P347:begin
\textbf{Supplied help:} Egypta land /\allowbreak{} tō Egypta lande /\allowbreak{} of Egyptum
% mdparagraph:P347:end

% mdparagraph:P348:begin
Land of the Egyptians; to their land; from the Egyptians. These foreign-name forms and their place-reference uses remain supplied.
% mdparagraph:P348:end

% mdparagraph:P349:begin
\textbf{Supplied words.} These forms and meanings remain available for this task.
% mdparagraph:P349:end

% mdtable:T68:begin
\par\begingroup
\fontsize{10.5}{13.2}\selectfont
\setlength{\LTleft}{0pt}\setlength{\LTright}{\fill}
\noindent
\begin{tabular}{@{}>{\raggedright\arraybackslash}p{19.8mm}>{\raggedright\arraybackslash}p{64.73mm}>{\raggedright\arraybackslash}p{33.85mm}@{}}
\toprule
\textbf{Form} & \textbf{Meaning} & \textbf{Role} \\
\midrule
Egypta & Egyptians; country-name phrases Egypta land and of Egyptum & genitive, plural, masculine \\
Egyptum & Egyptians; country-name phrases Egypta land and of Egyptum & dative, plural, masculine \\
on & during, with dative time & as used here \\
\bottomrule
\end{tabular}
\par\endgroup

% mdtable:T68:end
% output: EX-XII-XIII-33-03-O001
\end{minipage}\par\addvspace{4pt}

\par\begin{minipage}{\linewidth}\raggedright
% mdparagraph:P350:begin
\textbf{Cook XII, XIII.33.4}
% mdparagraph:P350:end

% item: EX-XII-XIII-33-04
% mdparagraph:P351:begin
What can they ask for [request; 156. b]?
% mdparagraph:P351:end

% output: EX-XII-XIII-33-04-O001
\end{minipage}\par\addvspace{4pt}

\par\begin{minipage}{\linewidth}\raggedright
% mdparagraph:P352:begin
\textbf{Cook XII, XIII.33.5}
% mdparagraph:P352:end

% item: EX-XII-XIII-33-05
% mdparagraph:P353:begin
The sun rules over the day (156. f).
% mdparagraph:P353:end

% output: EX-XII-XIII-33-05-O001
\end{minipage}\par\addvspace{4pt}

\par\begin{minipage}{\linewidth}\raggedright
% mdparagraph:P354:begin
\textbf{Cook XII, XIII.33.6}
% mdparagraph:P354:end

% item: EX-XII-XIII-33-06
% mdparagraph:P355:begin
They expected nothing but [except] death itself (156. g).
% mdparagraph:P355:end

% mdparagraph:P356:begin
\textbf{Supplied help:} nāhtes būtan dēaðes selfes
% mdparagraph:P356:end

% mdparagraph:P357:begin
nothing except death itself: the whole exception phrase is supplied with the genitive governed by the expectation construction.
% mdparagraph:P357:end

\end{minipage}\par\addvspace{4pt}

% teaching: NEW-U16-SUP-NAHT
% teaching: NEW-U16-VOC-NAHT
\reserve{9\baselineskip}
% mdparagraph:P358:begin
\textbf{Supplied words.} These forms and meanings remain available for this task.
% mdparagraph:P358:end

% mdtable:T69:begin
\par\begingroup
\fontsize{10.5}{13.2}\selectfont
\setlength{\LTleft}{0pt}\setlength{\LTright}{\fill}
\noindent
\begin{tabular}{@{}>{\raggedright\arraybackslash}p{31.32mm}>{\raggedright\arraybackslash}p{34.92mm}>{\raggedright\arraybackslash}p{52.14mm}@{}}
\toprule
\textbf{Form} & \textbf{Meaning} & \textbf{Role} \\
\midrule
nāhtes & nothing; naught & genitive, singular, neuter \\
būtan & except & as used here \\
\bottomrule
\end{tabular}
\par\endgroup

% mdtable:T69:end
% output: EX-XII-XIII-33-06-O001
\reserve{5\baselineskip}
\par\begin{minipage}{\linewidth}\raggedright
% mdparagraph:P359:begin
\textbf{Cook XII, XIII.33.7}
% mdparagraph:P359:end

% item: EX-XII-XIII-33-07
% mdparagraph:P360:begin
Thou wilt cease from thy work (156. k).
% mdparagraph:P360:end

% output: EX-XII-XIII-33-07-O001
\end{minipage}\par\addvspace{4pt}

\par\begin{minipage}{\linewidth}\raggedright
% mdparagraph:P361:begin
\textbf{Cook XII, XIII.33.8}
% mdparagraph:P361:end

% item: EX-XII-XIII-33-08
% mdparagraph:P362:begin
Perhaps he will find a friend (194. a).
% mdparagraph:P362:end

% output: EX-XII-XIII-33-08-O001
\end{minipage}\par\addvspace{4pt}

\par\begin{minipage}{\linewidth}\raggedright
% mdparagraph:P363:begin
\textbf{Cook XII, XIII.33.9}
% mdparagraph:P363:end

% item: EX-XII-XIII-33-09
% mdparagraph:P364:begin
He said to them, I will (139) go a fishing; they said, And we will go with thee.
% mdparagraph:P364:end

% mdparagraph:P365:begin
\textbf{Supplied help:} on fiscnoð
% mdparagraph:P365:end

% mdparagraph:P366:begin
a-fishing; engaged in fishing. Keep the source spelling fiscnoð, not unmarked replacement fiscnað.
% mdparagraph:P366:end

\end{minipage}\par\addvspace{4pt}

% teaching: NEW-U16-SUP-FISCNOTH
% teaching: NEW-U16-VOC-FISCNOTH
\reserve{9\baselineskip}
% mdparagraph:P367:begin
\textbf{Supplied words.} These forms and meanings remain available for this task.
% mdparagraph:P367:end

% mdtable:T70:begin
\par\begingroup
\fontsize{10.5}{13.2}\selectfont
\setlength{\LTleft}{0pt}\setlength{\LTright}{\fill}
\noindent
\begin{tabular}{@{}>{\raggedright\arraybackslash}p{21.1mm}>{\raggedright\arraybackslash}p{55.07mm}>{\raggedright\arraybackslash}p{42.2mm}@{}}
\toprule
\textbf{Form} & \textbf{Meaning} & \textbf{Role} \\
\midrule
on & in the fishing activity, in supplied on fiscnoð & as used here \\
fiscnoð & fishing & accusative, singular, masculine \\
\bottomrule
\end{tabular}
\par\endgroup

% mdtable:T70:end
% output: EX-XII-XIII-33-09-O001
\reserve{5\baselineskip}
\par\begin{minipage}{\linewidth}\raggedright
% mdparagraph:P368:begin
\textbf{Cook XII, XIII.33.10}
% mdparagraph:P368:end

% item: EX-XII-XIII-33-10
% mdparagraph:P369:begin
They went out (ūt), and went into (om. a) ship, and caught nothing (183) in (on) that night.
% mdparagraph:P369:end

% mdparagraph:P370:begin
\textbf{Supplied words.} These forms and meanings remain available for this task.
% mdparagraph:P370:end

% mdtable:T71:begin
\par\begingroup
\fontsize{10.5}{13.2}\selectfont
\setlength{\LTleft}{0pt}\setlength{\LTright}{\fill}
\noindent
\begin{tabular}{@{}>{\raggedright\arraybackslash}p{34.16mm}>{\raggedright\arraybackslash}p{50.05mm}>{\raggedright\arraybackslash}p{34.16mm}@{}}
\toprule
\textbf{Form} & \textbf{Meaning} & \textbf{Role} \\
\midrule
on & during, with dative time & as used here \\
\bottomrule
\end{tabular}
\par\endgroup

% mdtable:T71:end
% output: EX-XII-XIII-33-10-O001
\end{minipage}\par\addvspace{4pt}

\reserve{7\baselineskip}
% mdparagraph:P371:begin
\subsection*{Exercise 34}
% mdparagraph:P371:end

\reserve{5\baselineskip}
\par\begin{minipage}{\linewidth}\raggedright
% mdparagraph:P372:begin
\textbf{Cook XII, XIII.34.1}
% mdparagraph:P372:end

% item: EX-XII-XIII-34-01
% mdparagraph:P373:begin
Whoever you are (193. c), forsake me not.
% mdparagraph:P373:end

% output: EX-XII-XIII-34-01-O001
\end{minipage}\par\addvspace{4pt}

\par\begin{minipage}{\linewidth}\raggedright
% mdparagraph:P374:begin
\textbf{Cook XII, XIII.34.2}
% mdparagraph:P374:end

% item: EX-XII-XIII-34-02
% mdparagraph:P375:begin
If you find nobody (196. d) that will feed you, return hither.
% mdparagraph:P375:end

% output: EX-XII-XIII-34-02-O001
\end{minipage}\par\addvspace{4pt}

\par\begin{minipage}{\linewidth}\raggedright
% mdparagraph:P376:begin
\textbf{Cook XII, XIII.34.3}
% mdparagraph:P376:end

% item: EX-XII-XIII-34-03
% mdparagraph:P377:begin
Let my books be sufficient for both of us (see p. 167, l. 13).
% mdparagraph:P377:end

% mdparagraph:P378:begin
\textbf{Supplied help:} Genyhtsumien mīna bēc unc bǣm.
% mdparagraph:P378:end

% mdparagraph:P379:begin
Let my books suffice for us both.
% mdparagraph:P379:end

% mdparagraph:P380:begin
The verb is plural optative, agreeing with bēc. Unc is dual dative; bǣm is dative plural because both has no separate dual form. This whole sentence remains supplied.
% mdparagraph:P380:end

% mdparagraph:P381:begin
\textbf{Supplied words.} These forms and meanings remain available for this task.
% mdparagraph:P381:end

% mdtable:T72:begin
\par\begingroup
\fontsize{10.5}{13.2}\selectfont
\setlength{\LTleft}{0pt}\setlength{\LTright}{\fill}
\noindent
\begin{tabular}{@{}>{\raggedright\arraybackslash}p{25.97mm}>{\raggedright\arraybackslash}p{37.22mm}>{\raggedright\arraybackslash}p{55.18mm}@{}}
\toprule
\textbf{Form} & \textbf{Meaning} & \textbf{Role} \\
\midrule
genyhtsumien & suffice; be sufficient for & third person, present, subjunctive, plural \\
\bottomrule
\end{tabular}
\par\endgroup

% mdtable:T72:end
% output: EX-XII-XIII-34-03-O001
\end{minipage}\par\addvspace{4pt}

\par\begin{minipage}{\linewidth}\raggedright
% mdparagraph:P382:begin
\textbf{Cook XII, XIII.34.4}
% mdparagraph:P382:end

% item: EX-XII-XIII-34-04
% mdparagraph:P383:begin
Both the boys are well educated.
% mdparagraph:P383:end

% output: EX-XII-XIII-34-04-O001
\end{minipage}\par\addvspace{4pt}

\par\begin{minipage}{\linewidth}\raggedright
% mdparagraph:P384:begin
\textbf{Cook XII, XIII.34.5}
% mdparagraph:P384:end

% item: EX-XII-XIII-34-05
% mdparagraph:P385:begin
They were both righteous before (beforan, 166) God.
% mdparagraph:P385:end

% output: EX-XII-XIII-34-05-O001
\end{minipage}\par\addvspace{4pt}

\par\begin{minipage}{\linewidth}\raggedright
% mdparagraph:P386:begin
\textbf{Cook XII, XIII.34.6}
% mdparagraph:P386:end

% item: EX-XII-XIII-34-06
% mdparagraph:P387:begin
I conjure thee by (ðurh) the living God, that thou tell us if thou be Christ, (om. the) Son of God.
% mdparagraph:P387:end

% mdparagraph:P388:begin
\textbf{Supplied help:} ðurh ðone libbendan God
% mdparagraph:P388:end

% mdparagraph:P389:begin
by the living God
% mdparagraph:P389:end

% mdparagraph:P390:begin
Libbendan is the supplied weak masculine accusative singular participle agreeing with God after ðone. The finite forms of libban do not supply this adjective position.
% mdparagraph:P390:end

% mdparagraph:P391:begin
\textbf{Supplied words.} These forms and meanings remain available for this task.
% mdparagraph:P391:end

% mdtable:T73:begin
\par\begingroup
\fontsize{10.5}{13.2}\selectfont
\setlength{\LTleft}{0pt}\setlength{\LTright}{\fill}
\noindent
\begin{tabular}{@{}>{\raggedright\arraybackslash}p{16.74mm}>{\raggedright\arraybackslash}p{47.8mm}>{\raggedright\arraybackslash}p{53.83mm}@{}}
\toprule
\textbf{Form} & \textbf{Meaning} & \textbf{Role} \\
\midrule
libbendan & living, in the supplied weak participial adjective & present, participle, accusative, singular, masculine, weak \\
\bottomrule
\end{tabular}
\par\endgroup

% mdtable:T73:end
% output: EX-XII-XIII-34-06-O001
\end{minipage}\par\addvspace{4pt}

\par\begin{minipage}{\linewidth}\raggedright
% mdparagraph:P392:begin
\textbf{Cook XII, XIII.34.7}
% mdparagraph:P392:end

% item: EX-XII-XIII-34-07
% mdparagraph:P393:begin
I tell you, After these (om. things), ye shall see the Son of Man sitting on (on, 172) the right hand of (om. the) majesty of God, and coming in (om. the) clouds of heaven.
% mdparagraph:P393:end

% mdparagraph:P394:begin
\textbf{Supplied help:} on ðā swīðran hand Godes mægenðrymmes
% mdparagraph:P394:end

% mdparagraph:P395:begin
at the right hand of God’s majesty
% mdparagraph:P395:end

% mdparagraph:P396:begin
On has an accusative in this fixed position expression; it does not express motion by the seated person. The perceived ongoing actions use plain infinitives sittan and cuman. English -ing does not automatically require an Old English participle.
% mdparagraph:P396:end

% mdparagraph:P397:begin
\textbf{Supplied words.} These forms and meanings remain available for this task.
% mdparagraph:P397:end

% mdtable:T74:begin
\par\begingroup
\fontsize{10.5}{13.2}\selectfont
\setlength{\LTleft}{0pt}\setlength{\LTright}{\fill}
\noindent
\begin{tabular}{@{}>{\raggedright\arraybackslash}p{19.83mm}>{\raggedright\arraybackslash}p{78.72mm}>{\raggedright\arraybackslash}p{19.83mm}@{}}
\toprule
\textbf{Form} & \textbf{Meaning} & \textbf{Role} \\
\midrule
on & at a right/\allowbreak{}left side, with accusative hand in this bounded source expression & as used here \\
\bottomrule
\end{tabular}
\par\endgroup

% mdtable:T74:end
% output: EX-XII-XIII-34-07-O001
\end{minipage}\par\addvspace{4pt}

\par\begin{minipage}{\linewidth}\raggedright
% mdparagraph:P398:begin
\textbf{Cook XII, XIII.34.8}
% mdparagraph:P398:end

% item: EX-XII-XIII-34-08
% mdparagraph:P399:begin
Be (om. it unto) you two (dual, 81) according to (æfter) your faith.
% mdparagraph:P399:end

% output: EX-XII-XIII-34-08-O001
\end{minipage}\par\addvspace{4pt}

\par\begin{minipage}{\linewidth}\raggedright
% mdparagraph:P400:begin
\textbf{Cook XII, XIII.34.9}
% mdparagraph:P400:end

% item: EX-XII-XIII-34-09
% mdparagraph:P401:begin
Truly the Son of Man is to come (200) in (on, 166) his Father’s glory (wuldor), and then (ðonne) (om. shall) reward (āgyltan, 113) every one (164) according to (be) his own work.
% mdparagraph:P401:end

% mdparagraph:P402:begin
\textbf{Supplied help:} is tō cumanne
% mdparagraph:P402:end

% mdparagraph:P403:begin
is to come
% mdparagraph:P403:end

% mdparagraph:P404:begin
This supplied wesan + tō + inflected-infinitive phrase expresses impending coming. Knowing the ordinary purpose infinitive does not by itself establish this use. Retain the complete expression here.
% mdparagraph:P404:end

% mdparagraph:P405:begin
\textbf{Supplied help:} ðonne āgylt hē ǣlcum be his āgnum weorce
% mdparagraph:P405:end

% mdparagraph:P406:begin
then he rewards everyone according to his own work
% mdparagraph:P406:end

% mdparagraph:P407:begin
Āgylt is the contracted third singular of strong āgyldan, and ǣlcum is the dative recipient. The whole clause remains supplied; see the correction above.
% mdparagraph:P407:end

\end{minipage}\par\addvspace{4pt}

% teaching: NEW-U16-SUP-AGEN
% teaching: NEW-U16-VOC-AGEN
\reserve{10\baselineskip}
% mdparagraph:P408:begin
\textbf{Supplied words.} These forms and meanings remain available for this task.
% mdparagraph:P408:end

% mdtable:T75:begin
\par\begingroup
\fontsize{10.5}{13.2}\selectfont
\setlength{\LTleft}{0pt}\setlength{\LTright}{\fill}
\noindent
\begin{tabular}{@{}>{\raggedright\arraybackslash}p{21.06mm}>{\raggedright\arraybackslash}p{43.81mm}>{\raggedright\arraybackslash}p{53.5mm}@{}}
\toprule
\textbf{Form} & \textbf{Meaning} & \textbf{Role} \\
\midrule
āgylt & reward; repay & third person, present, indicative, singular \\
āgnum & own & dative, singular, neuter \\
on & in, in the expression in his glory & as used here \\
\bottomrule
\end{tabular}
\par\endgroup

% mdtable:T75:end
% output: EX-XII-XIII-34-09-O001
\reserve{5\baselineskip}
\par\begin{minipage}{\linewidth}\raggedright
% mdparagraph:P409:begin
\textbf{Cook XII, XIII.34.10}
% mdparagraph:P409:end

% item: EX-XII-XIII-34-10
% mdparagraph:P410:begin
When (ðonne) Christ cometh, then knoweth no man whence he is.
% mdparagraph:P410:end

% output: EX-XII-XIII-34-10-O001
\end{minipage}\par\addvspace{4pt}

\par\begin{minipage}{\linewidth}\raggedright
% mdparagraph:P411:begin
\textbf{Cook XII, XIII.34.11}
% mdparagraph:P411:end

% item: EX-XII-XIII-34-11
% mdparagraph:P412:begin
Jesus cried (cleopian) and taught (lǣran) in the temple, and said, Me ye know (cunnan), and ye know (witan) whence I am; and I came not from myself [me self], but he is true that (87. d) sent me, whom (87. a) ye know not.
% mdparagraph:P412:end

% mdparagraph:P413:begin
\textbf{Supplied help:} sōð
% mdparagraph:P413:end

% mdparagraph:P414:begin
True: supplied nominative masculine singular in predicate use. No complete new adjective system is required.
% mdparagraph:P414:end

\end{minipage}\par\addvspace{4pt}

% teaching: NEW-U16-SUP-SOTH-CITATION
% teaching: NEW-U16-VOC-SOTH-CITATION
\reserve{9\baselineskip}
% mdparagraph:P415:begin
\textbf{Supplied words.} These forms and meanings remain available for this task.
% mdparagraph:P415:end

% mdtable:T76:begin
\par\begingroup
\fontsize{10.5}{13.2}\selectfont
\setlength{\LTleft}{0pt}\setlength{\LTright}{\fill}
\noindent
\begin{tabular}{@{}>{\raggedright\arraybackslash}p{29.59mm}>{\raggedright\arraybackslash}p{29.59mm}>{\raggedright\arraybackslash}p{59.19mm}@{}}
\toprule
\textbf{Form} & \textbf{Meaning} & \textbf{Role} \\
\midrule
sōð & true & nominative, singular, masculine \\
sē & that; those & nominative, singular, masculine \\
\bottomrule
\end{tabular}
\par\endgroup

% mdtable:T76:end
% output: EX-XII-XIII-34-11-O001
\reserve{5\baselineskip}
\par\begin{minipage}{\linewidth}\raggedright
% mdparagraph:P416:begin
\textbf{Cook XII, XIII.34.12}
% mdparagraph:P416:end

% item: EX-XII-XIII-34-12
% mdparagraph:P417:begin
Ye have not known (pret.) him; I know him, and if I say that I know (194) him not (verb last), I (om. shall) be like (om. unto) you (165); but I know him, and I keep (healdan) his saying (sprǣc).
% mdparagraph:P417:end

% output: EX-XII-XIII-34-12-O001
\end{minipage}\par\addvspace{4pt}

\par\begin{minipage}{\linewidth}\raggedright
% mdparagraph:P418:begin
\textbf{Cook XII, XIII.34.13}
% mdparagraph:P418:end

% item: EX-XII-XIII-34-13
% mdparagraph:P419:begin
If I glorify myself, my glory is (183) naught (verb first); my Father (om. it) is, that (87. d) glorifieth me, concerning whom (87. a) ye say (cweðan) that he is (194) your God.
% mdparagraph:P419:end

% mdparagraph:P420:begin
\textbf{Supplied help:} nāwiht
% mdparagraph:P420:end

% mdparagraph:P421:begin
Nothing; naught. Here nāwiht is nominative in the predicate after nis, unlike its earlier accusative object use.
% mdparagraph:P421:end

\end{minipage}\par\addvspace{4pt}

% teaching: NEW-U16-SUP-NAWIHT
% teaching: NEW-U16-VOC-NAWIHT
\reserve{9\baselineskip}
% mdparagraph:P422:begin
\textbf{Supplied words.} These forms and meanings remain available for this task.
% mdparagraph:P422:end

% mdtable:T77:begin
\par\begingroup
\fontsize{10.5}{13.2}\selectfont
\setlength{\LTleft}{0pt}\setlength{\LTright}{\fill}
\noindent
\begin{tabular}{@{}>{\raggedright\arraybackslash}p{28.77mm}>{\raggedright\arraybackslash}p{32.08mm}>{\raggedright\arraybackslash}p{57.53mm}@{}}
\toprule
\textbf{Form} & \textbf{Meaning} & \textbf{Role} \\
\midrule
nāwiht & nothing; naught & nominative, singular, neuter \\
sē & that; those & nominative, singular, masculine \\
\bottomrule
\end{tabular}
\par\endgroup

% mdtable:T77:end
% output: EX-XII-XIII-34-13-O001
\reserve{7\baselineskip}
% mdparagraph:P423:begin
\subsection*{Exercise 35}
% mdparagraph:P423:end

\reserve{5\baselineskip}
\par\begin{minipage}{\linewidth}\raggedright
% mdparagraph:P424:begin
\textbf{Cook XII, XIII.35.1}
% mdparagraph:P424:end

% item: EX-XII-XIII-35-01
% mdparagraph:P425:begin
Catch this ball.
% mdparagraph:P425:end

% output: EX-XII-XIII-35-01-O001
\end{minipage}\par\addvspace{4pt}

\par\begin{minipage}{\linewidth}\raggedright
% mdparagraph:P426:begin
\textbf{Cook XII, XIII.35.2}
% mdparagraph:P426:end

% item: EX-XII-XIII-35-02
% mdparagraph:P427:begin
Who caught the black ball?
% mdparagraph:P427:end

% output: EX-XII-XIII-35-02-O001
\end{minipage}\par\addvspace{4pt}

\par\begin{minipage}{\linewidth}\raggedright
% mdparagraph:P428:begin
\textbf{Cook XII, XIII.35.3}
% mdparagraph:P428:end

% item: EX-XII-XIII-35-03
% mdparagraph:P429:begin
Go the way that (87. d) you came.
% mdparagraph:P429:end

% output: EX-XII-XIII-35-03-O001
\end{minipage}\par\addvspace{4pt}

\par\begin{minipage}{\linewidth}\raggedright
% mdparagraph:P430:begin
\textbf{Cook XII, XIII.35.4}
% mdparagraph:P430:end

% item: EX-XII-XIII-35-04
% mdparagraph:P431:begin
I never played better than I did to-day.
% mdparagraph:P431:end

% output: EX-XII-XIII-35-04-O001
\end{minipage}\par\addvspace{4pt}

\par\begin{minipage}{\linewidth}\raggedright
% mdparagraph:P432:begin
\textbf{Cook XII, XIII.35.5}
% mdparagraph:P432:end

% item: EX-XII-XIII-35-05
% mdparagraph:P433:begin
She never sang better than she did now.
% mdparagraph:P433:end

% output: EX-XII-XIII-35-05-O001
\end{minipage}\par\addvspace{4pt}

\par\begin{minipage}{\linewidth}\raggedright
% mdparagraph:P434:begin
\textbf{Cook XII, XIII.35.6}
% mdparagraph:P434:end

% item: EX-XII-XIII-35-06
% mdparagraph:P435:begin
Thrice (ðrīwa) I suffered shipwreck; a night (170) and a day I was in this sea.
% mdparagraph:P435:end

% output: EX-XII-XIII-35-06-O001
\end{minipage}\par\addvspace{4pt}

\par\begin{minipage}{\linewidth}\raggedright
% mdparagraph:P436:begin
\textbf{Cook XII, XIII.35.7}
% mdparagraph:P436:end

% item: EX-XII-XIII-35-07
% mdparagraph:P437:begin
I knew (witan) a certain man that (87. d) was carried (lǣdan) unto (ōð, 172) the third heaven, and there (ðǣr) heard the words that no earthly man may (not) speak (sprecan).
% mdparagraph:P437:end

% mdparagraph:P438:begin
\textbf{Supplied help:} ōð ðone ðriddan heofon
% mdparagraph:P438:end

% mdparagraph:P439:begin
as far as the third heaven; exact preposition plus accusative goal supplied.
% mdparagraph:P439:end

\end{minipage}\par\addvspace{4pt}

% teaching: NEW-U16-SUP-OTH
% teaching: NEW-U16-VOC-OTH
\reserve{9\baselineskip}
% mdparagraph:P440:begin
\textbf{Supplied words.} These forms and meanings remain available for this task.
% mdparagraph:P440:end

% mdtable:T78:begin
\par\begingroup
\fontsize{10.5}{13.2}\selectfont
\setlength{\LTleft}{0pt}\setlength{\LTright}{\fill}
\noindent
\begin{tabular}{@{}>{\raggedright\arraybackslash}p{22.74mm}>{\raggedright\arraybackslash}p{50.15mm}>{\raggedright\arraybackslash}p{45.49mm}@{}}
\toprule
\textbf{Form} & \textbf{Meaning} & \textbf{Role} \\
\midrule
ōð & as far as; up to & invariable \\
sē & who; the one who, referring to the man & nominative, singular, masculine \\
\bottomrule
\end{tabular}
\par\endgroup

% mdtable:T78:end
% output: EX-XII-XIII-35-07-O001
\reserve{5\baselineskip}
\par\begin{minipage}{\linewidth}\raggedright
% mdparagraph:P441:begin
\textbf{Cook XII, XIII.35.8}
% mdparagraph:P441:end

% item: EX-XII-XIII-35-08
% mdparagraph:P442:begin
O, how narrow (nearu, 57. 5) is the gate and the way that (87. d) leads to life (verb last), and very few (om. there) are that find the way.
% mdparagraph:P442:end

% mdparagraph:P443:begin
\textbf{Supplied help:} hū nearu is ðæt geat and se weg ðe tō līfe lǣtt
% mdparagraph:P443:end

% mdparagraph:P444:begin
how narrow is the gate and the way that leads to life
% mdparagraph:P444:end

% mdparagraph:P445:begin
This wording treats gate and way as a combined figurative route and retains a singular predicate. Use this supplied coordination here; it is not a general singular-agreement rule for coordinated subjects.
% mdparagraph:P445:end

% mdparagraph:P446:begin
\textbf{Supplied help:} fēawe
% mdparagraph:P446:end

% mdparagraph:P447:begin
few people; masculine nominative plural substantive supplied, with no general new adjective paradigm.
% mdparagraph:P447:end

\end{minipage}\par\addvspace{4pt}

% teaching: NEW-U16-SUP-FEAWE
% teaching: NEW-U16-VOC-FEAWE
\reserve{8\baselineskip}
% mdparagraph:P448:begin
\textbf{Supplied words.} These forms and meanings remain available for this task.
% mdparagraph:P448:end

% mdtable:T79:begin
\par\begingroup
\fontsize{10.5}{13.2}\selectfont
\setlength{\LTleft}{0pt}\setlength{\LTright}{\fill}
\noindent
\begin{tabular}{@{}>{\raggedright\arraybackslash}p{30.28mm}>{\raggedright\arraybackslash}p{30.28mm}>{\raggedright\arraybackslash}p{57.83mm}@{}}
\toprule
\textbf{Form} & \textbf{Meaning} & \textbf{Role} \\
\midrule
fēawe & few people & nominative, plural, masculine \\
\bottomrule
\end{tabular}
\par\endgroup

% mdtable:T79:end
% output: EX-XII-XIII-35-08-O001
\reserve{5\baselineskip}
\par\begin{minipage}{\linewidth}\raggedright
% mdparagraph:P449:begin
\textbf{Cook XII, XIII.35.9}
% mdparagraph:P449:end

% item: EX-XII-XIII-35-09
% mdparagraph:P450:begin
They come to (tō, 166) you in sheep’s garments (gegierela), but within (innan) they are wolves.
% mdparagraph:P450:end

% mdparagraph:P451:begin
\textbf{Supplied help:} on scēapa gegierelum
% mdparagraph:P451:end

% mdparagraph:P452:begin
in sheep’s clothing; the local clothing use of on is supplied.
% mdparagraph:P452:end

% mdparagraph:P453:begin
\textbf{Supplied words.} These forms and meanings remain available for this task.
% mdparagraph:P453:end

% mdtable:T80:begin
\par\begingroup
\fontsize{10.5}{13.2}\selectfont
\setlength{\LTleft}{0pt}\setlength{\LTright}{\fill}
\noindent
\begin{tabular}{@{}>{\raggedright\arraybackslash}p{23.77mm}>{\raggedright\arraybackslash}p{70.84mm}>{\raggedright\arraybackslash}p{23.77mm}@{}}
\toprule
\textbf{Form} & \textbf{Meaning} & \textbf{Role} \\
\midrule
on & in; wearing clothing, in supplied on scēapa gegierelum & as used here \\
\bottomrule
\end{tabular}
\par\endgroup

% mdtable:T80:end
% output: EX-XII-XIII-35-09-O001
\end{minipage}\par\addvspace{4pt}

\par\begin{minipage}{\linewidth}\raggedright
% mdparagraph:P454:begin
\textbf{Cook XII, XIII.35.10}
% mdparagraph:P454:end

% item: EX-XII-XIII-35-10
% mdparagraph:P455:begin
When he came nigh (nēalǣcan) to the gate (163) of the city, then (ðā) was there (ðǣr) a dead man carried (beran).
% mdparagraph:P455:end

% output: EX-XII-XIII-35-10-O001
\end{minipage}\par\addvspace{4pt}

\par\begin{minipage}{\linewidth}\raggedright
% mdparagraph:P456:begin
\textbf{Cook XII, XIII.35.11}
% mdparagraph:P456:end

% item: EX-XII-XIII-35-11
% mdparagraph:P457:begin
I see a girl running (169) along (172) the street.
% mdparagraph:P457:end

% output: EX-XII-XIII-35-11-O001
\end{minipage}\par\addvspace{4pt}

\par\begin{minipage}{\linewidth}\raggedright
% mdparagraph:P458:begin
\textbf{Cook XII, XIII.35.12}
% mdparagraph:P458:end

% item: EX-XII-XIII-35-12
% mdparagraph:P459:begin
Anoint me with oil.
% mdparagraph:P459:end

% output: EX-XII-XIII-35-12-O001
\end{minipage}\par\addvspace{4pt}

\par\begin{minipage}{\linewidth}\raggedright
% mdparagraph:P460:begin
\textbf{Cook XII, XIII.35.13}
% mdparagraph:P460:end

% item: EX-XII-XIII-35-13
% mdparagraph:P461:begin
Give (sęllan) us of (of, 166) your oil.
% mdparagraph:P461:end

% mdparagraph:P462:begin
\textbf{Supplied help:} of ēowrum ęle
% mdparagraph:P462:end

% mdparagraph:P463:begin
of your oil; source partition/\allowbreak{}substance use of of is supplied.
% mdparagraph:P463:end

% mdparagraph:P464:begin
\textbf{Supplied words.} These forms and meanings remain available for this task.
% mdparagraph:P464:end

% mdtable:T81:begin
\par\begingroup
\fontsize{10.5}{13.2}\selectfont
\setlength{\LTleft}{0pt}\setlength{\LTright}{\fill}
\noindent
\begin{tabular}{@{}>{\raggedright\arraybackslash}p{19.31mm}>{\raggedright\arraybackslash}p{79.75mm}>{\raggedright\arraybackslash}p{19.31mm}@{}}
\toprule
\textbf{Form} & \textbf{Meaning} & \textbf{Role} \\
\midrule
of & of; from a portion of a whole, in the supplied phrases ān of ðǣm and of ēowrum ęle & as used here \\
\bottomrule
\end{tabular}
\par\endgroup

% mdtable:T81:end
% output: EX-XII-XIII-35-13-O001
\end{minipage}\par\addvspace{4pt}

\par\begin{minipage}{\linewidth}\raggedright
% mdparagraph:P465:begin
\textbf{Cook XII, XIII.35.14}
% mdparagraph:P465:end

% item: EX-XII-XIII-35-14
% mdparagraph:P466:begin
If the Son delivers you, ye (om. shall) indeed be free.
% mdparagraph:P466:end

% mdparagraph:P467:begin
\textbf{Supplied help:} frige
% mdparagraph:P467:end

% mdparagraph:P468:begin
free, nominative plural; Cook’s glossary explicitly gives this irregular plural.
% mdparagraph:P468:end

\end{minipage}\par\addvspace{4pt}

% teaching: NEW-U16-SUP-FREO
% teaching: NEW-U16-VOC-FREO
\reserve{8\baselineskip}
% mdparagraph:P469:begin
\textbf{Supplied words.} These forms and meanings remain available for this task.
% mdparagraph:P469:end

% mdtable:T82:begin
\par\begingroup
\fontsize{10.5}{13.2}\selectfont
\setlength{\LTleft}{0pt}\setlength{\LTright}{\fill}
\noindent
\begin{tabular}{@{}>{\raggedright\arraybackslash}p{30.28mm}>{\raggedright\arraybackslash}p{30.28mm}>{\raggedright\arraybackslash}p{57.83mm}@{}}
\toprule
\textbf{Form} & \textbf{Meaning} & \textbf{Role} \\
\midrule
frige & free & nominative, plural, masculine \\
\bottomrule
\end{tabular}
\par\endgroup

% mdtable:T82:end
% output: EX-XII-XIII-35-14-O001
\reserve{7\baselineskip}
% mdparagraph:P470:begin
\subsection*{Exercise 36}
% mdparagraph:P470:end

\reserve{5\baselineskip}
\par\begin{minipage}{\linewidth}\raggedright
% mdparagraph:P471:begin
\textbf{Cook XII, XIII.36.1}
% mdparagraph:P471:end

% item: EX-XII-XIII-36-01
% mdparagraph:P472:begin
Judge ye not according to (be) appearance (ansīen, 51. b), but judge right judgment (168. 1).
% mdparagraph:P472:end

% output: EX-XII-XIII-36-01-O001
\end{minipage}\par\addvspace{4pt}

\par\begin{minipage}{\linewidth}\raggedright
% mdparagraph:P473:begin
\textbf{Cook XII, XIII.36.2}
% mdparagraph:P473:end

% item: EX-XII-XIII-36-02
% mdparagraph:P474:begin
Then turned (bewęndan, 184. a) he to the woman, and said, Seest thou this woman? I went into (intō) thy house, thou gavest me not water for (tō) my feet; this (om. one hath) washed (101. VI.) my feet with (mid) her tears.
% mdparagraph:P474:end

% mdparagraph:P475:begin
\textbf{Supplied help:} ðēos
% mdparagraph:P475:end

% mdparagraph:P476:begin
This woman. The feminine pronoun refers to the woman as a person, although wīf is grammatically neuter. Retain this supplied reference in the exchange.
% mdparagraph:P476:end

% mdparagraph:P477:begin
\textbf{Supplied help:} wæter tō mīnum fōtum
% mdparagraph:P477:end

% mdparagraph:P478:begin
water for my feet: the purpose of washing/\allowbreak{}caring for the body part is supplied; tō does not mean motion toward a destination here.
% mdparagraph:P478:end

\end{minipage}\par\addvspace{4pt}

% teaching: NEW-U16-VOC-TO
\reserve{8\baselineskip}
% mdparagraph:P479:begin
\textbf{Supplied words.} These forms and meanings remain available for this task.
% mdparagraph:P479:end

% mdtable:T83:begin
\par\begingroup
\fontsize{10.5}{13.2}\selectfont
\setlength{\LTleft}{0pt}\setlength{\LTright}{\fill}
\noindent
\begin{tabular}{@{}>{\raggedright\arraybackslash}p{22.83mm}>{\raggedright\arraybackslash}p{72.72mm}>{\raggedright\arraybackslash}p{22.83mm}@{}}
\toprule
\textbf{Form} & \textbf{Meaning} & \textbf{Role} \\
\midrule
tō & for the care of a body part, in supplied wæter tō mīnum fōtum & as used here \\
\bottomrule
\end{tabular}
\par\endgroup

% mdtable:T83:end
% output: EX-XII-XIII-36-02-O001
\reserve{5\baselineskip}
\par\begin{minipage}{\linewidth}\raggedright
% mdparagraph:P480:begin
\textbf{Cook XII, XIII.36.3}
% mdparagraph:P480:end

% item: EX-XII-XIII-36-03
% mdparagraph:P481:begin
Thou gavest (verb last) me not (om. a) kiss (coss, 43); this (om. one), since (siððan) I entered, (om. hath) ceased not to kiss [that she kissed] my feet.
% mdparagraph:P481:end

% mdparagraph:P482:begin
\textbf{Supplied help:} ne geswāc ðæt hēo mīne fēt ne cyste
% mdparagraph:P482:end

% mdparagraph:P483:begin
did not stop kissing my feet
% mdparagraph:P483:end

% mdparagraph:P484:begin
Both negatives belong to this attested cessation construction: the second ne does not mean that she withheld kisses. Keep the complete finite clause supplied here.
% mdparagraph:P484:end

% mdparagraph:P485:begin
\textbf{Supplied help:} ðēos
% mdparagraph:P485:end

% mdparagraph:P486:begin
This woman. The feminine pronoun refers to the woman as a person, although wīf is grammatically neuter. Retain this supplied reference in the exchange.
% mdparagraph:P486:end

% output: EX-XII-XIII-36-03-O001
\end{minipage}\par\addvspace{4pt}

\par\begin{minipage}{\linewidth}\raggedright
% mdparagraph:P487:begin
\textbf{Cook XII, XIII.36.4}
% mdparagraph:P487:end

% item: EX-XII-XIII-36-04
% mdparagraph:P488:begin
My head thou with oils (om. didst) not anoint (smierwan, 113); this one anointed my feet with ointment (sealf, 51. b).
% mdparagraph:P488:end

% mdparagraph:P489:begin
\textbf{Supplied help:} ðēos
% mdparagraph:P489:end

% mdparagraph:P490:begin
This woman. The feminine pronoun refers to the woman as a person, although wīf is grammatically neuter. Retain this supplied reference in the exchange.
% mdparagraph:P490:end

% mdparagraph:P491:begin
\textbf{Supplied help:} mid ęlum
% mdparagraph:P491:end

% mdparagraph:P492:begin
With oils. The plural dative ęlum remains supplied.
% mdparagraph:P492:end

\end{minipage}\par\addvspace{4pt}

% teaching: NEW-U16-SUP-ELE
% teaching: NEW-U16-VOC-ELE
\reserve{8\baselineskip}
% mdparagraph:P493:begin
\textbf{Supplied words.} These forms and meanings remain available for this task.
% mdparagraph:P493:end

% mdtable:T84:begin
\par\begingroup
\fontsize{10.5}{13.2}\selectfont
\setlength{\LTleft}{0pt}\setlength{\LTright}{\fill}
\noindent
\begin{tabular}{@{}>{\raggedright\arraybackslash}p{32.7mm}>{\raggedright\arraybackslash}p{32.7mm}>{\raggedright\arraybackslash}p{52.97mm}@{}}
\toprule
\textbf{Form} & \textbf{Meaning} & \textbf{Role} \\
\midrule
ęlum & oil & dative, plural, masculine \\
\bottomrule
\end{tabular}
\par\endgroup

% mdtable:T84:end
% output: EX-XII-XIII-36-04-O001
\reserve{5\baselineskip}
\par\begin{minipage}{\linewidth}\raggedright
% mdparagraph:P494:begin
\textbf{Cook XII, XIII.36.5}
% mdparagraph:P494:end

% item: EX-XII-XIII-36-05
% mdparagraph:P495:begin
Therefore I say (sęcgean) to thee (164. b), Many sins are forgiven (forgiefan) her, because she loved (verb last) me much (swīðe).
% mdparagraph:P495:end

% output: EX-XII-XIII-36-05-O001
\end{minipage}\par\addvspace{4pt}

\par\begin{minipage}{\linewidth}\raggedright
% mdparagraph:P496:begin
\textbf{Cook XII, XIII.36.6}
% mdparagraph:P496:end

% item: EX-XII-XIII-36-06
% mdparagraph:P497:begin
Less (neut. acc. sing.) loveth (om. to) whom (87. b) less forgiven is.
% mdparagraph:P497:end

% mdparagraph:P498:begin
\textbf{Supplied help:} Lǣs lufað ðǣm ðe lǣs forgiefen is.
% mdparagraph:P498:end

% mdparagraph:P499:begin
The one to whom less is forgiven loves less.
% mdparagraph:P499:end

% mdparagraph:P500:begin
The matrix subject is understood from the attracted dative relative. The relative dative belongs to forgiveness; it must not be generalized into a dative subject of ordinary lufian. The prompt attaches neuter accusative singular specifically to its first Less: this is the adverbial degree accusative. The second lǣs is nominative neuter singular, the smaller amount that is forgiven, and is the passive subject. The entire compressed sentence is supplied.
% mdparagraph:P500:end

% mdparagraph:P501:begin
\textbf{Supplied help:} lǣs ... lǣs
% mdparagraph:P501:end

% mdparagraph:P502:begin
First lǣs is the specified adverbial accusative neuter singular; second lǣs is the nominative neuter singular amount forgiven. Both exact cells belong to the supplied clause.
% mdparagraph:P502:end

\end{minipage}\par\addvspace{4pt}

% teaching: NEW-U16-SUP-LAES
% teaching: NEW-U16-VOC-LAES
\reserve{9\baselineskip}
% mdparagraph:P503:begin
\textbf{Supplied words.} These forms and meanings remain available for this task.
% mdparagraph:P503:end

% mdtable:T85:begin
\par\begingroup
\fontsize{10.5}{13.2}\selectfont
\setlength{\LTleft}{0pt}\setlength{\LTright}{\fill}
\noindent
\begin{tabular}{@{}>{\raggedright\arraybackslash}p{27.92mm}>{\raggedright\arraybackslash}p{38.39mm}>{\raggedright\arraybackslash}p{52.07mm}@{}}
\toprule
\textbf{Form} & \textbf{Meaning} & \textbf{Role} \\
\midrule
lǣs & less; a smaller amount & accusative, singular, neuter \\
lǣs & less; a smaller amount & nominative, singular, neuter \\
\bottomrule
\end{tabular}
\par\endgroup

% mdtable:T85:end
% output: EX-XII-XIII-36-06-O001
\reserve{5\baselineskip}
\par\begin{minipage}{\linewidth}\raggedright
% mdparagraph:P504:begin
\textbf{Cook XII, XIII.36.7}
% mdparagraph:P504:end

% item: EX-XII-XIII-36-07
% mdparagraph:P505:begin
It (om. is) not permitted (ālīefan) thee to have her to wife.
% mdparagraph:P505:end

% mdparagraph:P506:begin
\textbf{Supplied help:} Hit ne ālīefð ðē hīe tō wīfe habban.
% mdparagraph:P506:end

% mdparagraph:P507:begin
It is not permitted for you to have her as wife.
% mdparagraph:P507:end

% mdparagraph:P508:begin
Impersonal hit with finite ālīefð, dative permitted person ðē, accusative hīe as object of habban, and tō wīfe as the wife-role phrase. Omit the English is; do not add wesan or turn hīe into a dative object.
% mdparagraph:P508:end

% mdparagraph:P509:begin
\textbf{Supplied help:} tō wīfe
% mdparagraph:P509:end

% mdparagraph:P510:begin
as wife; in the role of wife. This is part of the supplied permission construction, not a destination.
% mdparagraph:P510:end

% mdparagraph:P511:begin
\textbf{Supplied words.} These forms and meanings remain available for this task.
% mdparagraph:P511:end

% mdtable:T86:begin
\par\begingroup
\fontsize{10.5}{13.2}\selectfont
\setlength{\LTleft}{0pt}\setlength{\LTright}{\fill}
\noindent
\begin{tabular}{@{}>{\raggedright\arraybackslash}p{18.8mm}>{\raggedright\arraybackslash}p{51.81mm}>{\raggedright\arraybackslash}p{47.76mm}@{}}
\toprule
\textbf{Form} & \textbf{Meaning} & \textbf{Role} \\
\midrule
ālīefð & permit; allow, in impersonal permission & third person, present, indicative, singular \\
tō & as wife; in the role of wife, in supplied tō wīfe & as used here \\
\bottomrule
\end{tabular}
\par\endgroup

% mdtable:T86:end
% output: EX-XII-XIII-36-07-O001
\end{minipage}\par\addvspace{4pt}

\par\begin{minipage}{\linewidth}\raggedright
% mdparagraph:P512:begin
\textbf{Cook XII, XIII.36.8}
% mdparagraph:P512:end

% item: EX-XII-XIII-36-08
% mdparagraph:P513:begin
Then he promised with (mid) oath (āð, 43) to give (sęllan) her whatsoever she (om. should) ask (biddan) him.
% mdparagraph:P513:end

% mdparagraph:P514:begin
\textbf{Supplied help:} Ðā gehēt hē mid āðe hiere ðæt tō sęllanne, swā hwæs swā hēo hine bǣde.
% mdparagraph:P514:end

% mdparagraph:P515:begin
Then he promised with an oath to give her whatever she asked him for.
% mdparagraph:P515:end

% mdparagraph:P516:begin
Gehēt introduces the promise and the intended-action infinitive. Hiere is the recipient of giving, ðæt the accusative thing given. Within the free-choice request, biddan takes genitive hwæs and accusative hine; bǣde is preterite subjunctive in the open request. The resumptive ðæt lets each clause retain its own case requirement.
% mdparagraph:P516:end

\end{minipage}\par\addvspace{4pt}

% teaching: NEW-U16-VOC-MID
\reserve{8\baselineskip}
% mdparagraph:P517:begin
\textbf{Supplied words.} These forms and meanings remain available for this task.
% mdparagraph:P517:end

% mdtable:T87:begin
\par\begingroup
\fontsize{10.5}{13.2}\selectfont
\setlength{\LTleft}{0pt}\setlength{\LTright}{\fill}
\noindent
\begin{tabular}{@{}>{\raggedright\arraybackslash}p{20.53mm}>{\raggedright\arraybackslash}p{77.31mm}>{\raggedright\arraybackslash}p{20.53mm}@{}}
\toprule
\textbf{Form} & \textbf{Meaning} & \textbf{Role} \\
\midrule
mid & with; in the manner of the accompanying quality or mental disposition & as used here \\
\bottomrule
\end{tabular}
\par\endgroup

% mdtable:T87:end
% output: EX-XII-XIII-36-08-O001
\reserve{5\baselineskip}
\par\begin{minipage}{\linewidth}\raggedright
% mdparagraph:P518:begin
\textbf{Cook XII, XIII.36.9}
% mdparagraph:P518:end

% item: EX-XII-XIII-36-09
% mdparagraph:P519:begin
Then, admonished (myngian) by (fram) her mother, she said, Give me in a dish (disc, 43) the head of John (Johannes) the Baptist (fuluhtere, 44. 1).
% mdparagraph:P519:end

% mdparagraph:P520:begin
\textbf{Supplied help:} fram hiere mēder myngod
% mdparagraph:P520:end

% mdparagraph:P521:begin
admonished by her mother
% mdparagraph:P521:end

% mdparagraph:P522:begin
The past participle is supplied uninflected in its predicative attachment to hēo; fram plus dative marks the agent. This does not introduce a complete myngian system.
% mdparagraph:P522:end

% mdparagraph:P523:begin
\textbf{Supplied help:} Johannes /\allowbreak{} Johannem
% mdparagraph:P523:end

% mdparagraph:P524:begin
John: Johannes is genitive and Johannem accusative. The endings are retained from the attested forms; initial J follows Cook’s headword. See the note on names above.
% mdparagraph:P524:end

\end{minipage}\par\addvspace{4pt}

% teaching: NEW-U16-SUP-JOHANNES-NAME
% teaching: NEW-U16-VOC-JOHANNES-NAME
\reserve{10\baselineskip}
% mdparagraph:P525:begin
\textbf{Supplied words.} These forms and meanings remain available for this task.
% mdparagraph:P525:end

% mdtable:T88:begin
\par\begingroup
\fontsize{10.5}{13.2}\selectfont
\setlength{\LTleft}{0pt}\setlength{\LTright}{\fill}
\noindent
\begin{tabular}{@{}>{\raggedright\arraybackslash}p{26.28mm}>{\raggedright\arraybackslash}p{39.55mm}>{\raggedright\arraybackslash}p{52.55mm}@{}}
\toprule
\textbf{Form} & \textbf{Meaning} & \textbf{Role} \\
\midrule
myngod & admonished; prompted & past, participle \\
Johannes & John & genitive, singular, masculine \\
Johannem & John & accusative, singular, masculine \\
\bottomrule
\end{tabular}
\par\endgroup

% mdtable:T88:end
% output: EX-XII-XIII-36-09-O001
\reserve{5\baselineskip}
\par\begin{minipage}{\linewidth}\raggedright
% mdparagraph:P526:begin
\textbf{Cook XII, XIII.36.10}
% mdparagraph:P526:end

% item: EX-XII-XIII-36-10
% mdparagraph:P527:begin
Then he sent (āsęndan), and beheaded (behēafdian, 118) John in the prison.
% mdparagraph:P527:end

% mdparagraph:P528:begin
\textbf{Supplied help:} Johannes /\allowbreak{} Johannem
% mdparagraph:P528:end

% mdparagraph:P529:begin
John: Johannes is genitive and Johannem accusative. The endings are retained from the attested forms; initial J follows Cook’s headword. See the note on names above.
% mdparagraph:P529:end

% mdparagraph:P530:begin
\textbf{Supplied words.} These forms and meanings remain available for this task.
% mdparagraph:P530:end

% mdtable:T89:begin
\par\begingroup
\fontsize{10.5}{13.2}\selectfont
\setlength{\LTleft}{0pt}\setlength{\LTright}{\fill}
\noindent
\begin{tabular}{@{}>{\raggedright\arraybackslash}p{29.59mm}>{\raggedright\arraybackslash}p{29.59mm}>{\raggedright\arraybackslash}p{59.19mm}@{}}
\toprule
\textbf{Form} & \textbf{Meaning} & \textbf{Role} \\
\midrule
Johannes & John & genitive, singular, masculine \\
Johannem & John & accusative, singular, masculine \\
\bottomrule
\end{tabular}
\par\endgroup

% mdtable:T89:end
% output: EX-XII-XIII-36-10-O001
\end{minipage}\par\addvspace{4pt}

\par\begin{minipage}{\linewidth}\raggedright
% mdparagraph:P531:begin
\textbf{Cook XII, XIII.36.11}
% mdparagraph:P531:end

% item: EX-XII-XIII-36-11
% mdparagraph:P532:begin
And one brought his head in a dish.
% mdparagraph:P532:end

% output: EX-XII-XIII-36-11-O001
\end{minipage}\par\addvspace{4pt}

\reserve{7\baselineskip}
% mdparagraph:P533:begin
\section*{7. The Incidents in the Gymnasium}
% mdparagraph:P533:end

% teaching: NEW-U16-READING-SUPPORT
% mdparagraph:P534:begin
Apollonius has survived a shipwreck. A fisherman has shared clothing with him and shown him the road to the city. This scene begins as he follows that road; it ends with the king’s invitation. The following feast is outside the passage.
% mdparagraph:P534:end

% mdparagraph:P535:begin
The text is Cook’s complete scene at printed pages 168.1–170.16. A locator such as 169.6 means printed page 169, editorially counted line 6, anchored by the edition’s printed line numbers. P169N6 instead means original note 6 on page 169. Original notes and modern help are kept separate; a question in an original note is not an additional exercise.\textsuperscript{9}
% mdparagraph:P535:end

\reserve{7\baselineskip}
% mdparagraph:P536:begin
\subsection*{Supplied vocabulary and constructions}
% mdparagraph:P536:end

% mdparagraph:P537:begin
The following help is local to this reading. Retain the learned productive systems where they apply; the additional forms, spellings, meanings and constructions listed here may be consulted throughout the reading tasks.
% mdparagraph:P537:end

% teaching: NEW-U16-SUP-THES
% teaching: NEW-U16-VOC-THES
% teaching: NEW-U16-SUP-OTH-THAET
% teaching: NEW-U16-VOC-OTH-THAET
% teaching: NEW-U16-SUP-MID-THY
% teaching: NEW-U16-VOC-MID-THY
% teaching: NEW-U16-SUP-FULTUM
% teaching: NEW-U16-VOC-FULTUM
% teaching: NEW-U16-SUP-AN
% teaching: NEW-U16-VOC-AN
% teaching: NEW-U16-SUP-SCIETE
% teaching: NEW-U16-VOC-SCIETE
% teaching: NEW-U16-SUP-GEONG
% teaching: NEW-U16-VOC-GEONG
% teaching: NEW-U16-SUP-PLEGA
% teaching: NEW-U16-VOC-PLEGA
% teaching: NEW-U16-SUP-BAETHSTEDE
% teaching: NEW-U16-VOC-BAETHSTEDE
% teaching: NEW-U16-SUP-CEASTERWARAN
% teaching: NEW-U16-VOC-CEASTERWARAN
% teaching: NEW-U16-SUP-ELTHEODIG
% teaching: NEW-U16-VOC-ELTHEODIG
% teaching: NEW-U16-SUP-THEOW
% teaching: NEW-U16-VOC-THEOW
% teaching: NEW-U16-SUP-AETHELE
% teaching: NEW-U16-VOC-AETHELE
% teaching: NEW-U16-SUP-UNAETHELE
% teaching: NEW-U16-VOC-UNAETHELE
% teaching: NEW-U16-SUP-OPEN
% teaching: NEW-U16-VOC-OPEN
% teaching: NEW-U16-SUP-THA-CONJUNCTION
% teaching: NEW-U16-VOC-THA-CONJUNCTION
% teaching: NEW-U16-SUP-APOLLONIUS
% teaching: NEW-U16-VOC-APOLLONIUS
% teaching: NEW-U16-SUP-HEALF
% teaching: NEW-U16-VOC-HEALF
% teaching: NEW-U16-SUP-SCICCELS
% teaching: NEW-U16-VOC-SCICCELS
% teaching: NEW-U16-SUP-ON-ADVERB
% teaching: NEW-U16-VOC-ON-ADVERB
% teaching: NEW-U16-SUP-THWEAL
% teaching: NEW-U16-VOC-THWEAL
% teaching: NEW-U16-SUP-GELICA
% teaching: NEW-U16-VOC-GELICA
% teaching: NEW-U16-SUP-FLOCC
% teaching: NEW-U16-VOC-FLOCC
% teaching: NEW-U16-SUP-FAERINGA
% teaching: NEW-U16-VOC-FAERINGA
% teaching: NEW-U16-SUP-ARCESTRATES
% teaching: NEW-U16-VOC-ARCESTRATES
% teaching: NEW-U16-SUP-THEOD
% teaching: NEW-U16-VOC-THEOD
% teaching: NEW-U16-SUP-MENIGU
% teaching: NEW-U16-VOC-MENIGU
% teaching: NEW-U16-VOC-WITH
% teaching: NEW-U16-SUP-SWIFT
% teaching: NEW-U16-VOC-SWIFT
% teaching: NEW-U16-SUP-HRAEDNES
% teaching: NEW-U16-VOC-HRAEDNES
% teaching: NEW-U16-SUP-ONGEAN
% teaching: NEW-U16-VOC-ONGEAN
% teaching: NEW-U16-SUP-HRAEDLICE
% teaching: NEW-U16-VOC-HRAEDLICE
% teaching: NEW-U16-SUP-SNELNES
% teaching: NEW-U16-VOC-SNELNES
% teaching: NEW-U16-VOC-THAET-CONJUNCTION
% teaching: NEW-U16-SUP-HEONAN
% teaching: NEW-U16-VOC-HEONAN
% teaching: NEW-U16-SUP-CNIHT
% teaching: NEW-U16-VOC-CNIHT
% teaching: NEW-U16-VOC-THAES-THE
% teaching: NEW-U16-SUP-TOP
% teaching: NEW-U16-VOC-TOP
% teaching: NEW-U16-SUP-SWIFTNES
% teaching: NEW-U16-VOC-SWIFTNES
% teaching: NEW-U16-SUP-WESAN
% teaching: NEW-U16-VOC-WESAN
% teaching: NEW-U16-SUP-IELDU
% teaching: NEW-U16-VOC-IELDU
% teaching: NEW-U16-SUP-GEOGUTH
% teaching: NEW-U16-VOC-GEOGUTH
% teaching: NEW-U16-SUP-CYNESETL
% teaching: NEW-U16-VOC-CYNESETL
% teaching: NEW-U16-SUP-GECWEMLICE
% teaching: NEW-U16-VOC-GECWEMLICE
% teaching: NEW-U16-SUP-GEMAENE
% teaching: NEW-U16-VOC-GEMAENE
% teaching: NEW-U16-SUP-HAELU
% teaching: NEW-U16-VOC-HAELU
% teaching: NEW-U16-SUP-THENUNG
% teaching: NEW-U16-VOC-THENUNG
% teaching: NEW-U16-SUP-GESEON
% teaching: NEW-U16-VOC-GESEON
% teaching: NEW-U16-SUP-HORG
% teaching: NEW-U16-VOC-HORG
% teaching: NEW-U16-SUP-FORLIDEN
% teaching: NEW-U16-VOC-FORLIDEN
% teaching: NEW-U16-SUP-GEREORD
% teaching: NEW-U16-VOC-GEREORD
\reserve{28\baselineskip}
% mdparagraph:P538:begin
\textbf{Reading forms and meanings.} These forms and meanings remain available for this task.
% mdparagraph:P538:end

% mdtable:T90:begin
\par\begingroup
\fontsize{10.5}{13.2}\selectfont
\setlength{\LTleft}{0pt}\setlength{\LTright}{\fill}
\begin{longtable}{@{}>{\raggedright\arraybackslash}p{26mm}>{\raggedright\arraybackslash}p{44.5mm}>{\raggedright\arraybackslash}p{47.88mm}@{}}
\toprule
\textbf{Form} & \textbf{Meaning} & \textbf{Role} \\
\midrule
\endfirsthead
\toprule
\textbf{Form} & \textbf{Meaning} & \textbf{Role} \\
\midrule
\endhead
\midrule
\multicolumn{3}{r}{\footnotesize Continued on the following page.}\\
\endfoot
\bottomrule
\endlastfoot
ðisum & this; these & dative, plural, neuter \\
getǣht & shown; directed & past, participle \\
ōð-ðæt & until & invariable \\
becōm & come; arrive & third singular preterite indicative \\
inēode & went in & third person, preterite, indicative, singular \\
mid-ðȳ-ðe & when; while & invariable \\
ðōhte & thought & as used here \\
meahte & might; could & third singular preterite; indicative or subjunctive in form \\
fultumes & help; support & genitive, singular, masculine \\
ǣnne & one & accusative, singular, masculine \\
sē & that; those & nominative, singular, masculine \\
gesmierwed & anointed & past, participle \\
mid & with a cloth or cloak, accompanying girded or wrapped & as used here \\
scīetan & linen cloth & dative, singular, feminine \\
begyrd & girded; clothed around & past, participle \\
geongra & young & genitive, plural, masculine \\
plegan & play; game & accusative, singular, masculine \\
handa & hand; with right/\allowbreak{}left, side & dative, singular \\
tō & to, complement of belonging & invariable \\
bæðstęde & bathing place; here gymnasium & dative, singular, masculine \\
belimpende & belonging; relating & present, participle \\
Gehīere & hear; listen & second person, imperative, plural \\
ceasterwaran & townspeople; citizens & nominative, plural, masculine \\
ęlðēodige & foreign people; strangers & nominative, plural, masculine \\
frige & free & nominative, plural, masculine \\
ðēowe & unfree; enslaved people & nominative, plural, masculine \\
æðele & noble people & nominative, plural, masculine \\
unæðele & people not noble & nominative, plural, masculine \\
bæðstęde & bathing place; here gymnasium & nominative, singular, masculine \\
open & open & nominative, singular, masculine \\
ðā-ðā & when; factual past time & invariable \\
Apollonius & Apollonius & nominative, singular, masculine \\
gehīerde & heard, third singular preterite indicative of prefixed hear & singular \\
unscrȳdde & undressed; took off & third person, preterite, indicative, singular \\
healfan & half & dative, singular, masculine \\
sciccelse & cloak & dative, singular, masculine \\
on & on, in the relative clause referring to the cloak he wore & invariable \\
ðwēale & bath; bathing place & dative, singular, neuter \\
behēold & observed & third person, preterite, indicative, singular \\
ānra & one & genitive, plural, masculine, strong \\
on & at work; engaged in an activity, in exact source phrases & invariable \\
gelīcan & equal; peer & accusative, singular, masculine \\
findan & find & infinitive \\
flocce & group; company & dative, singular, masculine \\
fǣringa & suddenly & invariable \\
Arcestrates & Arcestrates, king of the region & nominative, singular, masculine \\
ðēode & region; country & genitive, singular, feminine \\
męnige & company; crowd & dative, singular, feminine \\
āgan & began & third person, preterite, indicative, singular \\
plegian & play & infinitive \\
wið & with; against fellow players, with dative in the exact source playing frame & as used here \\
gemęngde & joined; mingled & third person, preterite, indicative, singular \\
swā-swā & just as; like & invariable \\
iernende & run & present, participle, nominative, singular, masculine \\
gelæhte & caught & third person, preterite, indicative, singular \\
mid & with; in the manner of the accompanying quality or mental disposition & as used here \\
swiftre & swift & dative, singular, feminine \\
hrædnesse & speed; quickness & dative, singular, feminine \\
geslægene & struck & past participle, masculine accusative singular; exceptional ending (P169N2) \\
ongēan & back; again & invariable \\
gesęnde & sent back & third person, preterite, indicative, singular \\
plegiendan & play & present, participle, dative, singular, masculine, weak \\
āgēan & back; again & invariable \\
hrædlīce & quickly & invariable \\
slōh & struck & third person, preterite, indicative, singular \\
feallan & fall & infinitive \\
oncnēow & recognized; perceived & third person, preterite, indicative, singular \\
geongan & young & genitive, singular, masculine \\
snelnesse & skill; agility & accusative, singular, feminine \\
ðæt & that; the precise linkage at 169.6 remains uncertain (P169N6) & as used here \\
plegan & play; game & dative, singular, masculine \\
tō & to, the person addressed in speech & invariable \\
heonan & hence; from here & invariable \\
cniht & young man; youth & nominative, singular, masculine \\
ðæs-ðe & as, in the exact parenthesis ðæs-ðe mē ðyncð & as used here \\
gelīca & equal; peer & nominative, singular, masculine \\
hęrede & praised & third singular preterite; indicative reading here, with the same subjunctive spelling possible \\
genēalǣcte & approached & third person, preterite, indicative, singular \\
swang & swung; brandished & third person, preterite, indicative, singular \\
top & ball?, as Cook tentatively interprets it & accusative, singular, masculine \\
swiftnesse & swiftness & dative, singular, feminine \\
wæs & in the impersonal source expression it seemed to the king & third person, preterite, indicative, singular \\
geðūht & seemed & past, participle \\
swilce & as if, introducing the source comparison of apparent rejuvenation & invariable \\
of & from, the state of old age in the supplied rejuvenation comparison & dative \\
ielde & age; old age & dative, singular, feminine \\
tō & to, a resulting state of youth & invariable \\
geoguðe & youth & dative, singular, feminine \\
gewęnd & turned; changed back & past, participle \\
cynesętle & royal seat; throne & dative, singular, neuter \\
gecwēmlīce & pleasantly; agreeably & invariable \\
ðēnode & served; attended & third person, preterite, indicative, singular \\
ūtēode & went out & third person, preterite, indicative, singular \\
be & by, the body part by which a person is led & as used here \\
ðanan & thence; from there & as used here \\
gewęnde & turned; changed back & third person, preterite, indicative, singular \\
āgān & gone away & past, participle \\
wæs & auxiliary with an intransitive past participle & third person, preterite, indicative, singular \\
swęrie & swear & first person, present, indicative, singular \\
gemǣnan & common; shared & accusative, singular, feminine \\
hǣle & safety; rescue & accusative, singular, feminine \\
baðode & bathed & first person, preterite, indicative, singular \\
nāt & do not know & first person, present, indicative, singular \\
ðurh & through; by means of service, or by evidence in the exact source question Ðurh hwæt & as used here \\
geonges & young & genitive, singular, masculine \\
ðēnunge & service; attendance & accusative, singular, feminine \\
beseah & looked; turned his gaze & third person, preterite, indicative, singular \\
gewite & find out; learn & second person, imperative, singular \\
geonga & young & nominative, singular, masculine \\
gehīersumode & served; obeyed & third person, preterite, indicative, singular \\
æfter & after; in search of; about the person sought or inquired after & as used here \\
Apollonio & Apollonius & dative, singular, masculine \\
geseah & see that a stated circumstance holds & third person, preterite, indicative, singular \\
horgum & dirty; filthy & dative, singular, masculine \\
bewǣfed & wrapped; clothed & past, participle \\
węnde & turned; returned & third person, preterite, indicative, singular \\
āscodest & asked about & second person, preterite, indicative, singular \\
forliden & shipwrecked & nominative, singular, masculine \\
forswīge & keep silent about; conceal & third person, present, subjunctive, singular \\
gesweotolað & reveals; makes clear & third person, present, indicative, singular \\
sęge & say; tell & second person, imperative, singular \\
gereorde & meal; banquet & dative, singular, neuter \\
\end{longtable}
\par\endgroup

% mdtable:T90:end
% mdparagraph:P539:begin
\textbf{Further verbal help.} The following notes identify the action or construction in the exact source form.
% mdparagraph:P539:end

% mdparagraph:P540:begin
\textbf{168.02. getǣht — shown; directed.} Past participle in passive getǣht wæs; him is the person to whom the route was indicated.
% mdparagraph:P540:end

% mdparagraph:P541:begin
\textbf{168.03. inēode — went in.} The prefixed form means went in; its finite verb is part of the word.
% mdparagraph:P541:end

% mdparagraph:P542:begin
\textbf{168.03. meahte — might; could.} Inside the dependent question hwone ... biddan meahte; indicative/\allowbreak{}subjunctive surface coincides. The selected interpretation is possibility in the thought.
% mdparagraph:P542:end

% mdparagraph:P543:begin
\textbf{168.05. gesmierwed — anointed.} Exact source prefixed participle and retained w; distinct from selected unprefixed smierwan paradigm. Mid ęle marks the anointing material.
% mdparagraph:P543:end

% mdparagraph:P544:begin
\textbf{168.06. begyrd — girded; clothed around.} Passive/\allowbreak{}predicative participle with mid scīetan, a cloth.
% mdparagraph:P544:end

% mdparagraph:P545:begin
\textbf{168.07. belimpende — belonging; relating.} Tō + dative bæðstęde identifies what the games/\allowbreak{}equipment relate to; preserve P168N6.
% mdparagraph:P545:end

% mdparagraph:P546:begin
\textbf{168.07. cleopode — cried aloud.} Preterite third singular of cleopian in the cry-aloud sense. Micelre stefne expresses the manner or means, “with a loud voice”.
% mdparagraph:P546:end

% mdparagraph:P547:begin
\textbf{168.08. Gehīere — hear; listen.} Exact plural command with shortened -e before following gē, not singular by surface alone. The same form recurs later in the line.
% mdparagraph:P547:end

% mdparagraph:P548:begin
\textbf{168.11. unscrȳdde — undressed; took off.} Hine is the person undressed; ðǣm healfan sciccelse is the removed clothing, supplied case frame.
% mdparagraph:P548:end

% mdparagraph:P549:begin
\textbf{168.13. behēold — observed.} Accusative persons observed; no modern behold quotation formula implied.
% mdparagraph:P549:end

% mdparagraph:P550:begin
\textbf{168.17. inēode — went in.} On ðæt bæð gives the entry destination.
% mdparagraph:P550:end

% mdparagraph:P551:begin
\textbf{168.17. āgan — began.} Strong III preterite of āginnan, not āgan own/\allowbreak{}possess and not āgān depart.
% mdparagraph:P551:end

% mdparagraph:P552:begin
\textbf{168.17. plegian — play.} Infinitive after began; wið + dative companions and mid + dative ball are supported locally.
% mdparagraph:P552:end

% mdparagraph:P553:begin
\textbf{168.18. gemęngde — joined; mingled.} Reflexive hine; on ðæs cyninges plegan continues on 169.1. Plegan is an oblique noun governed by on; ðæs cyninges is its genitive possessor. It is not a finite verb.
% mdparagraph:P553:end

% mdparagraph:P554:begin
\textbf{169.01. iernende — running.} The supplied present participle accompanies Apollonius. The printed division iern-/\allowbreak{}ende is joined in the reading.
% mdparagraph:P554:end

% mdparagraph:P555:begin
\textbf{169.02. gelæhte — caught.} Exact source preterite has short æ as printed; preserve lexical spelling and do not silently import a guessed long form.
% mdparagraph:P555:end

% mdparagraph:P556:begin
\textbf{169.03. geslægene — struck.} The participial phrase refers to the accusative ðone ðoðor. Retain the unusual printed -e, to which Cook draws attention in P169N2.
% mdparagraph:P556:end

% mdparagraph:P557:begin
\textbf{169.03. gesęnde — sent back.} Tō + dative recipient; ongēan is the return direction.
% mdparagraph:P557:end

% mdparagraph:P558:begin
\textbf{169.03. plegiendan — playing.} Weak dative modifier of cyninge after ðǣm; the finite paradigm does not itself teach this participial cell.
% mdparagraph:P558:end

% mdparagraph:P559:begin
\textbf{169.04. slōh — struck.} The supplied h spelling is a valid singular variant of the previously selected slōg.
% mdparagraph:P559:end

% mdparagraph:P560:begin
\textbf{169.05. oncnēow — recognized; perceived.} Accusative snelnesse, followed by a result/\allowbreak{}content sequence whose pronoun referents remain genuinely underdetermined in P169N6.
% mdparagraph:P560:end

% mdparagraph:P561:begin
\textbf{169.09. hęrede — praised.} Hine is the person praised. The selected indicative interpretation fits the known event, but the third-singular preterite subjunctive has the same spelling; the ending alone does not distinguish them.
% mdparagraph:P561:end

% mdparagraph:P562:begin
\textbf{169.10. arn — ran.} Third singular preterite of iernan, as in the learned table.
% mdparagraph:P562:end

% mdparagraph:P563:begin
\textbf{169.10. genēalǣcte — approached.} The prefixed verb with tō + dative goal is supplied; do not rename it as the unprefixed nēalǣcan form.
% mdparagraph:P563:end

% mdparagraph:P564:begin
\textbf{169.11. swang — swung; brandished.} Its direct object is top, tentatively “ball”. The Latin parallel describes rubbing or anointing with ointment; retain the actual Old English action rather than replacing it by that parallel.
% mdparagraph:P564:end

% mdparagraph:P565:begin
\textbf{169.12. geðūht — seemed.} Supply ðǣm cyninge wæs geðūht swilce ... wǣre, it seemed to the king as though ...; dative experiencer and irrealis comparison are explicit.
% mdparagraph:P565:end

% mdparagraph:P566:begin
\textbf{169.13. gewęnd — turned; changed back.} In swilce hē of ielde tō geoguðe gewęnd wǣre: as though he had been turned from age to youth. Do not mechanically label every wesan+participle a motion perfect.
% mdparagraph:P566:end

% mdparagraph:P567:begin
\textbf{169.14. ðēnode — served; attended.} The dative him is the person served. Retain the source’s -ode spelling.
% mdparagraph:P567:end

% mdparagraph:P568:begin
\textbf{169.15. ūtēode — went out.} Prefixed motion verb with of + dative place left.
% mdparagraph:P568:end

% mdparagraph:P569:begin
\textbf{169.16. gewęnde — turned; went away.} Reflexive dative him; ðæs weges is the route genitive, matched by the composition 35.3. Pronoun references need discourse reasoning.
% mdparagraph:P569:end

% mdparagraph:P570:begin
\textbf{170.02. āgān — gone away.} Past participle in āgān wæs, “had gone away”. Its second vowel is long; the second vowel of āgan “began” at 168.17 is short.
% mdparagraph:P570:end

% mdparagraph:P571:begin
\textbf{170.02. swęrie — swear.} First-person quoted oath; ðurh + accusative oath standard; the narrated speaker now owns ic.
% mdparagraph:P571:end

% mdparagraph:P572:begin
\textbf{170.03. baðode — bathed.} Reflexive mē and adverbial comparison with dyde; this bath verb is supplied, substitute dōn is newly productive.
% mdparagraph:P572:end

% mdparagraph:P573:begin
\textbf{170.03. nāt — do not know.} The supplied contracted negative means “I do not know”. Only this local form of nytan is required.
% mdparagraph:P573:end

% mdparagraph:P574:begin
\textbf{170.04. beseah — looked; turned his gaze.} Reflexive hine and tō + dative person toward whom he turned; supply as a local phrase.
% mdparagraph:P574:end

% mdparagraph:P575:begin
\textbf{170.05. gewite — find out; learn.} The preterite-present knowing verb (§126), not strong I gewītan depart; the dependent question has hwæt ... sīe.
% mdparagraph:P575:end

% mdparagraph:P576:begin
\textbf{170.06. gehīersumode — served; obeyed.} Dative mē is the person served; hyphenation gehīer-/\allowbreak{}sumode is print line division only.
% mdparagraph:P576:end

% mdparagraph:P577:begin
\textbf{170.09. bewǣfed — wrapped; clothed.} Mid + dative garment in the passive/\allowbreak{}predicative description.
% mdparagraph:P577:end

% mdparagraph:P578:begin
\textbf{170.09. węnde — turned; returned.} Ongēan + tō destination; distinct from wēnan expected despite the similar modern gloss.
% mdparagraph:P578:end

% mdparagraph:P579:begin
\textbf{170.11. āscodest — asked about.} Æfter links to the relative antecedent; exact source spelling preserves -odest.
% mdparagraph:P579:end

% mdparagraph:P580:begin
\textbf{170.13. forswīge — keep silent about; conceal.} Ðēah concessive clause; hit is the thing withheld. Exact source ending is supplied rather than silently generalized to a weak-II table.
% mdparagraph:P580:end

% mdparagraph:P581:begin
\textbf{170.14. gesweotolað — reveals; makes clear.} Gegierela is the singular subject, hine the person whose condition is revealed.
% mdparagraph:P581:end

% mdparagraph:P582:begin
\textbf{170.15. sęge — say; tell.} Exact source shortened command spelling; the following mixed direct/\allowbreak{}indirect quotation is explicitly supported by P170N15, not presented as a regular report transformation.
% mdparagraph:P582:end

% mdparagraph:P583:begin
\textbf{168.01, 168.02: ðone weg ðe him getǣht wæs.} Relative ðe refers to the road; getǣht wæs is a past passive, with him the person shown the route.
% mdparagraph:P583:end

% mdparagraph:P584:begin
\textbf{168.02: ōð-ðæt hē becōm.} Until he arrived: an endpoint clause; retain the source macron.
% mdparagraph:P584:end

% mdparagraph:P585:begin
\textbf{168.03, 168.04: hwone hē biddan meahte līfes fultumes.} hwone is the person asked, while līfes fultumes names the requested support of life. meahte can be indicative or subjunctive in form; context is a question within his thought, not proof of a uniquely recognizable ending.
% mdparagraph:P585:end

% mdparagraph:P586:begin
\textbf{168.04, 168.05: geseah hē ǣnne nacodne cnapan ... iernan.} Accusative participant with a perception infinitive. Cook’s glossary allows nacod to mean wearing only a tunic here; do not require a claim of total nudity.
% mdparagraph:P586:end

% mdparagraph:P587:begin
\textbf{168.06, 168.07: geongra manna plegan ... tō ðǣm bæðstęde belimpende.} plegan literally denotes play/\allowbreak{}games; carrying equipment for the young men’s games is a contextual explanation, not a new unrestricted dictionary sense. belimpende modifies plegan, as P168N6 says; it is not another finite action of the boy.
% mdparagraph:P587:end

% mdparagraph:P588:begin
\textbf{168.08, 168.09: Gehīere, gē ceasterwaran ....} The call is addressed to the plural audience. Supply the whole source exhortation; do not replace its Gehīere with a silently regularized imperative.
% mdparagraph:P588:end

% mdparagraph:P589:begin
\textbf{168.11, 168.12: hine unscrȳdde ðǣm healfan sciccelse ðe hē on hæfde.} Reflexive undressing and the garment in the dative; the relative clause says he had the cloak on. Preserve the source construction as a whole.
% mdparagraph:P589:end

% mdparagraph:P590:begin
\textbf{168.13: hiera ānra gehwilcne on hiera weorce.} Each single one of them at their work: first hiera participates in the partitive expression; second is possessive. Keep the printed short spelling.
% mdparagraph:P590:end

% mdparagraph:P591:begin
\textbf{169.02, 169.03: mid swiftre hrædnesse geslægene.} P169N2 remarks on the participle’s awkward form. Supply the complete phrase as describing the ball struck with swift speed; retain geslægene rather than silently repairing agreement.
% mdparagraph:P591:end

% mdparagraph:P592:begin
\textbf{169.05, 169.06, 169.07: oncnēow ... ðæt hē wiste ðæt hē næfde his gelīcan.} The nested reference is unclear, as P169N6 states. næfde is negative: the text does not say he had an equal. The Latin comparison concerns the king’s previously unmatched skill; use it as help while retaining the OE uncertainty, not as an invisible replacement text.
% mdparagraph:P592:end

% mdparagraph:P593:begin
\textbf{169.08: ðæs-ðe mē ðyncð.} As it seems to me: exact parenthetic impersonal construction, dative experiencer mē.
% mdparagraph:P593:end

% mdparagraph:P594:begin
\textbf{169.11, 169.12, 169.13: swang ... ðone top ... swilce ... gewęnd wǣre.} Retain the Old English action with top, tentatively a ball. Cook’s Latin parallel describes rubbing or anointing with ointment and its rejuvenating effect; it names no shoulder or neck. Wǣre belongs to the imagined transformation comparison.
% mdparagraph:P594:end

% mdparagraph:P595:begin
\textbf{169.12: ðǣm cyninge wæs geðūht.} It seemed to the king: whole impersonal be-plus-participle construction supplied; do not read it as an ordinary king-as-subject clause.
% mdparagraph:P595:end

% mdparagraph:P596:begin
\textbf{169.14, 169.15, 169.16: him ... hē hine lǣdde be ðǣre handa ... him ... gewęnde.} Use P169N11 and P169N14–P169N16 for the pronoun chain. The later him is reflexive; the by-the-hand phrase specifies how one person leads another.
% mdparagraph:P596:end

% mdparagraph:P597:begin
\textbf{169.16, 169.17: ðæs weges ðe hē ǣr cōm.} The genitive names the route followed, as in the pattern taught above; the relative ðe identifies that route. Retain the past-time reference of the source wording.
% mdparagraph:P597:end

% mdparagraph:P598:begin
\textbf{170.01, 170.02: siððan Apollonius āgān wæs.} After Apollonius had gone away: intransitive completed departure with be, not a passive in which someone sends him away.
% mdparagraph:P598:end

% mdparagraph:P599:begin
\textbf{170.02: ðurh ðā gemǣnan hǣle.} An oath by the common safety or welfare. Hǣle is feminine accusative singular, here in the safety/\allowbreak{}rescue sense. No possessive “our” is printed.
% mdparagraph:P599:end

% mdparagraph:P600:begin
\textbf{170.04, 170.05, 170.06: ðurh hwilces geonges mannes ðēnunge ... hwæt se geonga mann sīe.} The service is the means of the benefit; hwilces ... mannes is its genitive dependent. hwæt asks what sort/\allowbreak{}status of man he is. sīe is the dependent interrogative be form.
% mdparagraph:P600:end

% mdparagraph:P601:begin
\textbf{170.08, 170.09, 170.10, 170.11: ēode æfter Apollonio ... ðe ðū æfter āscodest.} Going after a person, then asking about him: the relative marker is the complement of æfter within the clause; do not force temporal after for either use.
% mdparagraph:P601:end

% mdparagraph:P602:begin
\textbf{170.13: Ðēah hē hit self forswīge.} Concessive clause: even if he himself keeps it unspoken. Supply the source forswīge with its present subjunctive interpretation; do not demand a full new weak-verb paradigm.
% mdparagraph:P602:end

% mdparagraph:P603:begin
\textbf{170.14, 170.15, 170.16: sęge him ðæt ‘se cyning bitt ðē ðæt ðū cume ...’.} P170N15 explicitly diagnoses mixed direct/\allowbreak{}reported construction. First him is the messenger’s addressee; within the quoted message ðē is the person invited (accusative), and ðū/\allowbreak{}cume address him directly. Retain wording and speech levels; no silent repair.
% mdparagraph:P603:end

% mdparagraph:P604:begin
\textbf{The Latin parallels in English.} Original P169N6 compares the difficult knowledge/\allowbreak{}equal clause with a Latin account in which the king has long known his own skill to be without an equal. This parallel helps explain Cook’s concern but does not settle every Old English pronoun or cancel \textbf{næfde}. Original P169N10 quotes a Latin account of rubbing or anointing with an ointment used by wrestlers; it does not name a shoulder or neck. Keep the Old English \textbf{swang ðone top} and the note’s tentative interpretation distinct from that Latin action. No Latin translation is required.
% mdparagraph:P604:end

% mdparagraph:P605:begin
\textbf{Quotation punctuation.} The edition begins a quotation at 170.2 but does not print its expected closing mark before the next narrated action. The final invitation also mixes quoted and reported wording. The text below retains those signs and constructions; they are not silently regularized.
% mdparagraph:P605:end

\reserve{7\baselineskip}
% mdparagraph:P606:begin
\subsection*{Source text}
% mdparagraph:P606:end

% passage: FB-XII-APOLLONIUS-GYMNASIUM
% mdparagraph:P607:begin
\readingtext{Æfter ðisum wordum hē ēode on ðone weg ðe him getǣht\textsuperscript{P168N1} wæs, ōð-ðæt hē becōm tō ðǣre ceastre geate, and ðǣr inēode. Mid-ðȳ-ðe hē ðōhte hwone hē biddan meahte līfes fultumes,\textsuperscript{P168N2} ðā geseah hē ǣnne nacodne cnapan geond ðā strǣte iernan. Sē wæs mid\textsuperscript{P168N3} ęle gesmierwed, and mid scīetan begyrd, and bær geongra\textsuperscript{P168N4} manna\textsuperscript{P168N4} plegan\textsuperscript{P168N4} on handa, tō ðǣm bæðstęde\textsuperscript{P168N5} belimpende.\textsuperscript{P168N6} And cleopode\textsuperscript{P168N7} micelre stefne,\textsuperscript{P168N8} and cwæð: “Gehīere,\textsuperscript{P168N9} gē ceasterwaran\textsuperscript{P168N10}! Gehīere, gē ęlðēodige,\textsuperscript{P168N11} frige and ðēowe, æðele and unæðele! Se bæðstęde is open.”}
% mdparagraph:P607:end

% mdparagraph:P608:begin
\readingtext{Ðā-ðā Apollonius ðæt gehīerde, hē hine unscrȳdde ðǣm\textsuperscript{P168N12} healfan sciccelse ðe hē on hæfde, and ēode intō ðǣm ðwēale.\textsuperscript{P168N13} And mid-ðȳ-ðe hē behēold hiera ānra\textsuperscript{P168N14} gehwilcne on hiera weorce, hē sōhte his gelīcan,\textsuperscript{P168N15} ac hē ne meahte hine ðǣr findan on ðǣm flocce. Ðā fǣringa cōm Arcestrates, ealre ðǣre ðēode\textsuperscript{P168N16} cyning,\textsuperscript{P168N17} mid micelre męnige his manna,\textsuperscript{P168N18} and inēode on ðæt bæð. Ðā āgan se cyning plegian wið\textsuperscript{P168N19} his gefērum mid ðoðore.\textsuperscript{P168N20} And Apollonius hine\textsuperscript{P168N21} gemęngde,\textsuperscript{P168N21} swā-swā God wolde, on ðæs cyninges plegan, and, iernende, ðone ðoðor gelæhte,\textsuperscript{P169N1} and, mid swiftre hrædnesse geslægene,\textsuperscript{P169N2} ongēan gesęnde tō ðǣm plegiendan cyninge. Ęft hē āgēan āsęnde; hē hrædlīce slōh, swā hē hine\textsuperscript{P169N3} nǣfre feallan ne lēt. Se cyning ðā oncnēow ðæs geongan snelnesse,\textsuperscript{P169N4} ðæt hē wiste\textsuperscript{P169N5} ðæt hē næfde his gelīcan on ðǣm plegan.\textsuperscript{P169N6} Ðā cwæð hē tō his gefērum: “Gāð ēow heonan; ðes cniht, ðæs-ðe mē ðyncð,\textsuperscript{P169N7} is mīn gelīca.”}
% mdparagraph:P608:end

% mdparagraph:P609:begin
\readingtext{Ðā-ðā Apollonius gehīerde ðæt se cyning hine hęrede, hē arn hrædlīce, and genēalǣcte tō ðǣm cyninge, and mid gelǣredre\textsuperscript{P169N8} handa\textsuperscript{P169N9} hē swang\textsuperscript{P169N10} ðone top mid swā micelre swiftnesse ðæt ðǣm cyninge wæs geðūht swilce hē of ielde tō geoguðe gewęnd wǣre. And, æfter ðǣm, on his cynesętle hē him\textsuperscript{P169N11} gecwēmlīce\textsuperscript{P169N12} ðēnode\textsuperscript{P169N13}; and, ðā-ðā hē ūtēode of ðǣm bæðe, hē\textsuperscript{P169N14} hine\textsuperscript{P169N15} lǣdde be ðǣre handa, and him\textsuperscript{P169N16} ðā siððan ðanan gewęnde, ðæs weges\textsuperscript{P169N17} ðe hē ǣr cōm.}
% mdparagraph:P609:end

% mdparagraph:P610:begin
\readingtext{Ðā cwæð se cyning tō his mannum,\textsuperscript{P170N1} siððan Apollonius āgān\textsuperscript{P170N2} wæs: “Ic swęrie ðurh ðā gemǣnan hǣle\textsuperscript{P170N3} ðæt ic mē nǣfre bęt ne baðode ðonne ic dyde\textsuperscript{P170N4} tō-dæg, nāt ic ðurh\textsuperscript{P170N5} hwilces geonges mannes ðēnunge.\textsuperscript{P170N6} Ðā beseah hē hine tō ānum his manna, and cwæð: “Gā, and gewite hwæt\textsuperscript{P170N7} se geonga mann sīe, ðe mē tō-dæg swā wel gehīersumode.”}
% mdparagraph:P610:end

% mdparagraph:P611:begin
\readingtext{Se mann ðā ēode æfter Apollonio. Mid-ðȳ-ðe hē geseah ðæt hē\textsuperscript{P170N8} wæs mid horgum\textsuperscript{P170N9} sciccelse bewǣfed, ðā węnde hē ongēan tō ðǣm cyninge, and cwæð: “Se geonga mann ðe\textsuperscript{P170N10} ðū æfter āscodest is forliden\textsuperscript{P170N11} mann.\textsuperscript{P170N11}” Ðā cwæð se cyning: “Ðurh\textsuperscript{P170N12} hwæt\textsuperscript{P170N12} wāst\textsuperscript{P170N13} ðū ðæt?” Se mann him andswarode, and cwæð: “Ðēah hē hit self forswīge,\textsuperscript{P170N14} his gegierela hine gesweotolað.” Ðā cwæð se cyning: “Gā hrædlīce, and sęge him ðæt\textsuperscript{P170N15} ‘se cyning bitt ðē ðæt ðū cume\textsuperscript{P170N16} tō his gereorde.’”}
% mdparagraph:P611:end

\par
\reserve{7\baselineskip}
% mdparagraph:P612:begin
\subsection*{Original notes}
% mdparagraph:P612:end

% mdparagraph:P613:begin
\textsuperscript{P168N1}: See 187.
% mdparagraph:P613:end

% mdparagraph:P614:begin
\textsuperscript{P168N2}: MS. fultum. See 159. b.
% mdparagraph:P614:end

% mdparagraph:P615:begin
\textsuperscript{P168N3}: See 174.
% mdparagraph:P615:end

% mdparagraph:P616:begin
\textsuperscript{P168N4}: Lat. lusus juvenales.
% mdparagraph:P616:end

% mdparagraph:P617:begin
\textsuperscript{P168N5}: Lat. gymnasium.
% mdparagraph:P617:end

% mdparagraph:P618:begin
\textsuperscript{P168N6}: Modifies plegan.
% mdparagraph:P618:end

% mdparagraph:P619:begin
\textsuperscript{P168N7}: See 20.
% mdparagraph:P619:end

% mdparagraph:P620:begin
\textsuperscript{P168N8}: See 160. 1. This word is the Chaucerian steven.
% mdparagraph:P620:end

% mdparagraph:P621:begin
\textsuperscript{P168N9}: See 95, note.
% mdparagraph:P621:end

% mdparagraph:P622:begin
\textsuperscript{P168N10}: This is a compound word, formed of a Latin and an English element. Which is Latin, and what is its form in that language?
% mdparagraph:P622:end

% mdparagraph:P623:begin
\textsuperscript{P168N11}: Lat. peregrini. See 152.
% mdparagraph:P623:end

% mdparagraph:P624:begin
\textsuperscript{P168N12}: See 162.
% mdparagraph:P624:end

% mdparagraph:P625:begin
\textsuperscript{P168N13}: Lat. lavacrum.
% mdparagraph:P625:end

% mdparagraph:P626:begin
\textsuperscript{P168N14}: See 154. b.
% mdparagraph:P626:end

% mdparagraph:P627:begin
\textsuperscript{P168N15}: Lat. parem, Eng. peer.
% mdparagraph:P627:end

% mdparagraph:P628:begin
\textsuperscript{P168N16}: Lat. regionis.
% mdparagraph:P628:end

% mdparagraph:P629:begin
\textsuperscript{P168N17}: See 151.
% mdparagraph:P629:end

% mdparagraph:P630:begin
\textsuperscript{P168N18}: Lat. famularum. See 154. a.
% mdparagraph:P630:end

% mdparagraph:P631:begin
\textsuperscript{P168N19}: Why not mid?
% mdparagraph:P631:end

% mdparagraph:P632:begin
\textsuperscript{P168N20}: This curious word is very rare in Old English.
% mdparagraph:P632:end

% mdparagraph:P633:begin
\textsuperscript{P168N21}: Lat. miscuit se. See 184. b.
% mdparagraph:P633:end

% mdparagraph:P634:begin
\textsuperscript{P169N1}: See 114.
% mdparagraph:P634:end

% mdparagraph:P635:begin
\textsuperscript{P169N2}: Lat. subtili velocitate percussam. The OE. participle is a little awkward.
% mdparagraph:P635:end

% mdparagraph:P636:begin
\textsuperscript{P169N3}: The ball.
% mdparagraph:P636:end

% mdparagraph:P637:begin
\textsuperscript{P169N4}: Lat. velocitatem.
% mdparagraph:P637:end

% mdparagraph:P638:begin
\textsuperscript{P169N5}: See 126. What is the latest English quotation that you can find for this word?
% mdparagraph:P638:end

% mdparagraph:P639:begin
\textsuperscript{P169N6}: This clause is not very clear. The Latin has: et quia sciebat se (i.e. Archistrates) in pilæ lusu neminem parem habere, ad suos ait, famuli, recedite; hic enim juvenis, etc.
% mdparagraph:P639:end

% mdparagraph:P640:begin
\textsuperscript{P169N7}: Lat. ut suspicor. See 157. 1 and 164. l. What Mod. Eng. word comes from mē ðyncð?
% mdparagraph:P640:end

% mdparagraph:P641:begin
\textsuperscript{P169N8}: Lat. docta.
% mdparagraph:P641:end

% mdparagraph:P642:begin
\textsuperscript{P169N9}: See 51. 3.
% mdparagraph:P642:end

% mdparagraph:P643:begin
\textsuperscript{P169N10}: Here the English departs from the Latin: ceroma fricavit eum tanta subtilitate, ut de sene juvenem redderet. Top would seem to signify the same as ðoðor.
% mdparagraph:P643:end

% mdparagraph:P644:begin
\textsuperscript{P169N11}: See 164. e.
% mdparagraph:P644:end

% mdparagraph:P645:begin
\textsuperscript{P169N12}: Lat. gratissime.
% mdparagraph:P645:end

% mdparagraph:P646:begin
\textsuperscript{P169N13}: See 28.
% mdparagraph:P646:end

% mdparagraph:P647:begin
\textsuperscript{P169N14}: Apollonius.
% mdparagraph:P647:end

% mdparagraph:P648:begin
\textsuperscript{P169N15}: Archistrates.
% mdparagraph:P648:end

% mdparagraph:P649:begin
\textsuperscript{P169N16}: See 184. a.
% mdparagraph:P649:end

% mdparagraph:P650:begin
\textsuperscript{P169N17}: See 157.
% mdparagraph:P650:end

% mdparagraph:P651:begin
\textsuperscript{P170N1}: Lat. amicos.
% mdparagraph:P651:end

% mdparagraph:P652:begin
\textsuperscript{P170N2}: How is the sense of Mod. Eng. ago related to that of this word?
% mdparagraph:P652:end

% mdparagraph:P653:begin
\textsuperscript{P170N3}: This phrase shows Christian influence.
% mdparagraph:P653:end

% mdparagraph:P654:begin
\textsuperscript{P170N4}: Note this use of dōn to replace a verb of specific meaning.
% mdparagraph:P654:end

% mdparagraph:P655:begin
\textsuperscript{P170N5}: Governs ðēnunge.
% mdparagraph:P655:end

% mdparagraph:P656:begin
\textsuperscript{P170N6}: See 28.
% mdparagraph:P656:end

% mdparagraph:P657:begin
\textsuperscript{P170N7}: How does this, as here used, differ in meaning from hwā?
% mdparagraph:P657:end

% mdparagraph:P658:begin
\textsuperscript{P170N8}: Apollonius.
% mdparagraph:P658:end

% mdparagraph:P659:begin
\textsuperscript{P170N9}: Lat. sordido.
% mdparagraph:P659:end

% mdparagraph:P660:begin
\textsuperscript{P170N10}: Governed by æfter. See 87. c and 201. 1.
% mdparagraph:P660:end

% mdparagraph:P661:begin
\textsuperscript{P170N11}: Lat. naufragus.
% mdparagraph:P661:end

% mdparagraph:P662:begin
\textsuperscript{P170N12}: Lat. unde.
% mdparagraph:P662:end

% mdparagraph:P663:begin
\textsuperscript{P170N13}: See 126.
% mdparagraph:P663:end

% mdparagraph:P664:begin
\textsuperscript{P170N14}: See 196. e.
% mdparagraph:P664:end

% mdparagraph:P665:begin
\textsuperscript{P170N15}: Confusion of two constructions, the direct and the indirect.
% mdparagraph:P665:end

% mdparagraph:P666:begin
\textsuperscript{P170N16}: Lat. ut venias. Translate by the infinitive, as often in such cases.
% mdparagraph:P666:end

\reserve{7\baselineskip}
% mdparagraph:P667:begin
\subsection*{Reading tasks}
% mdparagraph:P667:end

% mdparagraph:P668:begin
Use the whole text, the original notes and the supplied modern help for all six tasks. A supported interpretation should preserve the actual negative and clause structure while acknowledging a genuine uncertainty. The line excerpts below are locating aids; interpret the complete clause in the continuous reading.
% mdparagraph:P668:end

\reserve{5\baselineskip}
\par\begin{minipage}{\linewidth}\raggedright
% mdparagraph:P669:begin
\textbf{R01}
% mdparagraph:P669:end

% item: NEW-U16-R01
% mdparagraph:P670:begin
Translate the whole gymnasium scene, from Æfter ðisum wordum through tō his gereorde. Keep the local help and original notes available. Preserve the changes of speaker, the complete action sequence and the places where the edition acknowledges uncertainty.
% mdparagraph:P670:end

% output: NEW-U16-R01-O001
% mdparagraph:P671:begin
(1) Whole translation: \_\_\_\_\_\_\_\_\_\_\_\_\_\_\_\_\_\_\_\_
% mdparagraph:P671:end

\end{minipage}\par\addvspace{4pt}

\par\begin{minipage}{\linewidth}\raggedright
% mdparagraph:P672:begin
\textbf{R02}
% mdparagraph:P672:end

% item: NEW-U16-R02
% mdparagraph:P673:begin
Explain the sequence and grounds of the encounter in five brief answers, citing the relevant words or source lines: what brings Apollonius into the bath; what attracts the king's notice during the game; what the king says about the young man; what the messenger infers and from what evidence; and what the king finally orders.
% mdparagraph:P673:end

% output: NEW-U16-R02-O001
% mdparagraph:P674:begin
(1) Entry into the bath: \_\_\_\_\_\_\_\_\_\_\_\_\_\_\_\_\_\_\_\_
% mdparagraph:P674:end

% output: NEW-U16-R02-O002
% mdparagraph:P675:begin
(2) The game: \_\_\_\_\_\_\_\_\_\_\_\_\_\_\_\_\_\_\_\_
% mdparagraph:P675:end

% output: NEW-U16-R02-O003
% mdparagraph:P676:begin
(3) The king's statement: \_\_\_\_\_\_\_\_\_\_\_\_\_\_\_\_\_\_\_\_
% mdparagraph:P676:end

% output: NEW-U16-R02-O004
% mdparagraph:P677:begin
(4) The messenger's conclusion: \_\_\_\_\_\_\_\_\_\_\_\_\_\_\_\_\_\_\_\_
% mdparagraph:P677:end

% output: NEW-U16-R02-O005
% mdparagraph:P678:begin
(5) Final order: \_\_\_\_\_\_\_\_\_\_\_\_\_\_\_\_\_\_\_\_
% mdparagraph:P678:end

\end{minipage}\par\addvspace{4pt}

\par\begin{minipage}{\linewidth}\raggedright
% mdparagraph:P679:begin
\textbf{R03}
% mdparagraph:P679:end

% item: NEW-U16-R03
% mdparagraph:P680:begin
Identify the referent and grammatical role of each selected pronoun. Use the local context and original notes; distinguish a contextual interpretation from what its ending alone proves.
% mdparagraph:P680:end

% mdparagraph:P681:begin
Selected passages: 169.4 hine after slōh; 169.14 him after hē; 169.15 hē immediately before hine lǣdde; 169.15 hine before lǣdde; 169.16 him before ðā siððan; 170.9 hē after ðæt.
% mdparagraph:P681:end

% mdparagraph:P682:begin
Source wording for the cited locations:
% mdparagraph:P682:end

% mdparagraph:P683:begin
169.04: Ęft hē āgēan āsęnde; hē hrædlīce slōh, swā hē hine
% mdparagraph:P683:end

% mdparagraph:P684:begin
169.14: on his cynesętle hē him gecwēmlīce ðēnode; and,
% mdparagraph:P684:end

% mdparagraph:P685:begin
169.15: ðā-ðā hē ūtēode of ðǣm bæðe, hē hine lǣdde be ðǣre
% mdparagraph:P685:end

% mdparagraph:P686:begin
169.16: handa, and him ðā siððan ðanan gewęnde, ðæs weges
% mdparagraph:P686:end

% mdparagraph:P687:begin
170.09: ðæt hē wæs mid horgum sciccelse bewǣfed, ðā węnde
% mdparagraph:P687:end

% output: NEW-U16-R03-O001
% mdparagraph:P688:begin
(1) 169.4 hine: \_\_\_\_\_\_\_\_\_\_\_\_\_\_\_\_\_\_\_\_
% mdparagraph:P688:end

% output: NEW-U16-R03-O002
% mdparagraph:P689:begin
(2) 169.14 him: \_\_\_\_\_\_\_\_\_\_\_\_\_\_\_\_\_\_\_\_
% mdparagraph:P689:end

% output: NEW-U16-R03-O003
% mdparagraph:P690:begin
(3) 169.15 hē before hine: \_\_\_\_\_\_\_\_\_\_\_\_\_\_\_\_\_\_\_\_
% mdparagraph:P690:end

% output: NEW-U16-R03-O004
% mdparagraph:P691:begin
(4) 169.15 hine: \_\_\_\_\_\_\_\_\_\_\_\_\_\_\_\_\_\_\_\_
% mdparagraph:P691:end

% output: NEW-U16-R03-O005
% mdparagraph:P692:begin
(5) 169.16 him: \_\_\_\_\_\_\_\_\_\_\_\_\_\_\_\_\_\_\_\_
% mdparagraph:P692:end

% output: NEW-U16-R03-O006
% mdparagraph:P693:begin
(6) 170.9 hē: \_\_\_\_\_\_\_\_\_\_\_\_\_\_\_\_\_\_\_\_
% mdparagraph:P693:end

\end{minipage}\par\addvspace{4pt}

\par\begin{minipage}{\linewidth}\raggedright
% mdparagraph:P694:begin
\textbf{R04}
% mdparagraph:P694:end

% item: NEW-U16-R04
% mdparagraph:P695:begin
For each passage give its relation to the surrounding clause, identify its finite verb or finite verb group, and state the form's mood with the help provided. Treat the source-specific constructions within their stated limits.
% mdparagraph:P695:end

% mdparagraph:P696:begin
Selected passages: 168.1–2 ðe him getǣht wæs; 168.2 ōð-ðæt hē becōm; 168.3 hwone hē biddan meahte; 169.9 ðæt se cyning hine hęrede; 170.6 hwæt se geonga mann sīe; 170.13 Ðēah hē hit self forswīge.
% mdparagraph:P696:end

% output: NEW-U16-R04-O001
% mdparagraph:P697:begin
(1) 168.1–2 relation: \_\_\_\_\_\_\_\_\_\_\_\_\_\_\_\_\_\_\_\_
% mdparagraph:P697:end

% output: NEW-U16-R04-O002
% mdparagraph:P698:begin
(2) 168.1–2 finite group: \_\_\_\_\_\_\_\_\_\_\_\_\_\_\_\_\_\_\_\_
% mdparagraph:P698:end

% output: NEW-U16-R04-O003
% mdparagraph:P699:begin
(3) 168.2 relation: \_\_\_\_\_\_\_\_\_\_\_\_\_\_\_\_\_\_\_\_
% mdparagraph:P699:end

% output: NEW-U16-R04-O004
% mdparagraph:P700:begin
(4) 168.2 finite form: \_\_\_\_\_\_\_\_\_\_\_\_\_\_\_\_\_\_\_\_
% mdparagraph:P700:end

% output: NEW-U16-R04-O005
% mdparagraph:P701:begin
(5) 168.3 relation: \_\_\_\_\_\_\_\_\_\_\_\_\_\_\_\_\_\_\_\_
% mdparagraph:P701:end

% output: NEW-U16-R04-O006
% mdparagraph:P702:begin
(6) 168.3 finite group: \_\_\_\_\_\_\_\_\_\_\_\_\_\_\_\_\_\_\_\_
% mdparagraph:P702:end

% output: NEW-U16-R04-O007
% mdparagraph:P703:begin
(7) 169.9 relation: \_\_\_\_\_\_\_\_\_\_\_\_\_\_\_\_\_\_\_\_
% mdparagraph:P703:end

% output: NEW-U16-R04-O008
% mdparagraph:P704:begin
(8) 169.9 finite form: \_\_\_\_\_\_\_\_\_\_\_\_\_\_\_\_\_\_\_\_
% mdparagraph:P704:end

% output: NEW-U16-R04-O009
% mdparagraph:P705:begin
(9) 170.6 relation: \_\_\_\_\_\_\_\_\_\_\_\_\_\_\_\_\_\_\_\_
% mdparagraph:P705:end

% output: NEW-U16-R04-O010
% mdparagraph:P706:begin
(10) 170.6 finite form: \_\_\_\_\_\_\_\_\_\_\_\_\_\_\_\_\_\_\_\_
% mdparagraph:P706:end

% output: NEW-U16-R04-O011
% mdparagraph:P707:begin
(11) 170.13 relation: \_\_\_\_\_\_\_\_\_\_\_\_\_\_\_\_\_\_\_\_
% mdparagraph:P707:end

% output: NEW-U16-R04-O012
% mdparagraph:P708:begin
(12) 170.13 finite form: \_\_\_\_\_\_\_\_\_\_\_\_\_\_\_\_\_\_\_\_
% mdparagraph:P708:end

\end{minipage}\par\addvspace{4pt}

\par\begin{minipage}{\linewidth}\raggedright
% mdparagraph:P709:begin
\textbf{R05}
% mdparagraph:P709:end

% item: NEW-U16-R05
% mdparagraph:P710:begin
Analyse the final invitation at 170.14–16. Identify the outer speaker, his immediate addressee, the person denoted by him, the person denoted by se cyning, and the person denoted by ðē and ðū. Then explain why the original note describes a mixture of direct and indirect construction. Preserve the actual wording while explaining it.
% mdparagraph:P710:end

% mdparagraph:P711:begin
Selected passages: 170.14–16.
% mdparagraph:P711:end

% mdparagraph:P712:begin
Source wording for the cited locations:
% mdparagraph:P712:end

% mdparagraph:P713:begin
170.14: gegierela hine gesweotolað.” Ðā cwæð se cyning: “Gā
% mdparagraph:P713:end

% mdparagraph:P714:begin
170.15: hrædlīce, and sęge him ðæt ‘se cyning bitt ðē ðæt ðū
% mdparagraph:P714:end

% mdparagraph:P715:begin
170.16: cume tō his gereorde.’”
% mdparagraph:P715:end

% output: NEW-U16-R05-O001
% mdparagraph:P716:begin
(1) Outer speaker: \_\_\_\_\_\_\_\_\_\_\_\_\_\_\_\_\_\_\_\_
% mdparagraph:P716:end

% output: NEW-U16-R05-O002
% mdparagraph:P717:begin
(2) Immediate addressee: \_\_\_\_\_\_\_\_\_\_\_\_\_\_\_\_\_\_\_\_
% mdparagraph:P717:end

% output: NEW-U16-R05-O003
% mdparagraph:P718:begin
(3) Him: \_\_\_\_\_\_\_\_\_\_\_\_\_\_\_\_\_\_\_\_
% mdparagraph:P718:end

% output: NEW-U16-R05-O004
% mdparagraph:P719:begin
(4) Se cyning: \_\_\_\_\_\_\_\_\_\_\_\_\_\_\_\_\_\_\_\_
% mdparagraph:P719:end

% output: NEW-U16-R05-O005
% mdparagraph:P720:begin
(5) Ðē and ðū: \_\_\_\_\_\_\_\_\_\_\_\_\_\_\_\_\_\_\_\_
% mdparagraph:P720:end

% output: NEW-U16-R05-O006
% mdparagraph:P721:begin
(6) Construction: \_\_\_\_\_\_\_\_\_\_\_\_\_\_\_\_\_\_\_\_
% mdparagraph:P721:end

\end{minipage}\par\addvspace{4pt}

\par\begin{minipage}{\linewidth}\raggedright
% mdparagraph:P722:begin
\textbf{R06}
% mdparagraph:P722:end

% item: NEW-U16-R06
% mdparagraph:P723:begin
Consult printed page 169, original notes 6 and 10 (not text lines 6 and 10), together with the modern help explaining their Latin parallels in English. For each location, give a defensible translation of the actual Old English and explain the limit the note places on a confident interpretation. No knowledge or translation of Latin is required; do not replace the Old English with the Latin account.
% mdparagraph:P723:end

% mdparagraph:P724:begin
Selected passages: 169.5–7; 169.11–13.
% mdparagraph:P724:end

% mdparagraph:P725:begin
Source wording for the cited locations:
% mdparagraph:P725:end

% mdparagraph:P726:begin
169.05: nǣfre feallan ne lēt. Se cyning ðā oncnēow ðæs geongan
% mdparagraph:P726:end

% mdparagraph:P727:begin
169.06: snelnesse, ðæt hē wiste ðæt hē næfde his gelīcan on
% mdparagraph:P727:end

% mdparagraph:P728:begin
169.07: ðǣm plegan. Ðā cwæð hē tō his gefērum: “Gāð ēow
% mdparagraph:P728:end

% mdparagraph:P729:begin
169.11: mid gelǣredre handa hē swang ðone top mid swā
% mdparagraph:P729:end

% mdparagraph:P730:begin
169.12: micelre swiftnesse ðæt ðǣm cyninge wæs geðūht swilce
% mdparagraph:P730:end

% mdparagraph:P731:begin
169.13: hē of ielde tō geoguðe gewęnd wǣre. And, æfter ðǣm,
% mdparagraph:P731:end

% output: NEW-U16-R06-O001
% mdparagraph:P732:begin
(1) 169.5–7 translation: \_\_\_\_\_\_\_\_\_\_\_\_\_\_\_\_\_\_\_\_
% mdparagraph:P732:end

% output: NEW-U16-R06-O002
% mdparagraph:P733:begin
(2) 169.5–7 limit: \_\_\_\_\_\_\_\_\_\_\_\_\_\_\_\_\_\_\_\_
% mdparagraph:P733:end

% output: NEW-U16-R06-O003
% mdparagraph:P734:begin
(3) 169.11–13 translation: \_\_\_\_\_\_\_\_\_\_\_\_\_\_\_\_\_\_\_\_
% mdparagraph:P734:end

% output: NEW-U16-R06-O004
% mdparagraph:P735:begin
(4) 169.11–13 limit: \_\_\_\_\_\_\_\_\_\_\_\_\_\_\_\_\_\_\_\_
% mdparagraph:P735:end

\end{minipage}\par\addvspace{4pt}

\reserve{8\baselineskip}
% mdparagraph:P736:begin
\section*{Sources and editorial notes}
% mdparagraph:P736:end

\begingroup\fontsize{10}{12.8}\selectfont
\setlength{\parskip}{1.5pt}
% mdparagraph:P737:begin
\textsuperscript{1}: Cook, \emph{A First Book in Old English}, §§43–53, with the noun entries in \emph{Exercises in Old English}, pp. 18–22. For nętt, Wright, \emph{Old English Grammar}, §§355–356 explicitly assigns nett to the cyn(n) type. The limited singular meanings are teaching limits, not claims of historical defectiveness.
% mdparagraph:P737:end

% mdparagraph:P738:begin
\textsuperscript{2}: Cook §§57–62 and 79; Wright §§435–436 for the u/\allowbreak{}o/\allowbreak{}w alternation of nearu. The tables retain the selected complete uncontracted systems where alternatives are possible.
% mdparagraph:P738:end

% mdparagraph:P739:begin
\textsuperscript{3}: Cook §§95–116 and the cited vocabulary entries; Wright pp. 262–263 for contracted slēan/\allowbreak{}ðwēan preterites and pp. 270–274 for the gierwan/\allowbreak{}smierwan pattern. The selected spellings are teaching forms; attested alternatives are not thereby declared ungrammatical.
% mdparagraph:P739:end

% mdparagraph:P740:begin
\textsuperscript{4}: Cook §95, note: the plural command has the shortened form before a following subject pronoun. Here the additional productive position is specifically dēme gē.
% mdparagraph:P740:end

% mdparagraph:P741:begin
\textsuperscript{5}: Bosworth–Toller, s.v. cunnan, distinguishes acquaintance with a person from the infinitival ability construction; personal-object witan is also attested and is expressly requested by Cook 35.7.
% mdparagraph:P741:end

% mdparagraph:P742:begin
\textsuperscript{6}: Cook, \emph{Exercises in Old English} (1895), joint Group XII, XIII, pp. 18–22, Exercises 32–36. All 59 numbered prompts and their controls are retained.
% mdparagraph:P742:end

% mdparagraph:P743:begin
\textsuperscript{7}: The printed hint in 34.9 confuses weak āgyltan with strong āgyldan “repay; reward”. Cook’s strong-verb treatment, Bosworth–Toller s.v. ā-gildan, and the West Saxon Matthew 16.27 support the correction and contracted āgylt. The English prompt itself is unchanged.
% mdparagraph:P743:end

% mdparagraph:P744:begin
\textsuperscript{8}: The West Saxon Gospels, Matthew 14.8 and 14.10, printed p. 114, have genitive Iohannes and accusative Iohannem. Initial I → J is the only declared normalization in the supplied name forms.
% mdparagraph:P744:end

% mdparagraph:P745:begin
\textsuperscript{9}: Cook, \emph{A First Book in Old English}, second edition (1895), pp. 168–170, “The Incidents in the Gymnasium”; one-based PDF pp. 184–186. All 54 original notes and 59 anchors are retained. Modern help interprets the exact reading without replacing it by its Latin parallels. For ceroma, the Latin lexicon meaning is wrestlers’ ointment.
% mdparagraph:P745:end

\endgroup
\end{document}
